%%%%%%%%%%%%%%%%%%%%%%%%%%%%%%%%%%%%%%%%%%%%%%%%%%%%%%%%%%%%%%%%%%%%
%  Real Analysis I — Course Notes
%  Part 1: Preamble through end of Chapter 2
%%%%%%%%%%%%%%%%%%%%%%%%%%%%%%%%%%%%%%%%%%%%%%%%%%%%%%%%%%%%%%%%%%%%
\documentclass[12pt,a4paper,openany]{book}

%% ── Encoding and fonts ──────────────────────────────────────────
\usepackage[utf8]{inputenc}
\usepackage[T1]{fontenc}
\usepackage{lmodern}
\usepackage{microtype}

%% ── Mathematics ─────────────────────────────────────────────────
\usepackage{amsmath,amssymb,amsthm,mathtools}

%% ── Graphics and colour ────────────────────────────────────────
\usepackage[dvipsnames]{xcolor}
\usepackage{tikz}
\usetikzlibrary{arrows.meta,positioning,calc,decorations.pathreplacing,
                 patterns,shapes.geometric}

%% ── Boxes ───────────────────────────────────────────────────────
\usepackage[most]{tcolorbox}
\tcbuselibrary{theorems,skins,breakable}

%% ── Page layout ─────────────────────────────────────────────────
\usepackage[margin=2.5cm]{geometry}
\usepackage{fancyhdr}
\pagestyle{fancy}
\fancyhf{}
\fancyhead[LE]{\slshape\leftmark}
\fancyhead[RO]{\slshape\rightmark}
\fancyfoot[C]{\thepage}
\renewcommand{\headrulewidth}{0.4pt}

%% ── Lists ───────────────────────────────────────────────────────
\usepackage{enumitem}

%% ── Index and hyperlinks ────────────────────────────────────────
\usepackage{makeidx}
\makeindex
\usepackage[colorlinks=true,linkcolor=blue!70!black,
            urlcolor=blue!60!black,citecolor=green!50!black]{hyperref}

%% ── Captions outside floats ──────────────────────────────────────
\usepackage{caption}

%% ── Mini table of contents per chapter ──────────────────────────
\usepackage{minitoc}

%% ── Drop caps ───────────────────────────────────────────────────
\usepackage{lettrine}

%% ── Miscellaneous ───────────────────────────────────────────────
\usepackage{booktabs,array}

%% ═══════════════════════════════════════════════════════════════
%%  CUSTOM COMMANDS
%% ═══════════════════════════════════════════════════════════════
\newcommand{\N}{\mathbb{N}}
\newcommand{\Z}{\mathbb{Z}}
\newcommand{\Q}{\mathbb{Q}}
\newcommand{\R}{\mathbb{R}}
\newcommand{\C}{\mathbb{C}}
\newcommand{\eps}{\varepsilon}
\newcommand{\abs}[1]{\left\lvert #1 \right\rvert}
\newcommand{\norm}[1]{\left\lVert #1 \right\rVert}
\newcommand{\set}[1]{\left\{#1\right\}}
\newcommand{\setcond}[2]{\left\{\,#1\;\middle|\;#2\,\right\}}

%% ═══════════════════════════════════════════════════════════════
%%  THEOREM ENVIRONMENTS
%% ═══════════════════════════════════════════════════════════════

% --- Theorem environments (amsthm + tcolorbox wrapping) ---
\theoremstyle{definition}
\newtheorem{definition}{Definition}[chapter]
\newtheorem{theorem}{Theorem}[chapter]
\newtheorem{proposition}{Proposition}[chapter]
\newtheorem{lemma}{Lemma}[chapter]
\newtheorem{corollary}{Corollary}[chapter]
\theoremstyle{remark}
\newtheorem{remark}{Remark}[chapter]
\newtheorem{example}{Example}[chapter]
\newtheorem{exercise}{Exercise}[chapter]

\tcolorboxenvironment{definition}{enhanced,breakable,colback=blue!5,colframe=blue!70!black,top=4pt,bottom=4pt,before skip=10pt,after skip=10pt}
\tcolorboxenvironment{theorem}{enhanced,breakable,colback=red!5,colframe=red!70!black,top=4pt,bottom=4pt,before skip=10pt,after skip=10pt}
\tcolorboxenvironment{proposition}{enhanced,breakable,colback=orange!5,colframe=orange!70!black,top=4pt,bottom=4pt,before skip=10pt,after skip=10pt}
\tcolorboxenvironment{lemma}{enhanced,breakable,colback=yellow!5,colframe=yellow!60!black,top=4pt,bottom=4pt,before skip=10pt,after skip=10pt}
\tcolorboxenvironment{corollary}{enhanced,breakable,colback=green!5,colframe=green!60!black,top=4pt,bottom=4pt,before skip=10pt,after skip=10pt}

% --- Common errors box ---
\newtcolorbox{commonerrors}{%
  enhanced, breakable,
  colback=red!3, colframe=red!80!black,
  title={\textbf{Common Errors --- Watch Out!}},
  fonttitle=\bfseries\color{white},
  coltitle=white,
  top=4pt, bottom=4pt,
  before skip=10pt, after skip=10pt
}

% --- Proof strategy box ---
\newtcolorbox{strategy}[1][]{%
  enhanced, breakable,
  colback=violet!5, colframe=violet!70!black,
  title={\textbf{Strategy: }#1},
  fonttitle=\bfseries,
  top=4pt, bottom=4pt,
  before skip=10pt, after skip=10pt
}

% --- Chapter summary box ---
\newtcolorbox{chaptersummary}{%
  enhanced, breakable,
  colback=teal!5, colframe=teal!70!black,
  title={\textbf{Chapter Summary}},
  fonttitle=\bfseries\color{white},
  coltitle=white,
  top=4pt, bottom=4pt,
  before skip=10pt, after skip=10pt
}

% --- Difficulty stars ---
\newcommand{\easy}{\textnormal{[$\star$]}}
\newcommand{\medium}{\textnormal{[$\star\star$]}}
\newcommand{\hard}{\textnormal{[$\star\star\star$]}}

%% ═══════════════════════════════════════════════════════════════
\begin{document}
%% ═══════════════════════════════════════════════════════════════

%% ── Title page ──────────────────────────────────────────────────
\begin{titlepage}
\begin{tikzpicture}[remember picture, overlay]
  \fill[left color=blue!15!white, right color=blue!5!white]
    (current page.south west) rectangle (current page.north east);
  \fill[color=orange!70!yellow]
    ([yshift=-2cm]current page.north west) rectangle ([yshift=-2.3cm]current page.north east);
  \fill[color=blue!80!black, rounded corners=8pt]
    ([xshift=3cm, yshift=-4cm]current page.north west) rectangle
    ([xshift=-3cm, yshift=-10cm]current page.north east);
  \node[anchor=center, text=white, font=\Huge\bfseries]
    at ([yshift=-6cm]current page.north) {Real Analysis I};
  \node[anchor=center, text=orange!80!yellow, font=\Large]
    at ([yshift=-7.5cm]current page.north) {Lecture Notes};
  \node[anchor=center, text=white!80, font=\large]
    at ([yshift=-8.8cm]current page.north) {Licence L2 --- 2025--2026};
  \node[anchor=center, text=blue!80!black, font=\Large\itshape]
    at ([yshift=-12cm]current page.north) {Ya\'e Ulrich Gaba};
  \draw[orange!70!yellow, line width=2pt]
    ([xshift=5cm, yshift=-13.5cm]current page.north west) --
    ([xshift=-5cm, yshift=-13.5cm]current page.north east);
  \node[anchor=center, text width=12cm, align=center, text=blue!60!black, font=\itshape]
    at ([yshift=-16cm]current page.north) {``Rigour is to the mathematician what morality is to man.''\\[4pt]--- Henri Poincar\'e};
  \node[anchor=south, text=gray, font=\small]
    at ([yshift=1.5cm]current page.south) {\today};
\end{tikzpicture}
\end{titlepage}

%% ── Table of contents ──────────────────────────────────────────
\frontmatter
\tableofcontents
\dominitoc

%% ═══════════════════════════════════════════════════════════════
%%  PREFACE
%% ═══════════════════════════════════════════════════════════════
\chapter*{Preface}
\addcontentsline{toc}{chapter}{Preface}
\markboth{Preface}{Preface}

Welcome to Real Analysis.

If you are reading these notes, you are most likely beginning your first year
of university mathematics.  You may have enjoyed mathematics at school ---
perhaps you found satisfaction in solving equations, computing integrals, or
sketching curves.  The course you are about to undertake will be \emph{very
different} from what you have experienced before, and that is precisely what
makes it exciting.

\medskip

\textbf{Why rigour?}\quad
At school you were told that $\sqrt{2}$ is irrational, that continuous
functions satisfy the Intermediate Value Theorem, and that every
differentiable function is continuous.  You likely accepted these facts on
faith or saw informal justifications.  In this course we will \emph{prove}
every single one of them, starting from clearly stated axioms.  This is not
pedantry; it is the only way to be certain that our conclusions are correct.
History is full of examples where intuition alone led mathematicians astray
--- plausible-sounding statements that turn out to be false, and
counter-intuitive results that turn out to be true.  Rigorous proof is our
shield against error.

\medskip

\textbf{What to expect.}\quad
The first weeks may feel slow.  We will spend a significant amount of time on
logic, on the axioms of the real numbers, and on the precise definition of a
limit --- before we even begin to compute anything.  This is intentional.
Think of it as learning the grammar of a language before writing essays: the
investment pays enormous dividends later.  Once the foundations are solid, the
major theorems of the course will unfold naturally and almost inevitably.

\medskip

\textbf{The journey.}\quad
Here is a brief overview of the path ahead.

\begin{itemize}[leftmargin=2em]
\item \textbf{Chapter~1} sets up the language of mathematics: logic,
  quantifiers, proof techniques, sets, and functions.  If you have never
  written a formal proof before, this chapter is your essential training
  ground.
\item \textbf{Chapter~2} introduces the real number system and, most
  importantly, the \emph{completeness axiom} --- the single property that
  distinguishes $\R$ from $\Q$.  This axiom is the engine that powers nearly
  every deep result in the course.
\item Later chapters (not included in this part) will develop the theory of
  sequences, series, continuous functions, and differentiation.
\end{itemize}

\textbf{How to study.}\quad
Mathematics is not a spectator sport.  Reading these notes is necessary but
far from sufficient.  You must \emph{write proofs yourself}, struggle with
exercises, make mistakes, and correct them.  When you read a proof, close
the notes and try to reproduce the argument from memory.  When you get stuck
on an exercise, do not look at the solution immediately --- give yourself at
least thirty minutes of honest effort.  The understanding that comes from
wrestling with a problem is far deeper than the understanding that comes
from reading someone else's solution.

\medskip

\textbf{A word of encouragement.}\quad
If you find the first proofs difficult, that is completely normal.  Every
mathematician alive today once struggled with their first $\eps$-$\delta$
argument.  Persist, ask questions, and trust the process.  By the end of the
semester, you will look back at these first chapters and wonder why they ever
seemed hard.

\medskip

Good luck, and enjoy the journey.

%% ═══════════════════════════════════════════════════════════════
%%  NOTATION
%% ═══════════════════════════════════════════════════════════════
\chapter*{Notation}
\addcontentsline{toc}{chapter}{Notation}
\markboth{Notation}{Notation}

The following symbols and conventions are used throughout these notes.

\bigskip

\begin{tabular}{@{}l l@{}}
\toprule
\textbf{Symbol} & \textbf{Meaning} \\
\midrule
$\N$ & The set of natural numbers $\{0,1,2,3,\dots\}$ \\
$\N^*$ & The set of positive natural numbers $\{1,2,3,\dots\}$ \\
$\Z$ & The set of integers \\
$\Q$ & The set of rational numbers \\
$\R$ & The set of real numbers \\
$\R^*$ & $\R \setminus \{0\}$ \\
$\R_+$ & $\{x \in \R : x \geq 0\}$ \\
$\C$ & The set of complex numbers \\
$\emptyset$ & The empty set \\
$\eps$ & A strictly positive real number (usually ``small'') \\
$\abs{x}$ & Absolute value of $x$ \\
$\norm{x}$ & Norm of $x$ \\
$\forall$ & ``For all'' (universal quantifier) \\
$\exists$ & ``There exists'' (existential quantifier) \\
$\exists!$ & ``There exists a unique'' \\
$\implies$ & ``Implies'' (logical implication) \\
$\iff$ & ``If and only if'' (logical equivalence) \\
$A \subset B$ & $A$ is a subset of $B$ (possibly equal) \\
$A \subsetneq B$ & $A$ is a strict subset of $B$ \\
$A \cup B$ & Union of $A$ and $B$ \\
$A \cap B$ & Intersection of $A$ and $B$ \\
$A \setminus B$ & Set difference \\
$A^c$ & Complement of $A$ \\
$A \times B$ & Cartesian product of $A$ and $B$ \\
$f : A \to B$ & Function from $A$ to $B$ \\
$f \circ g$ & Composition of $f$ and $g$ \\
$\sup A$ & Supremum (least upper bound) of $A$ \\
$\inf A$ & Infimum (greatest lower bound) of $A$ \\
$\max A$, $\min A$ & Maximum, minimum of $A$ (when they exist) \\
$\lfloor x \rfloor$ & Floor (integer part) of $x$ \\
\bottomrule
\end{tabular}

\bigskip

\noindent\textbf{Convention.}\quad Unless otherwise stated, the letters
$m,n,k,p$ denote natural numbers, the letters $x,y,z,a,b,c$ denote real
numbers, and $\eps,\delta,\eta$ denote \emph{strictly positive} real numbers.

%% ═══════════════════════════════════════════════════════════════
%%  MAIN MATTER
%% ═══════════════════════════════════════════════════════════════
\mainmatter

%%%%%%%%%%%%%%%%%%%%%%%%%%%%%%%%%%%%%%%%%%%%%%%%%%%%%%%%%%%%%%%%%%%%
%%%%%%%%%%%%%%%%%%%%%%%%%%%%%%%%%%%%%%%%%%%%%%%%%%%%%%%%%%%%%%%%%%%%
%%
%%  CHAPTER 1 : Logic, Sets, and Mathematical Reasoning
%%
%%%%%%%%%%%%%%%%%%%%%%%%%%%%%%%%%%%%%%%%%%%%%%%%%%%%%%%%%%%%%%%%%%%%
%%%%%%%%%%%%%%%%%%%%%%%%%%%%%%%%%%%%%%%%%%%%%%%%%%%%%%%%%%%%%%%%%%%%
\chapter{Logic, Sets, and Mathematical Reasoning}
\label{ch:logic}
\index{logic}

Before we can prove anything about real numbers, sequences, or functions,
we need a precise language in which to express our statements and a toolkit
of proof techniques to establish their truth.  This chapter provides both.
It is the single most important chapter in the course: every subsequent
proof relies on the material developed here.

%% ──────────────────────────────────────────────────────────────
\section{Propositions and Logical Connectives}
\label{sec:propositions}
\index{proposition}

\begin{definition}[Proposition]\label{def:proposition}
A \textbf{proposition}\index{proposition} (or \textbf{statement}) is a
sentence that is either \emph{true}~(T) or \emph{false}~(F), but not both.
\end{definition}

\begin{example}\label{ex:prop-examples}
The following are propositions:
\begin{enumerate}[label=(\alph*)]
\item ``$2+3=5$'' \quad (true).
\item ``Every even number greater than $2$ is the sum of two primes''
  \quad (this is Goldbach's conjecture --- it is either true or false, even
  though we do not currently know which).
\item ``$\sqrt{2}$ is rational'' \quad (false, as we shall prove in
  Section~\ref{sec:contradiction}).
\end{enumerate}
The sentence ``Is $7$ a prime number?''\ is \emph{not} a proposition
because it is a question, not a declarative statement.  Similarly, ``$x>3$''
is not a proposition by itself because its truth value depends on the
unspecified variable~$x$; it is a \textbf{predicate}\index{predicate}.
\end{example}

\begin{definition}[Logical connectives]\label{def:connectives}
\index{connective!logical}
Let $P$ and $Q$ be propositions.  We define the following connectives:
\begin{enumerate}[label=(\roman*)]
  \item \textbf{Negation}\index{negation} $\neg P$ (``not $P$''): true when
    $P$ is false, and false when $P$ is true.
  \item \textbf{Conjunction}\index{conjunction} $P \land Q$ (``$P$ and
    $Q$''): true when both $P$ and $Q$ are true, false otherwise.
  \item \textbf{Disjunction}\index{disjunction} $P \lor Q$ (``$P$ or
    $Q$''): true when at least one of $P,Q$ is true, false only when both
    are false.  (This is the \emph{inclusive} ``or''.)
  \item \textbf{Implication}\index{implication} $P \implies Q$ (``if $P$
    then $Q$'', or ``$P$ implies $Q$''): false \emph{only} when $P$ is true
    and $Q$ is false; true in all other cases.
  \item \textbf{Equivalence}\index{equivalence} $P \iff Q$ (``$P$ if and
    only if $Q$''): true when $P$ and $Q$ have the same truth value, false
    otherwise.
\end{enumerate}
\end{definition}

\begin{remark}\label{rem:implication-vacuous}
The definition of implication often surprises students.  In particular, the
statement $P \implies Q$ is \emph{true} whenever $P$ is false, regardless
of~$Q$.  This is called a \textbf{vacuously true}\index{vacuously true}
implication.  For instance, ``If $1=0$, then the moon is made of cheese'' is
a \emph{true} implication.  This convention is not arbitrary; it is the only
definition that makes the logical calculus consistent and useful.
\end{remark}

\paragraph{Truth tables.}\index{truth table}
The following table summarises all the connectives.

\medskip
\begin{center}
\begin{tabular}{cc|ccccc}
\toprule
$P$ & $Q$ & $\neg P$ & $P\land Q$ & $P\lor Q$ & $P\implies Q$ & $P\iff Q$\\
\midrule
T & T & F & T & T & T & T \\
T & F & F & F & T & F & F \\
F & T & T & F & T & T & F \\
F & F & T & F & F & T & T \\
\bottomrule
\end{tabular}
\end{center}

\begin{remark}\label{rem:converse-contrapositive}
Given the implication $P \implies Q$:
\begin{itemize}
  \item The \textbf{converse}\index{converse} is $Q \implies P$.  It is
    \emph{not} logically equivalent to the original implication.
  \item The \textbf{contrapositive}\index{contrapositive} is $\neg Q
    \implies \neg P$.  It \emph{is} logically equivalent to $P \implies Q$
    (check the truth table!).  This fact is the basis for proof by
    contrapositive.
\end{itemize}
\end{remark}

%% ──────────────────────────────────────────────────────────────
\section{Quantifiers}
\label{sec:quantifiers}
\index{quantifier}

In mathematics we constantly make statements about all elements of a set or
about the existence of elements with a particular property.  The symbols
$\forall$ and $\exists$ make such statements precise.

\begin{definition}[Universal and existential quantifiers]\label{def:quantifiers}
Let $P(x)$ be a predicate depending on a variable $x$ ranging over a set~$E$.
\begin{enumerate}[label=(\roman*)]
  \item \textbf{Universal quantifier}\index{quantifier!universal}:
    $\forall x \in E,\; P(x)$ means ``for every element $x$ of $E$,
    the property $P(x)$ holds.''
  \item \textbf{Existential quantifier}\index{quantifier!existential}:
    $\exists x \in E,\; P(x)$ means ``there exists at least one element
    $x$ of $E$ such that $P(x)$ holds.''
\end{enumerate}
\end{definition}

\subsection{Negation of Quantified Statements}
\label{subsec:negation-quantifiers}
\index{negation!of quantified statements}

This is one of the most important skills you will develop in this course.
The rules are:
\begin{align}
  \neg\bigl(\forall x \in E,\; P(x)\bigr)
    &\;\iff\; \exists x \in E,\; \neg P(x),
    \label{eq:neg-forall} \\
  \neg\bigl(\exists x \in E,\; P(x)\bigr)
    &\;\iff\; \forall x \in E,\; \neg P(x).
    \label{eq:neg-exists}
\end{align}
In words: to negate a $\forall$, change it to $\exists$ and negate the
predicate.  To negate an $\exists$, change it to $\forall$ and negate the
predicate.  When there are \emph{multiple} quantifiers, you proceed from
left to right, flipping each one.

\begin{strategy}[Negating a multiply-quantified statement]
  Given a statement of the form
  $\forall x,\; \exists y,\; \forall z,\; P(x,y,z)$,
  its negation is obtained by flipping every quantifier and negating the
  innermost predicate:
  $\exists x,\; \forall y,\; \exists z,\; \neg P(x,y,z)$.
\end{strategy}

\begin{example}[Negating the definition of convergence]
\label{ex:negate-convergence}
\index{convergence!negation of definition}
Consider the statement: ``The sequence $(a_n)$ converges to $\ell$'', which
is formally written as
\[
  \forall \eps > 0,\; \exists N \in \N,\; \forall n \geq N,\;
  \abs{a_n - \ell} < \eps.
  \tag{$\mathcal{S}$}
\]
We negate this step by step.

\textbf{Step~1.}\quad The outermost quantifier is ``$\forall \eps>0$''.
Flip it to ``$\exists \eps>0$'' and negate everything that follows:
\[
  \neg(\mathcal{S}) \;\iff\;
  \exists \eps > 0,\; \neg\Bigl(
    \exists N \in \N,\; \forall n \geq N,\;
    \abs{a_n - \ell} < \eps
  \Bigr).
\]

\textbf{Step~2.}\quad The next quantifier is ``$\exists N \in \N$''.
Flip it to ``$\forall N \in \N$'' and negate the rest:
\[
  \neg(\mathcal{S}) \;\iff\;
  \exists \eps > 0,\; \forall N \in \N,\; \neg\Bigl(
    \forall n \geq N,\;
    \abs{a_n - \ell} < \eps
  \Bigr).
\]

\textbf{Step~3.}\quad The next quantifier is ``$\forall n \geq N$''.
Flip it to ``$\exists n \geq N$'' and negate the predicate:
\[
  \neg(\mathcal{S}) \;\iff\;
  \exists \eps > 0,\; \forall N \in \N,\; \exists n \geq N,\;
  \abs{a_n - \ell} \geq \eps.
\]

\textbf{Interpretation.}\quad The sequence $(a_n)$ does \emph{not} converge
to $\ell$ if and only if: \emph{there exists} some $\eps > 0$ such that,
\emph{no matter how large} we choose $N$, we can always \emph{find} some
$n \geq N$ for which $\abs{a_n-\ell} \geq \eps$.  In other words, there
are terms of the sequence that remain ``far'' from $\ell$ indefinitely.
\end{example}

\begin{example}\label{ex:negate-bounded}
Consider the statement ``The function $f$ is bounded above on $[0,1]$'':
\[
  \exists M \in \R,\; \forall x \in [0,1],\; f(x) \leq M.
\]
Its negation is:
\[
  \forall M \in \R,\; \exists x \in [0,1],\; f(x) > M.
\]
This says: no matter how large a bound $M$ you propose, I can find a point
$x$ in $[0,1]$ where $f(x)$ exceeds~$M$.
\end{example}

%% ──────────────────────────────────────────────────────────────
\section{Proof Techniques}
\label{sec:proof-techniques}
\index{proof techniques}

We now present the main methods of proof.  For each method, we first explain
the logical principle, then give worked examples.

%% ── Direct proof ────────────────────────────────────────────────
\subsection{Direct Proof}
\label{subsec:direct-proof}
\index{proof techniques!direct proof}

\begin{strategy}[Direct proof]
To prove ``$P \implies Q$'', assume $P$ is true and deduce, through a chain
of logical steps, that $Q$ is true.
\end{strategy}

\begin{proposition}[Sum of two even integers]\label{prop:sum-even}
The sum of two even integers is even.
\end{proposition}

\begin{proof}
Let $a$ and $b$ be even integers.  By definition of ``even'',\index{even
integer} there exist integers $k$ and $m$ such that $a = 2k$ and $b = 2m$.
Then
\[
  a + b = 2k + 2m = 2(k + m).
\]
Since $k + m$ is an integer, $a+b$ is of the form $2 \cdot (\text{integer})$,
hence $a+b$ is even.
\end{proof}

\begin{proposition}[Product of rational and irrational]\label{prop:rat-times-irr}
Let $q \in \Q$ with $q \neq 0$, and let $\alpha \notin \Q$.  Then
$q\alpha \notin \Q$.
\end{proposition}

\begin{proof}
We prove the contrapositive in the next subsection; here we give a direct
argument.  Suppose for the sake of deriving the conclusion directly.  We have
$q = \frac{a}{b}$ with $a,b \in \Z$, $b \neq 0$, $a \neq 0$.  Suppose, aiming
for a contradiction with the hypothesis that $\alpha$ is irrational, that
$q\alpha$ were rational.  Then $q\alpha = \frac{p}{r}$ for some $p,r\in\Z$,
$r\neq0$.  But then $\alpha = \frac{q\alpha}{q} = \frac{p}{r}\cdot\frac{b}{a}
= \frac{pb}{ra} \in \Q$, contradicting the assumption that $\alpha$ is
irrational.  Therefore $q\alpha \notin \Q$.
\end{proof}

%% ── Proof by contrapositive ─────────────────────────────────────
\subsection{Proof by Contrapositive}
\label{subsec:contrapositive}
\index{proof techniques!contrapositive}

Recall that $P \implies Q$ is logically equivalent to $\neg Q \implies \neg P$.
Sometimes it is easier to prove the contrapositive.

\begin{strategy}[Proof by contrapositive]
To prove ``$P \implies Q$'', assume $\neg Q$ is true and deduce $\neg P$.
\end{strategy}

\begin{proposition}\label{prop:n-squared-even}
Let $n$ be an integer.  If $n^2$ is even, then $n$ is even.
\end{proposition}

\begin{proof}
We prove the contrapositive: ``If $n$ is odd, then $n^2$ is odd.''

Assume $n$ is odd.  Then there exists an integer $k$ such that $n = 2k+1$.
Therefore
\[
  n^2 = (2k+1)^2 = 4k^2 + 4k + 1 = 2(2k^2 + 2k) + 1.
\]
Since $2k^2+2k$ is an integer, $n^2$ is of the form $2\cdot(\text{integer})+1$,
so $n^2$ is odd.  This completes the proof of the contrapositive, and hence
the original statement is proved.
\end{proof}

\begin{proposition}\label{prop:contrapositive-div}
Let $n \in \Z$.  If $n^2$ is not divisible by~$3$, then $n$ is not divisible
by~$3$.
\end{proposition}

\begin{proof}
We prove the contrapositive: ``If $n$ is divisible by~$3$, then $n^2$ is
divisible by~$3$.''

Assume $3 \mid n$.  Then $n = 3k$ for some $k \in \Z$.  Hence
$n^2 = 9k^2 = 3(3k^2)$, so $3 \mid n^2$.
\end{proof}

%% ── Proof by contradiction ──────────────────────────────────────
\subsection{Proof by Contradiction}
\label{sec:contradiction}
\index{proof techniques!contradiction}

\begin{strategy}[Proof by contradiction]
To prove a statement~$P$, assume $\neg P$ and derive a logical contradiction
(i.e., show that $\neg P$ implies something that is both true and false).
Since mathematics cannot contain contradictions, our assumption $\neg P$ must
have been wrong, so $P$ is true.
\end{strategy}

\begin{theorem}[$\sqrt{2}$ is irrational]\label{thm:sqrt2}
\index{$\sqrt{2}$!irrationality}
There is no rational number whose square equals~$2$.  In other words,
$\sqrt{2} \notin \Q$.
\end{theorem}

\begin{proof}
Suppose, for the sake of contradiction, that $\sqrt{2} \in \Q$.  Then we can
write
\[
  \sqrt{2} = \frac{p}{q},
\]
where $p,q \in \N^*$ and the fraction $\frac{p}{q}$ is in lowest terms
(i.e., $\gcd(p,q)=1$).

Squaring both sides gives $2 = \frac{p^2}{q^2}$, hence
\begin{equation}\label{eq:sqrt2-step}
  p^2 = 2q^2.
\end{equation}
This tells us that $p^2$ is even.  By Proposition~\ref{prop:n-squared-even}
(or its contrapositive, which we proved above), $p$ itself must be even.
So we can write $p = 2k$ for some $k \in \N^*$.

Substituting into \eqref{eq:sqrt2-step}:
\[
  (2k)^2 = 2q^2 \implies 4k^2 = 2q^2 \implies q^2 = 2k^2.
\]
By the same reasoning, $q^2$ is even, so $q$ is even.

But now both $p$ and $q$ are even, which contradicts our assumption that
$\gcd(p,q)=1$ (since they share the common factor~$2$).

This contradiction shows that the initial assumption ``$\sqrt{2}\in\Q$''
is false.  Therefore $\sqrt{2} \notin \Q$.
\end{proof}

\begin{proposition}[Infinitude of primes]\label{prop:inf-primes}
\index{prime numbers!infinitude}
There are infinitely many prime numbers.
\end{proposition}

\begin{proof}
Suppose, for the sake of contradiction, that there are only finitely many
primes, say $p_1, p_2, \dots, p_k$.  Consider the number
\[
  N = p_1 \cdot p_2 \cdots p_k + 1.
\]
Since $N > 1$, either $N$ is itself prime or $N$ has a prime factor.  In
either case, there exists a prime $p$ that divides~$N$.  Since $p_1, \dots,
p_k$ are all the primes, we must have $p = p_i$ for some~$i$.

But $p_i$ divides the product $p_1 p_2 \cdots p_k$, and $p_i$ divides $N$,
so $p_i$ divides
\[
  N - p_1 p_2 \cdots p_k = 1.
\]
No prime divides~$1$, so we have a contradiction.  Therefore the set of
primes is infinite.
\end{proof}

%% ── Mathematical induction ──────────────────────────────────────
\subsection{Mathematical Induction}
\label{subsec:induction}
\index{induction}

\begin{theorem}[Principle of Mathematical Induction]\label{thm:induction}
Let $P(n)$ be a statement depending on a natural number $n \geq n_0$.
Suppose:
\begin{enumerate}[label=(\roman*)]
  \item \textbf{Base case:}\index{induction!base case} $P(n_0)$ is true.
  \item \textbf{Inductive step:}\index{induction!inductive step} For every
    $n \geq n_0$, $P(n) \implies P(n+1)$.
\end{enumerate}
Then $P(n)$ is true for every $n \geq n_0$.
\end{theorem}

\begin{remark}
Think of induction like an infinite row of dominoes.  The base case knocks
over the first domino.  The inductive step guarantees that each falling
domino knocks over the next.  Together, they ensure that \emph{every} domino
falls.
\end{remark}

\begin{example}[Sum of the first $n$ natural numbers]\label{ex:induction-sum}
\index{induction!example}
We prove that for every $n \geq 1$,
\begin{equation}\label{eq:sum-n}
  1 + 2 + \cdots + n = \frac{n(n+1)}{2}.
\end{equation}

\textbf{Base case} ($n=1$).\quad The left-hand side is~$1$.  The right-hand
side is $\frac{1 \cdot 2}{2} = 1$.  So $P(1)$ holds.

\textbf{Inductive step.}\quad Let $n \geq 1$ and assume $P(n)$ is true,
i.e., $1 + 2 + \cdots + n = \frac{n(n+1)}{2}$.  We must show $P(n+1)$:
\begin{align*}
  1 + 2 + \cdots + n + (n+1)
    &= \frac{n(n+1)}{2} + (n+1)
       \quad\text{(by the inductive hypothesis)} \\
    &= \frac{n(n+1) + 2(n+1)}{2} \\
    &= \frac{(n+1)(n+2)}{2},
\end{align*}
which is exactly the formula with $n$ replaced by $n+1$.  So $P(n+1)$ holds.

By the Principle of Mathematical Induction, \eqref{eq:sum-n} holds for all
$n \geq 1$.
\end{example}

\begin{example}[Bernoulli's inequality]\label{ex:bernoulli}
\index{Bernoulli's inequality}
Let $x \geq -1$.  We prove by induction that for every $n \geq 1$,
\begin{equation}\label{eq:bernoulli}
  (1+x)^n \geq 1 + nx.
\end{equation}

\textbf{Base case} ($n=1$).\quad $(1+x)^1 = 1+x = 1+1\cdot x$.  The
inequality holds with equality.

\textbf{Inductive step.}\quad Assume $(1+x)^n \geq 1+nx$ for some $n\geq1$.
Then
\begin{align*}
  (1+x)^{n+1}
    &= (1+x)^n \cdot (1+x) \\
    &\geq (1+nx)(1+x) \quad\text{(inductive hypothesis; note $1+x \geq 0$)} \\
    &= 1 + nx + x + nx^2 \\
    &= 1 + (n+1)x + nx^2 \\
    &\geq 1 + (n+1)x \quad\text{(since $nx^2 \geq 0$)}.
\end{align*}
This is exactly the desired inequality with $n$ replaced by $n+1$.  By
induction, \eqref{eq:bernoulli} holds for all $n \geq 1$.
\end{example}

\begin{theorem}[Principle of Strong Induction]\label{thm:strong-induction}
\index{induction!strong}
Let $P(n)$ be a statement depending on a natural number $n \geq n_0$.
Suppose:
\begin{enumerate}[label=(\roman*)]
  \item $P(n_0)$ is true.
  \item For every $n \geq n_0$: if $P(k)$ is true for \emph{all}
    $n_0 \leq k \leq n$, then $P(n+1)$ is true.
\end{enumerate}
Then $P(n)$ is true for every $n \geq n_0$.
\end{theorem}

\begin{example}[Every integer $\geq 2$ has a prime factor]
\label{ex:strong-ind-prime}
We prove by strong induction that every integer $n \geq 2$ has at least one
prime factor.

\textbf{Base case} ($n=2$).\quad $2$ is itself prime, so it is its own prime
factor.

\textbf{Inductive step.}\quad Let $n \geq 2$ and assume that every integer
$k$ with $2 \leq k \leq n$ has a prime factor.  Consider $n+1$.  There are
two cases:
\begin{itemize}
  \item If $n+1$ is prime, then $n+1$ is its own prime factor.
  \item If $n+1$ is not prime, then $n+1 = ab$ where $2 \leq a,b \leq n$.
    By the strong inductive hypothesis, $a$ has a prime factor~$p$.  Since
    $p \mid a$ and $a \mid (n+1)$, we have $p \mid (n+1)$.
\end{itemize}
In both cases, $n+1$ has a prime factor.  By strong induction, the result
holds for all $n \geq 2$.
\end{example}

%% ──────────────────────────────────────────────────────────────
\section{Sets and Set Operations}
\label{sec:sets}
\index{set}

\begin{definition}[Set, element, subset]\label{def:set}
A \textbf{set}\index{set} is a collection of distinct objects, called its
\textbf{elements}\index{element}.  We write $x \in A$ to mean ``$x$ is an
element of~$A$'', and $x \notin A$ for its negation.

A set $A$ is a \textbf{subset}\index{subset} of $B$, written $A \subset B$,
if every element of $A$ is also an element of~$B$:
\[
  A \subset B \;\iff\; \forall x,\; (x \in A \implies x \in B).
\]
Two sets are \textbf{equal} if and only if $A \subset B$ and $B \subset A$.
\end{definition}

\begin{definition}[Set operations]\label{def:set-operations}
Let $A$ and $B$ be subsets of a universal set~$\Omega$.\index{set!operations}
\begin{enumerate}[label=(\roman*)]
  \item \textbf{Union}\index{union}: $A \cup B = \setcond{x \in \Omega}{x
    \in A \text{ or } x \in B}$.
  \item \textbf{Intersection}\index{intersection}: $A \cap B =
    \setcond{x \in \Omega}{x \in A \text{ and } x \in B}$.
  \item \textbf{Set difference}\index{set difference}: $A \setminus B =
    \setcond{x \in \Omega}{x \in A \text{ and } x \notin B}$.
  \item \textbf{Complement}\index{complement}: $A^c = \Omega \setminus A =
    \setcond{x \in \Omega}{x \notin A}$.
  \item \textbf{Cartesian product}\index{Cartesian product}: $A \times B =
    \setcond{(a,b)}{a \in A \text{ and } b \in B}$.
\end{enumerate}
\end{definition}

\medskip

\begin{center}
\begin{tikzpicture}[scale=0.8]
  % Union
  \begin{scope}[shift={(-5,0)}]
    \draw[thick, fill=blue!15] (-0.8,0) circle (1.2);
    \draw[thick, fill=blue!15] (0.8,0) circle (1.2);
    \node at (-1.3,0) {$A$};
    \node at (1.3,0) {$B$};
    \node at (0,-2) {$A \cup B$};
    \draw[thick] (-2.5,-1.8) rectangle (2.5,1.8);
    \node at (2.1,1.5) {$\Omega$};
  \end{scope}
  % Intersection
  \begin{scope}[shift={(0,0)}]
    \draw[thick] (-0.8,0) circle (1.2);
    \draw[thick] (0.8,0) circle (1.2);
    \begin{scope}
      \clip (-0.8,0) circle (1.2);
      \fill[blue!25] (0.8,0) circle (1.2);
    \end{scope}
    \draw[thick] (-0.8,0) circle (1.2);
    \draw[thick] (0.8,0) circle (1.2);
    \node at (-1.3,0) {$A$};
    \node at (1.3,0) {$B$};
    \node at (0,-2) {$A \cap B$};
    \draw[thick] (-2.5,-1.8) rectangle (2.5,1.8);
    \node at (2.1,1.5) {$\Omega$};
  \end{scope}
  % Complement
  \begin{scope}[shift={(5,0)}]
    \fill[blue!15] (-2.5,-1.8) rectangle (2.5,1.8);
    \fill[white] (0,0) circle (1.2);
    \draw[thick] (0,0) circle (1.2);
    \node at (0,0) {$A$};
    \node at (0,-2) {$A^c$};
    \draw[thick] (-2.5,-1.8) rectangle (2.5,1.8);
    \node at (2.1,1.5) {$\Omega$};
  \end{scope}
\end{tikzpicture}
\end{center}

\begin{theorem}[De Morgan's laws]\label{thm:demorgan}
\index{De Morgan's laws}
Let $A$ and $B$ be subsets of a universal set~$\Omega$.  Then
\begin{enumerate}[label=(\roman*)]
  \item $(A \cup B)^c = A^c \cap B^c$,\label{item:dm1}
  \item $(A \cap B)^c = A^c \cup B^c$.\label{item:dm2}
\end{enumerate}
\end{theorem}

\begin{proof}
We prove~\ref{item:dm1}; the proof of~\ref{item:dm2} is similar and left as
an exercise.

We show the two sets are equal by showing each is a subset of the other.

\textbf{($\subset$)}\quad Let $x \in (A \cup B)^c$.  Then $x \notin A \cup B$.
This means $x \notin A$ \emph{and} $x \notin B$ (because if $x$ were in
either $A$ or $B$, then $x$ would be in $A\cup B$).  So $x \in A^c$ and
$x \in B^c$, hence $x \in A^c \cap B^c$.

\textbf{($\supset$)}\quad Let $x \in A^c \cap B^c$.  Then $x \in A^c$ and
$x \in B^c$, i.e., $x \notin A$ and $x \notin B$.  Therefore $x \notin
A \cup B$, so $x \in (A\cup B)^c$.

Since each set is a subset of the other, they are equal.
\end{proof}

%% ──────────────────────────────────────────────────────────────
\section{Functions}
\label{sec:functions}
\index{function}

\begin{definition}[Function]\label{def:function}
A \textbf{function}\index{function} $f : A \to B$ is a rule that assigns to
each element $a \in A$ a \emph{unique} element $f(a) \in B$.  The set $A$
is called the \textbf{domain}\index{domain} and $B$ the
\textbf{codomain}\index{codomain}.
\end{definition}

\begin{definition}[Injection, surjection, bijection]\label{def:inj-surj-bij}
Let $f : A \to B$ be a function.
\begin{enumerate}[label=(\roman*)]
  \item $f$ is \textbf{injective}\index{injective} (one-to-one) if
    \[
      \forall a_1, a_2 \in A,\quad f(a_1)=f(a_2) \implies a_1 = a_2.
    \]
  \item $f$ is \textbf{surjective}\index{surjective} (onto) if
    \[
      \forall b \in B,\; \exists a \in A,\quad f(a) = b.
    \]
  \item $f$ is \textbf{bijective}\index{bijective} if it is both injective
    and surjective.
\end{enumerate}
\end{definition}

\bigskip

\begin{center}
\begin{tikzpicture}[
    >=Stealth,
    set/.style={draw, ellipse, minimum width=2cm, minimum height=3cm,
                thick},
    every node/.style={font=\small}
  ]
  % Injection
  \begin{scope}[shift={(-5,0)}]
    \node[set,label=above:{$A$}] (A1) at (0,0) {};
    \node[set,label=above:{$B$}] (B1) at (2.8,0) {};
    \foreach \y/\lbl in {0.7/a, 0/b, -0.7/c} {
      \node[fill=blue!60,circle,inner sep=1.5pt,label=left:{\lbl}] (a\lbl)
        at (0,\y) {};
    }
    \foreach \y/\lbl in {0.9/1, 0.3/2, -0.3/3, -0.9/4} {
      \node[fill=red!60,circle,inner sep=1.5pt,label=right:{\lbl}] (b\lbl)
        at (2.8,\y) {};
    }
    \draw[->] (aa) -- (b1);
    \draw[->] (ab) -- (b3);
    \draw[->] (ac) -- (b4);
    \node at (1.4,-2.2) {\textbf{Injective}};
  \end{scope}
  % Surjection
  \begin{scope}[shift={(0,0)}]
    \node[set,label=above:{$A$}] (A2) at (0,0) {};
    \node[set,label=above:{$B$}] (B2) at (2.8,0) {};
    \foreach \y/\lbl in {0.9/a, 0.3/b, -0.3/c, -0.9/d} {
      \node[fill=blue!60,circle,inner sep=1.5pt,label=left:{\lbl}] (a\lbl)
        at (0,\y) {};
    }
    \foreach \y/\lbl in {0.7/1, 0/2, -0.7/3} {
      \node[fill=red!60,circle,inner sep=1.5pt,label=right:{\lbl}] (b\lbl)
        at (2.8,\y) {};
    }
    \draw[->] (aa) -- (b1);
    \draw[->] (ab) -- (b2);
    \draw[->] (ac) -- (b2);
    \draw[->] (ad) -- (b3);
    \node at (1.4,-2.2) {\textbf{Surjective}};
  \end{scope}
  % Bijection
  \begin{scope}[shift={(5,0)}]
    \node[set,label=above:{$A$}] (A3) at (0,0) {};
    \node[set,label=above:{$B$}] (B3) at (2.8,0) {};
    \foreach \y/\lbl in {0.7/a, 0/b, -0.7/c} {
      \node[fill=blue!60,circle,inner sep=1.5pt,label=left:{\lbl}] (a\lbl)
        at (0,\y) {};
    }
    \foreach \y/\lbl in {0.7/1, 0/2, -0.7/3} {
      \node[fill=red!60,circle,inner sep=1.5pt,label=right:{\lbl}] (b\lbl)
        at (2.8,\y) {};
    }
    \draw[->] (aa) -- (b2);
    \draw[->] (ab) -- (b3);
    \draw[->] (ac) -- (b1);
    \node at (1.4,-2.2) {\textbf{Bijective}};
  \end{scope}
\end{tikzpicture}
\end{center}

%% ──────────────────────────────────────────────────────────────
\section{Equivalence Relations and Order Relations}
\label{sec:relations}

\begin{definition}[Equivalence relation]\label{def:equiv-rel}
\index{equivalence relation}
A relation $\sim$ on a set $E$ is an \textbf{equivalence relation} if it is:
\begin{enumerate}[label=(\roman*)]
  \item \textbf{Reflexive}\index{reflexive}: $\forall x \in E,\; x \sim x$.
  \item \textbf{Symmetric}\index{symmetric}: $\forall x,y \in E,\;
    x \sim y \implies y \sim x$.
  \item \textbf{Transitive}\index{transitive}: $\forall x,y,z \in E,\;
    (x \sim y \text{ and } y \sim z) \implies x \sim z$.
\end{enumerate}
The \textbf{equivalence class}\index{equivalence class} of $x$ is
$[x] = \setcond{y \in E}{y \sim x}$.
\end{definition}

\begin{example}\label{ex:congruence}
Fix $n \in \N^*$.  Define $a \sim b$ on $\Z$ by $n \mid (a-b)$ (congruence
modulo~$n$\index{congruence}).  This is an equivalence relation:
\begin{itemize}
  \item Reflexive: $n \mid (a-a) = 0$.
  \item Symmetric: if $n \mid (a-b)$, then $n \mid (-(a-b)) = (b-a)$.
  \item Transitive: if $n \mid (a-b)$ and $n \mid (b-c)$, then
    $n \mid ((a-b)+(b-c)) = (a-c)$.
\end{itemize}
\end{example}

\begin{definition}[Order relation]\label{def:order-rel}
\index{order relation}
A relation $\leq$ on a set $E$ is a \textbf{(partial) order} if it is:
\begin{enumerate}[label=(\roman*)]
  \item \textbf{Reflexive}: $\forall x \in E,\; x \leq x$.
  \item \textbf{Antisymmetric}\index{antisymmetric}: $\forall x,y \in E,\;
    (x \leq y \text{ and } y \leq x) \implies x = y$.
  \item \textbf{Transitive}: $\forall x,y,z \in E,\;
    (x \leq y \text{ and } y \leq z) \implies x \leq z$.
\end{enumerate}
The order is \textbf{total}\index{total order} if additionally $\forall
x,y \in E$, either $x \leq y$ or $y \leq x$.
\end{definition}

\begin{example}
The usual ordering $\leq$ on $\R$ is a total order.  The inclusion relation
$\subset$ on subsets of a given set is a partial order that is \emph{not}
total in general (e.g., $\{1\} \not\subset \{2\}$ and $\{2\} \not\subset
\{1\}$).
\end{example}

%% ──────────────────────────────────────────────────────────────
\section{Common Errors in Logic}
\label{sec:common-errors-logic}

\begin{commonerrors}
\begin{enumerate}[label=\textbf{\arabic*.}]
  \item \textbf{Confusing a statement with its converse.}\quad
    ``$P \implies Q$'' is \emph{not} the same as ``$Q \implies P$''.
    Example: ``If it rains, the ground is wet'' does not mean ``If the
    ground is wet, it rained'' (someone might have used a hose).

  \item \textbf{Incorrect negation of quantifiers.}\quad
    The negation of ``$\forall x,\; P(x)$'' is \emph{not}
    ``$\forall x,\; \neg P(x)$''.  It is ``$\exists x,\; \neg P(x)$''.
    You only need \emph{one} counterexample to disprove a universal
    statement.

  \item \textbf{Using a specific example as a proof.}\quad
    Checking that a formula works for $n=1,2,3$ does \emph{not} prove it
    works for all~$n$.  You need induction or a general argument.

  \item \textbf{Assuming what you want to prove.}\quad
    In a direct proof of ``$A \implies B$'', you may assume $A$ but you
    must \emph{derive}~$B$.  You may not write $B$ and then ``verify'' it.

  \item \textbf{Confusing ``or'' with ``exclusive or''.}\quad
    In mathematics, ``$P$ or $Q$'' means ``at least one of $P,Q$ is true''
    and allows both to be true.
\end{enumerate}
\end{commonerrors}

%% ──────────────────────────────────────────────────────────────
\section{Exercises}
\label{sec:exercises-ch1}

\begin{exercise}[Truth tables]\label{exo:truth-tables}\easy
\index{truth table!exercise}
Construct truth tables for each of the following compound propositions:
\begin{enumerate}[label=(\alph*)]
  \item $P \implies (Q \lor R)$,
  \item $(P \land Q) \implies P$,
  \item $(P \implies Q) \iff (\neg Q \implies \neg P)$.
\end{enumerate}
For~(c), verify that the result is always true (a \emph{tautology}).
\end{exercise}

\begin{exercise}[Negations]\label{exo:negations}\easy
\index{negation!exercise}
Write the negation of each of the following statements:
\begin{enumerate}[label=(\alph*)]
  \item $\forall x \in \R,\; x^2 \geq 0$.
  \item $\exists n \in \N,\; n^2 = n$.
  \item $\forall \eps > 0,\; \exists \delta > 0,\; \forall x \in \R,\;
    \abs{x-a}<\delta \implies \abs{f(x)-f(a)}<\eps$.
  \item $\forall x \in \R,\; (x > 0 \implies \exists n \in \N,\;
    \frac{1}{n} < x)$.
\end{enumerate}
\end{exercise}

\begin{exercise}[Direct proof]\label{exo:direct1}\easy
Prove that the product of two odd integers is odd.
\end{exercise}

\begin{exercise}[Contrapositive]\label{exo:contra1}\medium
Let $n \in \Z$.  Prove that if $n^2$ is divisible by~$5$, then $n$ is
divisible by~$5$.  \textit{(Hint: consider the possible remainders of $n$
modulo~$5$.)}
\end{exercise}

\begin{exercise}[Contradiction]\label{exo:contra2}\medium
Prove that $\sqrt{3}$ is irrational.  \textit{(Model your proof on the
proof that $\sqrt{2}$ is irrational.)}
\end{exercise}

\begin{exercise}[Induction: sum of squares]\label{exo:ind-squares}\easy
Prove by induction that for every $n \geq 1$:
\[
  \sum_{k=1}^{n} k^2 = \frac{n(n+1)(2n+1)}{6}.
\]
\end{exercise}

\begin{exercise}[Induction: divisibility]\label{exo:ind-div}\medium
Prove that for every $n \geq 1$, $6$ divides $n^3 - n$.
\end{exercise}

\begin{exercise}[Induction: geometric sum]\label{exo:ind-geom}\easy
Let $q \neq 1$.  Prove by induction that for every $n \geq 0$:
\[
  \sum_{k=0}^{n} q^k = \frac{1 - q^{n+1}}{1-q}.
\]
\end{exercise}

\begin{exercise}[Sets: De Morgan]\label{exo:demorgan2}\medium
Prove the second De Morgan law: $(A \cap B)^c = A^c \cup B^c$.
\end{exercise}

\begin{exercise}[Functions]\label{exo:injection}\medium
Let $f : \R \to \R$ be defined by $f(x) = 2x + 3$.  Prove that $f$ is
bijective.  Find $f^{-1}$.
\end{exercise}

\begin{exercise}[Composition and injectivity]\label{exo:comp-inj}\hard
Let $f : A \to B$ and $g : B \to C$.  Prove:
\begin{enumerate}[label=(\alph*)]
  \item If $g \circ f$ is injective, then $f$ is injective.
  \item If $g \circ f$ is surjective, then $g$ is surjective.
  \item Give an example showing that $g \circ f$ can be bijective even if
    neither $f$ nor $g$ is bijective.
\end{enumerate}
\end{exercise}

\begin{exercise}[Equivalence relation]\label{exo:equiv-rel}\medium
Define a relation on $\Z \times \Z^*$ by $(a,b) \sim (c,d) \iff ad = bc$.
Prove that $\sim$ is an equivalence relation.  \textit{(This is the
construction of~$\Q$.)}
\end{exercise}

%% ──────────────────────────────────────────────────────────────
\section*{Chapter Summary}
\addcontentsline{toc}{section}{Chapter Summary}

\begin{chaptersummary}
\begin{itemize}[leftmargin=1.5em]
  \item A \textbf{proposition} is a statement that is either true or false.
  \item The logical connectives are $\neg$, $\land$, $\lor$, $\implies$, $\iff$.
  \item The \textbf{contrapositive} of $P \implies Q$ is $\neg Q \implies \neg P$,
    and it is logically equivalent.
  \item To \textbf{negate quantifiers}: flip $\forall \leftrightarrow \exists$
    and negate the innermost predicate.
  \item \textbf{Proof methods}: direct proof, contrapositive, contradiction,
    induction (ordinary and strong).
  \item \textbf{Set operations}: $\cup$, $\cap$, $\setminus$, complement,
    Cartesian product.  \textbf{De Morgan's laws} relate unions, intersections,
    and complements.
  \item A function is \textbf{injective} (one-to-one), \textbf{surjective}
    (onto), or \textbf{bijective} (both).
  \item An \textbf{equivalence relation} is reflexive, symmetric, and transitive.
    An \textbf{order relation} is reflexive, antisymmetric, and transitive.
\end{itemize}
\end{chaptersummary}


%%%%%%%%%%%%%%%%%%%%%%%%%%%%%%%%%%%%%%%%%%%%%%%%%%%%%%%%%%%%%%%%%%%%
%%%%%%%%%%%%%%%%%%%%%%%%%%%%%%%%%%%%%%%%%%%%%%%%%%%%%%%%%%%%%%%%%%%%
%%
%%  CHAPTER 2 : The Real Numbers — Axioms and Completeness
%%
%%%%%%%%%%%%%%%%%%%%%%%%%%%%%%%%%%%%%%%%%%%%%%%%%%%%%%%%%%%%%%%%%%%%
%%%%%%%%%%%%%%%%%%%%%%%%%%%%%%%%%%%%%%%%%%%%%%%%%%%%%%%%%%%%%%%%%%%%
\chapter{The Real Numbers --- Axioms and Completeness}
\label{ch:reals}
\index{real numbers}

The real numbers are the natural setting for analysis.  In this chapter we
state the axioms that define $\R$, and we single out the one axiom that
makes $\R$ special: the \emph{completeness axiom}.  From this single
principle, we will derive the Archimedean property, the density of $\Q$ in
$\R$, and the nested intervals theorem --- tools that we will use over and
over again in subsequent chapters.

%% ──────────────────────────────────────────────────────────────
\section{Why \texorpdfstring{$\Q$}{Q} Is Not Enough}
\label{sec:Q-not-enough}
\index{rational numbers!insufficiency of}

We proved in Theorem~\ref{thm:sqrt2} that $\sqrt{2} \notin \Q$.  This means
that the equation $x^2 = 2$ has no solution in $\Q$.  But the problem runs
deeper: the rational number line has ``holes'' in it.

Consider the set
\[
  A = \setcond{r \in \Q}{r > 0 \text{ and } r^2 < 2}.
\]
This set is nonempty ($1 \in A$) and is bounded above in $\Q$ (for instance,
$2$ is an upper bound, since $r \in A$ implies $r < 2$).  However, $A$ has
\emph{no least upper bound in~$\Q$}.  Intuitively, the least upper bound
``should be'' $\sqrt{2}$, but $\sqrt{2} \notin \Q$.

This is not just an isolated curiosity.  Many fundamental constructions in
analysis --- limits of sequences, the Intermediate Value Theorem, the
existence of integrals --- require the number system to have ``no holes.''
The rational numbers fail this requirement.  The real numbers $\R$ are
designed precisely to fill every such hole.

%% ──────────────────────────────────────────────────────────────
\section{The Ordered Field Axioms}
\label{sec:field-axioms}
\index{ordered field axioms}

The real numbers $\R$ satisfy three groups of axioms.

\begin{definition}[Field axioms]\label{def:field-axioms}
\index{field axioms}
$(\R, +, \cdot)$ is a \textbf{field}, meaning it satisfies the following
for all $a,b,c \in \R$:

\medskip
\textbf{Addition axioms:}
\begin{enumerate}[label=(A\arabic*)]
  \item \textbf{Associativity:} $(a+b)+c = a+(b+c)$.
  \item \textbf{Commutativity:} $a+b = b+a$.
  \item \textbf{Identity:} There exists $0 \in \R$ such that $a+0=a$.
  \item \textbf{Inverses:} For each $a$, there exists $-a \in \R$ such that
    $a+(-a)=0$.
\end{enumerate}

\textbf{Multiplication axioms:}
\begin{enumerate}[label=(M\arabic*)]
  \item \textbf{Associativity:} $(a \cdot b)\cdot c = a \cdot (b \cdot c)$.
  \item \textbf{Commutativity:} $a \cdot b = b \cdot a$.
  \item \textbf{Identity:} There exists $1 \in \R$, $1 \neq 0$, such that
    $a \cdot 1 = a$.
  \item \textbf{Inverses:} For each $a \neq 0$, there exists $a^{-1} \in \R$
    such that $a \cdot a^{-1} = 1$.
\end{enumerate}

\textbf{Distributivity:}
\begin{enumerate}[label=(D)]
  \item $a \cdot (b+c) = a\cdot b + a \cdot c$.
\end{enumerate}
\end{definition}

\begin{definition}[Order axioms]\label{def:order-axioms}
\index{order axioms}
$\R$ is equipped with a total order $\leq$ satisfying for all $a,b,c \in \R$:
\begin{enumerate}[label=(O\arabic*)]
  \item \textbf{Compatibility with addition:} $a \leq b \implies a+c \leq b+c$.
  \item \textbf{Compatibility with multiplication:} $a \leq b$ and $0 \leq c$
    $\implies$ $ac \leq bc$.
\end{enumerate}
A field equipped with such an order is called an \textbf{ordered field}%
\index{ordered field}.
\end{definition}

\begin{remark}
$\Q$ also satisfies all of the above axioms.  What distinguishes $\R$ from
$\Q$ is the \emph{completeness axiom}, which we state in
Section~\ref{sec:completeness}.
\end{remark}

%% ──────────────────────────────────────────────────────────────
\section{Absolute Value}
\label{sec:absolute-value}
\index{absolute value}

\begin{definition}[Absolute value]\label{def:abs-value}
For $x \in \R$, the \textbf{absolute value} of $x$ is
\[
  \abs{x} = \begin{cases} x & \text{if } x \geq 0, \\ -x & \text{if } x < 0. \end{cases}
\]
\end{definition}

\begin{proposition}[Properties of absolute value]\label{prop:abs-props}
\index{absolute value!properties}
For all $x,y \in \R$:
\begin{enumerate}[label=(\roman*)]
  \item $\abs{x} \geq 0$, and $\abs{x}=0 \iff x=0$.
  \item $\abs{-x}=\abs{x}$.
  \item $\abs{xy}=\abs{x}\abs{y}$.
  \item \label{item:abs-ineq} $-\abs{x} \leq x \leq \abs{x}$.
  \item $\abs{x} \leq a \iff -a \leq x \leq a$ \quad (for $a \geq 0$).
\end{enumerate}
\end{proposition}

\begin{proof}
We prove each property by case analysis on the sign of $x$ (and $y$ where
relevant).

\textbf{(i)} If $x \geq 0$, then $\abs{x}=x\geq 0$.  If $x<0$, then
$\abs{x}=-x>0$.  In either case $\abs{x}\geq 0$.  Moreover, $\abs{x}=0$
implies either $x=0$ (if $x\geq 0$) or $-x=0$ (if $x<0$), both giving $x=0$.
Conversely, $\abs{0}=0$.

\textbf{(ii)} If $x \geq 0$, then $-x \leq 0$, so
$\abs{-x}=-(-x)=x=\abs{x}$.  If $x < 0$, then $-x > 0$, so
$\abs{-x}=-x=\abs{x}$.

\textbf{(iii)} Consider the four sign combinations.  For instance, if
$x \geq 0$ and $y \geq 0$, then $xy \geq 0$ and
$\abs{xy}=xy=\abs{x}\abs{y}$.  The other cases are similar.

\textbf{(iv)} If $x \geq 0$, then $\abs{x}=x$, and $-\abs{x}=-x \leq 0
\leq x$.  If $x<0$, then $\abs{x}=-x$, and $x < 0 \leq -x = \abs{x}$,
while $-\abs{x}=x$.  In both cases, $-\abs{x}\leq x \leq \abs{x}$.

\textbf{(v)} ($\Rightarrow$) If $\abs{x}\leq a$, then by~(iv),
$-a \leq -\abs{x} \leq x \leq \abs{x} \leq a$.
($\Leftarrow$) If $-a\leq x \leq a$, then $x \leq a$ and $-x \leq a$.
Since $\abs{x}$ equals either $x$ or $-x$, we get $\abs{x}\leq a$.
\end{proof}

\begin{theorem}[Triangle inequality]\label{thm:triangle-ineq}
\index{triangle inequality}
For all $x,y \in \R$,
\[
  \abs{x+y} \leq \abs{x} + \abs{y}.
\]
\end{theorem}

\begin{proof}
By Proposition~\ref{prop:abs-props}\ref{item:abs-ineq}, we have
\[
  -\abs{x} \leq x \leq \abs{x}
  \qquad\text{and}\qquad
  -\abs{y} \leq y \leq \abs{y}.
\]
Adding these two inequalities:
\[
  -(\abs{x}+\abs{y}) \leq x + y \leq \abs{x}+\abs{y}.
\]
By the characterisation in Proposition~\ref{prop:abs-props}(v) (with
$a = \abs{x}+\abs{y} \geq 0$), this is equivalent to
$\abs{x+y} \leq \abs{x}+\abs{y}$.
\end{proof}

\begin{corollary}[Reverse triangle inequality]\label{cor:reverse-triangle}
\index{triangle inequality!reverse}
For all $x,y \in \R$,
\[
  \bigl|\abs{x}-\abs{y}\bigr| \leq \abs{x-y}.
\]
\end{corollary}

\begin{proof}
By the triangle inequality applied to $x = (x-y)+y$:
\[
  \abs{x} = \abs{(x-y)+y} \leq \abs{x-y}+\abs{y},
\]
so $\abs{x}-\abs{y} \leq \abs{x-y}$.  Swapping the roles of $x$ and $y$:
$\abs{y}-\abs{x} \leq \abs{y-x} = \abs{x-y}$.  Since both
$\abs{x}-\abs{y}$ and $\abs{y}-\abs{x}$ are $\leq \abs{x-y}$, we conclude
$\bigl|\abs{x}-\abs{y}\bigr| \leq \abs{x-y}$.
\end{proof}

%% ──────────────────────────────────────────────────────────────
\section{Bounded Sets, Supremum, and Infimum}
\label{sec:sup-inf}
\index{supremum}\index{infimum}\index{bounded set}

\begin{definition}[Bounded sets]\label{def:bounded-set}
Let $A \subset \R$ be a nonempty set.
\begin{enumerate}[label=(\roman*)]
  \item $A$ is \textbf{bounded above}\index{bounded above} if there exists
    $M \in \R$ such that $\forall a \in A,\; a \leq M$.  Such an $M$ is
    called an \textbf{upper bound}\index{upper bound} of~$A$.
  \item $A$ is \textbf{bounded below}\index{bounded below} if there exists
    $m \in \R$ such that $\forall a \in A,\; a \geq m$.  Such an $m$ is
    called a \textbf{lower bound}\index{lower bound} of~$A$.
  \item $A$ is \textbf{bounded}\index{bounded} if it is bounded both above
    and below.
\end{enumerate}
\end{definition}

\begin{definition}[Supremum and infimum]\label{def:sup-inf}
Let $A \subset \R$ be nonempty.
\begin{enumerate}[label=(\roman*)]
  \item The \textbf{supremum}\index{supremum} (or \textbf{least upper bound})
    of~$A$, denoted $\sup A$, is a number $s \in \R$ such that:
    \begin{enumerate}[label=(\alph*)]
      \item $s$ is an upper bound of $A$: $\forall a \in A,\; a \leq s$.
      \item $s$ is the \emph{least} such bound: if $M$ is any upper bound
        of $A$, then $s \leq M$.
    \end{enumerate}
  \item The \textbf{infimum}\index{infimum} (or \textbf{greatest lower
    bound}) of~$A$, denoted $\inf A$, is defined analogously: the greatest
    number $t$ that is a lower bound of~$A$.
\end{enumerate}
\end{definition}

\begin{remark}\label{rem:sup-unique}
If a supremum exists, it is unique.  Indeed, if $s_1$ and $s_2$ are both
suprema of~$A$, then $s_1$ is an upper bound so $s_2 \leq s_1$ (since
$s_2$ is the \emph{least} upper bound), and $s_2$ is an upper bound so
$s_1 \leq s_2$.  Hence $s_1 = s_2$.
\end{remark}

\begin{remark}[Sup vs.\ Max]\label{rem:sup-vs-max}
\index{maximum!vs supremum}
The \textbf{maximum} of~$A$ is an element $a_0 \in A$ such that
$a_0 \geq a$ for all $a \in A$.  If $\max A$ exists, then
$\sup A = \max A$.  But $\sup A$ may exist even when $\max A$ does not.
For instance, $A = (0,1)$ has $\sup A = 1$, but $1 \notin A$, so $A$ has
no maximum.
\end{remark}

\medskip

\begin{center}
\begin{tikzpicture}[>=Stealth, scale=1]
  % Number line
  \draw[thick,->] (-0.5,0) -- (7,0);
  % Set A
  \draw[very thick, blue] (1.5,0) -- (4.5,0);
  \fill[blue] (1.5,0) circle (2.5pt);
  \draw[blue, thick] (4.5,0) circle (2.5pt); % open circle
  \node[above, blue] at (3,0.15) {$A = [a, b)$};
  % inf/sup labels
  \draw[thick, red, dashed] (1.5,-0.4) -- (1.5,0.6);
  \node[below, red] at (1.5,-0.45) {$\inf A = a$};
  \draw[thick, red, dashed] (4.5,-0.4) -- (4.5,0.6);
  \node[below, red] at (4.5,-0.45) {$\sup A = b$};
  % Upper bounds region
  \draw[thick, ForestGreen, decorate, decoration={brace, amplitude=5pt,
    mirror}] (4.5,-1.1) -- (6.8,-1.1);
  \node[below, ForestGreen] at (5.65,-1.3) {\small upper bounds};
  % Lower bounds region
  \draw[thick, ForestGreen, decorate, decoration={brace, amplitude=5pt,
    mirror}] (-0.3,-1.1) -- (1.5,-1.1);
  \node[below, ForestGreen] at (0.6,-1.3) {\small lower bounds};
\end{tikzpicture}
\end{center}

\begin{theorem}[$\eps$-characterisation of the supremum]\label{thm:eps-sup}
\index{supremum!$\eps$-characterisation}
Let $A \subset \R$ be nonempty and bounded above, and let $s \in \R$.
Then $s = \sup A$ if and only if:
\begin{enumerate}[label=(\roman*)]
  \item $\forall a \in A,\; a \leq s$ \quad ($s$ is an upper bound), and
  \item $\forall \eps > 0,\; \exists a \in A,\; a > s - \eps$
    \quad (no number smaller than $s$ is an upper bound).
\end{enumerate}
\end{theorem}

\begin{proof}
\textbf{($\Rightarrow$)}\quad Suppose $s = \sup A$.  Condition~(i) holds
because the supremum is an upper bound.  For~(ii), let $\eps > 0$.  Then
$s - \eps < s$.  Since $s$ is the \emph{least} upper bound, $s-\eps$ is
\emph{not} an upper bound of~$A$.  Therefore, there exists $a \in A$ such
that $a > s - \eps$.

\textbf{($\Leftarrow$)}\quad Suppose $s$ satisfies (i) and~(ii).
By~(i), $s$ is an upper bound of $A$, so $\sup A \leq s$ (since $\sup A$
is the least upper bound).  We now show $s \leq \sup A$.  Suppose for
contradiction that $s > \sup A$.  Set $\eps = s - \sup A > 0$.  By~(ii),
there exists $a \in A$ with $a > s - \eps = \sup A$.  But this contradicts
the fact that $\sup A$ is an upper bound of~$A$.  Therefore $s \leq \sup A$,
and combined with $\sup A \leq s$, we get $s = \sup A$.
\end{proof}

\begin{remark}
There is an analogous characterisation for the infimum: $t = \inf A$ if and
only if $t$ is a lower bound of $A$ and for every $\eps > 0$ there exists
$a \in A$ with $a < t + \eps$.
\end{remark}

%% ──────────────────────────────────────────────────────────────
\section{The Completeness Axiom}
\label{sec:completeness}
\index{completeness axiom}

We now state the axiom that makes $\R$ fundamentally different from $\Q$.

\begin{definition}[Completeness axiom (Least Upper Bound Property)]\label{def:completeness}
\index{least upper bound property}
Every nonempty subset of $\R$ that is bounded above has a supremum in~$\R$.
\end{definition}

\begin{remark}
Equivalently (by considering $-A = \{-a : a \in A\}$), every nonempty
subset of $\R$ that is bounded below has an infimum in~$\R$.
\end{remark}

\medskip

\textbf{Why is this axiom so important?}\quad Without completeness, we
cannot guarantee that limits exist.  For instance, consider the sequence of
rational numbers
\[
  1,\; 1.4,\; 1.41,\; 1.414,\; 1.4142,\; \dots
\]
which gives better and better decimal approximations to $\sqrt{2}$.
Intuitively this sequence ``converges,'' but its limit is $\sqrt{2}$, which
is not in~$\Q$.  In $\Q$, this sequence has no limit.  The completeness
axiom ensures that such pathologies do not occur in~$\R$: every ``hole'' is
filled.

\medskip

\textbf{$\Q$ lacks the completeness property.}\quad We saw in
Section~\ref{sec:Q-not-enough} that the set
$A = \setcond{r \in \Q}{r>0,\; r^2<2}$ is nonempty and bounded above
in~$\Q$, but has no least upper bound in~$\Q$.  This shows that $\Q$ does
not satisfy the completeness axiom.  The real numbers $\R$ are, in a precise
sense, the ``completion'' of~$\Q$: we adjoin all the missing limits.

%% ──────────────────────────────────────────────────────────────
\section{Consequences of Completeness}
\label{sec:consequences}

\subsection{The Archimedean Property}
\label{subsec:archimedean}
\index{Archimedean property}

\begin{theorem}[Archimedean property]\label{thm:archimedean}
For every $x \in \R$, there exists $n \in \N$ such that $n > x$.

Equivalently: for every $\eps > 0$ and every $M > 0$, there exists
$n \in \N^*$ such that $n\eps > M$ (i.e., the multiples of any positive
number are unbounded).
\end{theorem}

\begin{proof}
We prove the first formulation.  Suppose for contradiction that there
exists $x \in \R$ such that $n \leq x$ for all $n \in \N$.  Then $\N$ is
bounded above by~$x$.  Since $\N$ is nonempty and bounded above, the
completeness axiom guarantees that $s = \sup \N$ exists in~$\R$.

Now consider $s - 1$.  Since $s-1 < s$ and $s$ is the \emph{least} upper
bound of~$\N$, the number $s-1$ is not an upper bound of~$\N$.  Therefore
there exists $n_0 \in \N$ with $n_0 > s-1$, i.e., $n_0 + 1 > s$.

But $n_0 + 1 \in \N$ (since $\N$ is closed under the successor operation),
and $n_0 + 1 > s$ contradicts the fact that $s$ is an upper bound of~$\N$.

This contradiction shows that $\N$ is not bounded above, proving the result.
\end{proof}

\begin{remark}
An immediate consequence: for every $\eps > 0$, there exists $n \in \N^*$
such that $\frac{1}{n} < \eps$.  Indeed, by the Archimedean property applied
to $x = \frac{1}{\eps}$, there exists $n > \frac{1}{\eps}$, hence
$\frac{1}{n} < \eps$.  This fact is used constantly in analysis.
\end{remark}

\subsection{Density of \texorpdfstring{$\Q$}{Q} in
\texorpdfstring{$\R$}{R}}
\label{subsec:density-Q}
\index{density!of $\Q$ in $\R$}

\begin{theorem}[Density of $\Q$ in $\R$]\label{thm:density-Q}
Between any two distinct real numbers, there exists a rational number.
Formally: if $a, b \in \R$ with $a < b$, then there exists $q \in \Q$ such
that $a < q < b$.
\end{theorem}

\begin{proof}
Since $a < b$, we have $b - a > 0$.  By the Archimedean property, there
exists $n \in \N^*$ such that
\begin{equation}\label{eq:density-n}
  n(b-a) > 1, \quad\text{i.e.,}\quad \frac{1}{n} < b - a.
\end{equation}
This means the intervals of length $\frac{1}{n}$ are narrow enough to ``fit
between'' $a$ and~$b$.

Now we need to find an integer $m$ such that $a < \frac{m}{n} < b$, i.e.,
$na < m < nb$.  We must show that there is an integer in the open interval
$(na, nb)$.

Again by the Archimedean property, there exists a positive integer $M_1$
with $M_1 > na$ and a positive integer $M_2$ with $M_2 > -na$.  So the set
of integers greater than $na$ is nonempty.  Let
\[
  m = \min\setcond{k \in \Z}{k > na}.
\]
(This minimum exists because the set of integers greater than $na$ is
nonempty and bounded below by $na$, and every nonempty set of integers
bounded below has a minimum.)

By definition of $m$, we have $m > na$, i.e., $\frac{m}{n} > a$.  Also,
$m - 1 \leq na$ (since $m$ is the \emph{smallest} integer exceeding $na$),
so $m \leq na + 1$.  Therefore
\[
  m \leq na + 1 < na + n(b-a) = nb,
\]
where we used \eqref{eq:density-n} in the strict inequality.  Hence
$\frac{m}{n} < b$.

Setting $q = \frac{m}{n} \in \Q$, we have $a < q < b$.
\end{proof}

\begin{corollary}[Density of irrationals]\label{cor:density-irr}
\index{density!of irrationals}
Between any two distinct real numbers, there exists an irrational number.
\end{corollary}

\begin{proof}
Let $a < b$.  Then $a - \sqrt{2} < b - \sqrt{2}$.  By the density of
$\Q$ (Theorem~\ref{thm:density-Q}), there exists $q \in \Q$ with
$a - \sqrt{2} < q < b - \sqrt{2}$.  Set $\alpha = q + \sqrt{2}$.  Then
$a < \alpha < b$.  We claim $\alpha$ is irrational: if $\alpha$ were
rational, then $\alpha - q = \sqrt{2}$ would be rational (as the difference
of two rationals), contradicting Theorem~\ref{thm:sqrt2}.  Therefore
$\alpha \notin \Q$.
\end{proof}

\subsection{The Floor Function}
\label{subsec:floor}
\index{floor function}

\begin{proposition}[Existence and uniqueness of the floor]\label{prop:floor}
For every $x \in \R$, there exists a unique integer $n \in \Z$ such that
$n \leq x < n+1$.  This integer is called the \textbf{floor} of $x$ and is
denoted $\lfloor x \rfloor$.
\end{proposition}

\begin{proof}
\textbf{Existence.}\quad By the Archimedean property, the set
$S = \setcond{k \in \Z}{k > x}$ is nonempty.  Similarly, the set
$\setcond{k \in \Z}{k \leq x}$ is nonempty (take $k$ to be an integer
less than or equal to $x$, which exists by applying the Archimedean
property to $-x$).

Let $m = \min S$ (the smallest integer strictly greater than $x$; this
exists because $S$ is a nonempty set of integers bounded below).  Set
$n = m - 1$.  Then $n = m - 1 \leq x$ (since $m$ is the smallest integer
$>x$, we have $m-1 \leq x$) and $n + 1 = m > x$.  So $n \leq x < n + 1$.

\textbf{Uniqueness.}\quad Suppose $n$ and $n'$ both satisfy
$n \leq x < n+1$ and $n' \leq x < n'+1$.  Then $n \leq x < n'+1$ gives
$n < n'+1$, so $n \leq n'$.  Similarly $n' \leq n$.  Hence $n = n'$.
\end{proof}

\subsection{The Nested Intervals Theorem}
\label{subsec:nested-intervals}
\index{nested intervals theorem}

\begin{theorem}[Nested intervals theorem]\label{thm:nested}
Let $(I_n)_{n \geq 1}$ be a sequence of closed bounded intervals
$I_n = [a_n, b_n]$ such that $I_{n+1} \subset I_n$ for all $n$ (i.e., the
intervals are nested).  Then
\[
  \bigcap_{n=1}^{\infty} I_n \neq \emptyset.
\]
Moreover, if $\lim_{n\to\infty}(b_n - a_n) = 0$, then the intersection
contains exactly one point.
\end{theorem}

\begin{center}
\begin{tikzpicture}[>=Stealth, scale=1]
  \foreach \i/\a/\b in {0/0.5/6, 1/1.2/5.2, 2/1.8/4.5, 3/2.3/4.0, 4/2.7/3.6} {
    \pgfmathsetmacro{\yy}{-0.7*\i}
    \draw[thick, blue!80!black] (\a,\yy) -- (\b,\yy);
    \fill[blue!80!black] (\a,\yy) circle (2pt);
    \fill[blue!80!black] (\b,\yy) circle (2pt);
    \node[left, font=\small] at (\a-0.1,\yy) {$a_{\the\numexpr\i+1}$};
    \node[right, font=\small] at (\b+0.1,\yy) {$b_{\the\numexpr\i+1}$};
    \node[right, font=\small] at (6.5,\yy) {$I_{\the\numexpr\i+1}$};
  }
  \node at (3.1, -4.1) {$\vdots$};
  % Intersection point
  \draw[thick, red, dashed] (3.1, 0.4) -- (3.1, -3.8);
  \node[above, red, font=\small] at (3.1, 0.45) {$\displaystyle\bigcap I_n$};
\end{tikzpicture}
\end{center}

\begin{proof}
Since $I_{n+1} \subset I_n$, we have $a_n \leq a_{n+1}$ and
$b_{n+1} \leq b_n$ for all~$n$.  Moreover, $a_n \leq b_n$ for all $n$.  In
fact, $a_m \leq b_n$ for \emph{all} $m,n \geq 1$ (indeed, if $m \leq n$
then $a_m \leq a_n \leq b_n$, and if $m > n$ then $a_m \leq b_m \leq b_n$).

\textbf{Step~1: The set of left endpoints is bounded above.}\quad
The set $A = \{a_n : n \geq 1\}$ is nonempty and bounded above by every
$b_n$ (by the observation above).  By the completeness axiom, $s = \sup A$
exists.

\textbf{Step~2: $s$ belongs to every interval.}\quad
For each $n$, $b_n$ is an upper bound of~$A$, so $s \leq b_n$ (since $s$
is the \emph{least} upper bound).  Also, $a_n \in A$ and $s = \sup A \geq
a_n$.  Therefore $a_n \leq s \leq b_n$, which means $s \in [a_n, b_n] =
I_n$ for every~$n$.  Hence $s \in \bigcap_{n=1}^{\infty} I_n$, and this
intersection is nonempty.

\textbf{Step~3: Uniqueness when lengths tend to zero.}\quad
Suppose $\lim_{n\to\infty}(b_n-a_n)=0$ and let $x,y \in \bigcap I_n$.
Then for every $n$, $a_n \leq x \leq b_n$ and $a_n \leq y \leq b_n$, so
$\abs{x-y} \leq b_n - a_n$.  Since $b_n - a_n \to 0$, we get
$\abs{x-y}=0$, i.e., $x=y$.
\end{proof}

%% ──────────────────────────────────────────────────────────────
\section{Historical Remarks}
\label{sec:historical-reals}

\index{Dedekind cuts}\index{Cantor, Georg}
The rigorous construction of the real numbers was one of the great
achievements of nineteenth-century mathematics.  Two independent approaches
emerged around 1872:

\begin{itemize}
  \item \textbf{Dedekind cuts} (Richard Dedekind, 1872): A real number is
    defined as a partition of $\Q$ into two nonempty sets $L$ and $R$ such
    that every element of $L$ is less than every element of~$R$, and $L$ has
    no greatest element.  The ``cut'' between $L$ and $R$ represents the
    real number.  Irrational numbers correspond to cuts where $R$ has no
    least element either.

  \item \textbf{Cauchy sequences} (Georg Cantor, 1872): A real number is
    defined as an equivalence class of Cauchy sequences of rational numbers,
    where two sequences are equivalent if their difference converges to zero.
    This approach makes the completeness of $\R$ immediate: every Cauchy
    sequence of real numbers converges (this will be proved in a later
    chapter).
\end{itemize}

Both constructions yield the same object: the unique (up to isomorphism)
complete ordered field $\R$.  In this course, we take the axiomatic approach:
we \emph{postulate} that $\R$ exists and satisfies the field axioms, the
order axioms, and the completeness axiom, and we derive everything else from
these.

%% ──────────────────────────────────────────────────────────────
\section{Common Errors}
\label{sec:common-errors-ch2}

\begin{commonerrors}
\begin{enumerate}[label=\textbf{\arabic*.}]
  \item \textbf{Confusing $\sup A$ with $\max A$.}\quad
    The supremum of a set \emph{need not belong to the set}.  For example,
    $\sup(0,1)=1$, but $1 \notin (0,1)$, so $(0,1)$ has no maximum.
    Always check whether the supremum is attained before calling it a maximum.

  \item \textbf{Claiming $\sup A \in A$ without justification.}\quad
    This is a frequent error.  If you need the supremum to be an element of
    the set, you must \emph{prove} it, or state that $A$ has a maximum.

  \item \textbf{Forgetting that $\sup$ applies to \emph{nonempty} bounded
    sets.}\quad The completeness axiom requires $A$ to be nonempty and
    bounded above.  The supremum of the empty set is not defined.

  \item \textbf{Applying the Archimedean property to non-positive numbers.}%
    \quad The statement ``for every $\eps>0$ there exists $n$ with
    $\frac{1}{n}<\eps$'' requires $\eps > 0$.  For $\eps \leq 0$ the
    statement is meaningless.

  \item \textbf{Assuming $\R$ is ``just $\Q$ with more digits.''}\quad
    The decimal expansion viewpoint is useful for intuition but dangerous
    for proofs.  The completeness axiom is the correct foundation.
\end{enumerate}
\end{commonerrors}

%% ──────────────────────────────────────────────────────────────
\section{Exercises}
\label{sec:exercises-ch2}

\begin{exercise}[Supremum practice]\label{exo:sup1}\easy
Find $\sup A$ and $\inf A$ for each of the following sets, and determine
whether the supremum (respectively, infimum) is a maximum (respectively,
minimum).
\begin{enumerate}[label=(\alph*)]
  \item $A = \left\{\frac{1}{n} : n \in \N^*\right\}$.
  \item $A = \left\{\frac{n}{n+1} : n \in \N\right\}$.
  \item $A = \left\{(-1)^n\left(1 + \frac{1}{n}\right) : n \in \N^*\right\}$.
\end{enumerate}
\end{exercise}

\begin{exercise}[Supremum of a union]\label{exo:sup-union}\medium
Let $A, B \subset \R$ be nonempty and bounded above.  Prove that
\[
  \sup(A \cup B) = \max\{\sup A,\, \sup B\}.
\]
\end{exercise}

\begin{exercise}[$\eps$-characterisation]\label{exo:eps-char}\medium
Let $A = \left\{1 - \frac{1}{n} : n \in \N^*\right\}$.  Use the
$\eps$-characterisation (Theorem~\ref{thm:eps-sup}) to prove that
$\sup A = 1$.
\end{exercise}

\begin{exercise}[Supremum of a sum]\label{exo:sup-sum}\hard
Let $A, B \subset \R$ be nonempty and bounded above.  Define
$A + B = \{a + b : a \in A,\; b \in B\}$.  Prove that
$\sup(A+B) = \sup A + \sup B$.
\end{exercise}

\begin{exercise}[Archimedean consequences]\label{exo:arch}\easy
Using the Archimedean property, prove:
\begin{enumerate}[label=(\alph*)]
  \item For every $x \in \R$, there exists $n \in \Z$ such that $n \leq x < n+1$.
  \item $\inf\left\{\frac{1}{n} : n \in \N^*\right\} = 0$.
\end{enumerate}
\end{exercise}

\begin{exercise}[Density]\label{exo:density}\medium
\begin{enumerate}[label=(\alph*)]
  \item Prove that between any two distinct real numbers, there exist
    infinitely many rational numbers.
  \item Prove that between any two distinct real numbers, there exist
    infinitely many irrational numbers.
\end{enumerate}
\end{exercise}

\begin{exercise}[Completeness is essential]\label{exo:completeness}\hard
Let $A = \setcond{r \in \Q}{r > 0,\; r^2 < 2}$.
\begin{enumerate}[label=(\alph*)]
  \item Prove that $A$ is nonempty and bounded above in $\Q$.
  \item Prove that $A$ has no maximum in $\Q$.  \textit{(Hint: given
    $r \in A$, show that $r' = r + \frac{2-r^2}{r+2}$ satisfies
    $r' \in A$ and $r' > r$.)}
  \item Explain why $A$ has no least upper bound in $\Q$.
\end{enumerate}
\end{exercise}

\begin{exercise}[Nested intervals]\label{exo:nested}\medium
Let $I_n = \left[1-\frac{1}{n},\; 2+\frac{1}{n}\right]$ for $n \geq 1$.
Determine $\bigcap_{n=1}^{\infty} I_n$.
\end{exercise}

\begin{exercise}[Nested intervals --- singleton]\label{exo:nested-single}\medium
Let $I_n = \left[\frac{n}{n+1},\; \frac{n+2}{n+1}\right]$ for $n \geq 1$.
\begin{enumerate}[label=(\alph*)]
  \item Show that the intervals are nested.
  \item Compute $b_n - a_n$ and its limit.
  \item Determine $\bigcap_{n=1}^{\infty} I_n$.
\end{enumerate}
\end{exercise}

\begin{exercise}[Triangle inequality applications]\label{exo:triangle}\easy
Let $a,b,c \in \R$.  Prove:
\begin{enumerate}[label=(\alph*)]
  \item $\abs{a-c} \leq \abs{a-b} + \abs{b-c}$.
  \item $\abs{a+b+c} \leq \abs{a}+\abs{b}+\abs{c}$.
\end{enumerate}
\end{exercise}

\begin{exercise}[Absolute value equations]\label{exo:abs-eq}\easy
Solve and prove your answer:
\begin{enumerate}[label=(\alph*)]
  \item $\abs{2x-3}=5$.
  \item $\abs{x-1}+\abs{x+1}\leq 4$.
\end{enumerate}
\end{exercise}

\begin{exercise}[Floor function properties]\label{exo:floor-props}\medium
Let $x \in \R$ and $n \in \Z$.  Prove:
\begin{enumerate}[label=(\alph*)]
  \item $\lfloor x + n \rfloor = \lfloor x \rfloor + n$.
  \item $\lfloor x \rfloor + \lfloor -x \rfloor =
    \begin{cases} 0 & \text{if } x \in \Z, \\ -1 & \text{if } x \notin \Z. \end{cases}$
\end{enumerate}
\end{exercise}

%% ──────────────────────────────────────────────────────────────
\section*{Chapter Summary}
\addcontentsline{toc}{section}{Chapter Summary}

\begin{chaptersummary}
\begin{itemize}[leftmargin=1.5em]
  \item $\Q$ is an ordered field but is \textbf{incomplete}: bounded sets
    may lack a supremum in $\Q$.
  \item $\R$ is the unique \textbf{complete ordered field}: every nonempty
    subset bounded above has a supremum in $\R$ (completeness axiom).
  \item The \textbf{absolute value} $\abs{x}$ satisfies the triangle
    inequality: $\abs{x+y}\leq\abs{x}+\abs{y}$.
  \item The \textbf{supremum} $\sup A$ is the least upper bound.  The
    $\eps$-characterisation: $s = \sup A$ iff $s$ is an upper bound and
    $\forall\eps>0$, $\exists a\in A$, $a>s-\eps$.
  \item \textbf{Archimedean property}: $\N$ is unbounded above.
    Equivalently, for all $\eps>0$, $\exists n \in \N^*$,
    $\frac{1}{n}<\eps$.
  \item \textbf{Density of $\Q$}: between any two reals lies a rational
    (and an irrational).
  \item The \textbf{floor function} $\lfloor x \rfloor$ is the unique
    integer with $\lfloor x\rfloor \leq x < \lfloor x\rfloor+1$.
  \item \textbf{Nested intervals theorem}: a decreasing sequence of closed
    bounded intervals has nonempty intersection; if the lengths tend to zero,
    the intersection is a singleton.
\end{itemize}
\end{chaptersummary}
%%%%%%%%%%%%%%%%%%%%%%%%%%%%%%%%%%%%%%%%%%%%%%%%%%%%%%%%%%%%
%% CHAPTER 3 — SEQUENCES
%%%%%%%%%%%%%%%%%%%%%%%%%%%%%%%%%%%%%%%%%%%%%%%%%%%%%%%%%%%%
\chapter{Sequences of Real Numbers}\label{ch:sequences}
\index{sequence}

\section*{Motivation: what does ``approaching'' mean?}

Consider the sequence $1, \frac{1}{2}, \frac{1}{3}, \frac{1}{4}, \ldots$ The terms get
\emph{closer and closer} to $0$, yet no term actually equals~$0$. But what does
``closer and closer'' mean \emph{precisely}? How close is close enough?

The answer---one of the great achievements of 19th-century mathematics---is the
$\eps$--$N$ definition of convergence. It replaces the vague intuition of
``approaching'' with a rigorous, quantified statement that can be proved or
disproved. Mastering this definition is the single most important step in a
first course in analysis.

%===================================================================
\section{Sequences: first definitions}\label{sec:seq-def}
%===================================================================

\begin{definition}[Sequence of real numbers]\label{def:sequence}
\index{sequence!definition}
A \textbf{sequence of real numbers} is a function
\[
   u \colon \N \longrightarrow \R, \qquad n \longmapsto u_n.
\]
We write $(u_n)_{n\in\N}$, or simply $(u_n)$, and call $u_n$ the
\textbf{$n$-th term}\index{sequence!term} (or \textbf{general term}) of the
sequence.
\end{definition}

\begin{remark}
Sometimes the index starts at $n=1$ or another integer; the theory is the same.
We may also write $(a_n)$, $(x_n)$, etc.
\end{remark}

\begin{example}\label{ex:first-sequences}
\leavevmode
\begin{enumerate}
\item $u_n = \dfrac{1}{n+1}$: the terms are $1, \frac{1}{2}, \frac{1}{3}, \ldots$
\item $v_n = (-1)^n$: the terms alternate between $1$ and $-1$.
\item $w_n = \dfrac{n}{n+1}$: the terms are $0, \frac{1}{2}, \frac{2}{3}, \frac{3}{4}, \ldots$
\item $a_n = 2^n$: the terms grow without bound.
\end{enumerate}
\end{example}

%===================================================================
\section{Convergence of a sequence}\label{sec:convergence}
%===================================================================
\index{sequence!convergence}\index{convergence!of a sequence}

The following definition is the cornerstone of the entire course. Read it
slowly, several times.

\begin{tcolorbox}[
   colback=blue!3!white,
   colframe=blue!60!black,
   title={\textbf{Definition (Convergence of a sequence)}},
   fonttitle=\bfseries,
   arc=2mm
]
\label{def:convergence}
Let $(u_n)_{n\in\N}$ be a sequence of real numbers. We say that $(u_n)$
\textbf{converges}\index{convergence} to $\ell \in \R$ if
\[
   \boxed{%
   \forall \eps > 0, \quad \exists N \in \N, \quad \forall n \geq N, \quad
   \abs{u_n - \ell} < \eps.}
\]
We then write $\displaystyle \lim_{n\to\infty} u_n = \ell$, or
$u_n \xrightarrow[n\to\infty]{} \ell$.

\medskip
\textbf{Parsing each quantifier:}
\begin{enumerate}
\item \emph{``$\forall \eps > 0$''}: for \textbf{every} tolerance $\eps$, no
      matter how small---you do not choose~$\eps$; it is given to you by an
      adversary.
\item \emph{``$\exists N \in \N$''}: there \textbf{exists} a rank~$N$ (which
      \textbf{may depend on~$\eps$})---your job is to find (or prove the
      existence of) such an~$N$.
\item \emph{``$\forall n \geq N$''}: for \textbf{all} indices $n$ at least as
      large as~$N$---that is, from rank~$N$ onward, with \textbf{no exceptions}.
\item \emph{``$\abs{u_n - \ell} < \eps$''}: the distance from $u_n$ to $\ell$
      is strictly less than~$\eps$.
\end{enumerate}

\medskip
In plain language: \emph{no matter how narrow a band $(\ell-\eps,\,\ell+\eps)$
you draw around the limit, all but finitely many terms of the sequence lie
inside that band.}
\end{tcolorbox}

\begin{remark}
The order of quantifiers is critical.
``$\forall \eps > 0,\; \exists N$'' means $N$ may depend on~$\eps$.
If the order were reversed, $N$ would have to work for \emph{every} $\eps$
simultaneously---an impossibly strong demand.
\end{remark}

\begin{definition}[Divergent sequence]\label{def:divergent}
\index{sequence!divergent}
A sequence that does not converge is called \textbf{divergent}.
\end{definition}

%-------------------------------------------------------------------
\subsection*{Visualization: the $\eps$-band}
%-------------------------------------------------------------------

\begin{center}
\begin{tikzpicture}[scale=1.0, >=stealth]
  % axes
  \draw[->] (-0.5,0) -- (11,0) node[right] {$n$};
  \draw[->] (0,-0.8) -- (0,3.5) node[above] {$u_n$};
  % limit line
  \draw[dashed, blue!70!black, thick] (-0.3,1.5) -- (10.8,1.5)
        node[right] {$\ell$};
  % epsilon band
  \fill[blue!8] (-0.3,0.8) rectangle (10.8,2.2);
  \draw[blue!60!black] (-0.3,2.2) -- (10.8,2.2)
        node[right, font=\small] {$\ell+\eps$};
  \draw[blue!60!black] (-0.3,0.8) -- (10.8,0.8)
        node[right, font=\small] {$\ell-\eps$};
  % N marker
  \draw[red, thick, dashed] (4,-0.3) -- (4,3.2);
  \node[red, below] at (4,-0.3) {$N$};
  % terms before N (some outside the band)
  \foreach \x/\y in {0.5/3.0, 1.2/0.4, 1.9/2.8, 2.6/0.6, 3.3/2.5} {
    \fill[red!70!black] (\x,\y) circle (2.2pt);
  }
  % terms after N (all inside the band)
  \foreach \x/\y in {4.5/1.7, 5.2/1.3, 5.9/1.6, 6.6/1.4, 7.3/1.55,
                      8.0/1.45, 8.7/1.52, 9.4/1.48, 10.1/1.50} {
    \fill[green!50!black] (\x,\y) circle (2.2pt);
  }
  % brace for eps
  \draw[decorate, decoration={brace, amplitude=5pt, mirror},
        thick, blue!60!black]
        (10.9,0.8) -- (10.9,1.5) node[midway, right=6pt] {$\eps$};
  \draw[decorate, decoration={brace, amplitude=5pt},
        thick, blue!60!black]
        (10.9,1.5) -- (10.9,2.2) node[midway, right=6pt] {$\eps$};
  % labels
  \node[font=\footnotesize, red!70!black] at (1.8,3.3)
        {finitely many terms outside};
  \node[font=\footnotesize, green!50!black] at (8,2.6)
        {all terms inside from rank $N$ on};
\end{tikzpicture}
\end{center}

%===================================================================
\section{First proofs of convergence}\label{sec:first-eps-proofs}
%===================================================================

\subsection{Strategy: scratch work then formal proof}

Every $\eps$--$N$ proof proceeds in two stages:

\begin{enumerate}
\item \textbf{Scratch work (rough computation).}  Assume $\abs{u_n - \ell} < \eps$
      and work backward to find what condition on~$n$ guarantees this.  This
      is \emph{not} part of the final proof---it is how you \emph{discover} the
      right~$N$.
\item \textbf{Formal proof.}  Start with ``Let $\eps > 0$'', choose $N$
      (the value you found in the scratch work), and verify that
      $n \geq N \implies \abs{u_n - \ell} < \eps$.
\end{enumerate}

%-------------------------------------------------------------------
\subsection{Worked example: $\lim_{n\to\infty}\frac{1}{n} = 0$}
\label{sec:ex-1-over-n}
%-------------------------------------------------------------------
\index{sequence!1/n@$1/n$}

\begin{example}[Scratch work]\label{ex:1n-scratch}
We want $\abs{\frac{1}{n} - 0} < \eps$, i.e.\ $\frac{1}{n} < \eps$.
This is equivalent to $n > \frac{1}{\eps}$.
So any integer $N > \frac{1}{\eps}$ will work. By the Archimedean property,
such an $N$ exists.
\end{example}

\begin{proposition}\label{prop:lim-1-over-n}
$\displaystyle \lim_{n\to\infty}\frac{1}{n} = 0$.
\end{proposition}

\begin{proof}
Let $\eps > 0$. By the Archimedean property of~$\R$, there exists
$N \in \N$ such that $N > \dfrac{1}{\eps}$, i.e.\ $\dfrac{1}{N} < \eps$.

Let $n \geq N$. Then $n \geq N > \dfrac{1}{\eps}$, so
\[
   \abs{\frac{1}{n} - 0}
   = \frac{1}{n}
   \leq \frac{1}{N}
   < \eps.
\]
This shows that for every $\eps > 0$ there exists $N \in \N$ such that
$n \geq N \implies \abs{\frac{1}{n} - 0} < \eps$. Therefore
$\lim_{n\to\infty}\frac{1}{n} = 0$.
\end{proof}

%-------------------------------------------------------------------
\subsection{Worked example: $\lim_{n\to\infty}\frac{n}{n+1} = 1$}
\label{sec:ex-n-over-n+1}
%-------------------------------------------------------------------

\begin{example}[Scratch work]\label{ex:nn1-scratch}
We compute
\[
   \abs{\frac{n}{n+1} - 1}
   = \abs{\frac{n - (n+1)}{n+1}}
   = \frac{1}{n+1}.
\]
We need $\dfrac{1}{n+1} < \eps$, i.e.\ $n+1 > \dfrac{1}{\eps}$, i.e.\
$n > \dfrac{1}{\eps} - 1$.

So it suffices to take $N$ any integer with $N > \dfrac{1}{\eps} - 1$
(or, more simply, $N > \dfrac{1}{\eps}$, since that is larger).
\end{example}

\begin{proposition}\label{prop:lim-n-over-n+1}
$\displaystyle \lim_{n\to\infty}\frac{n}{n+1} = 1$.
\end{proposition}

\begin{proof}
Let $\eps > 0$. Choose $N \in \N$ with $N > \dfrac{1}{\eps}$ (Archimedean
property). For every $n \geq N$,
\[
   \abs{\frac{n}{n+1} - 1}
   = \frac{1}{n+1}
   \leq \frac{1}{n}
   \leq \frac{1}{N}
   < \eps.
\]
Hence $\dfrac{n}{n+1} \to 1$ as $n \to \infty$.
\end{proof}

%-------------------------------------------------------------------
\subsection{Worked example: $((-1)^n)$ diverges}
\label{sec:ex-minus1-n}
%-------------------------------------------------------------------
\index{sequence!divergent!(-1)\^{}n@$(-1)^n$}

\begin{proposition}\label{prop:minus1n-diverges}
The sequence $((-1)^n)_{n\in\N}$ does not converge.
\end{proposition}

\begin{proof}
Suppose for contradiction that $((-1)^n)$ converges to some $\ell \in \R$.
Then, by the definition of convergence, applied with $\eps = \frac{1}{2}$,
there exists $N \in \N$ such that
\[
   \forall n \geq N, \quad \abs{(-1)^n - \ell} < \tfrac{1}{2}.
\]
In particular, taking $n = 2N$ (which is even and $\geq N$) and
$n = 2N+1$ (which is odd and $\geq N$):
\[
   \abs{1 - \ell} < \tfrac{1}{2}
   \quad \text{and} \quad
   \abs{-1 - \ell} < \tfrac{1}{2}.
\]
By the triangle inequality,
\[
   2 = \abs{1 - (-1)}
     = \abs{(1 - \ell) + (\ell - (-1))}
     \leq \abs{1 - \ell} + \abs{-1 - \ell}
     < \tfrac{1}{2} + \tfrac{1}{2} = 1.
\]
This gives $2 < 1$, a contradiction. Therefore $((-1)^n)$ diverges.
\end{proof}

%===================================================================
\section{Uniqueness of the limit}\label{sec:uniqueness}
%===================================================================
\index{limit!uniqueness}

\begin{theorem}[Uniqueness of limits]\label{thm:limit-unique}
If a sequence $(u_n)$ converges, then its limit is unique. That is, if
$u_n \to \ell$ and $u_n \to \ell'$, then $\ell = \ell'$.
\end{theorem}

\begin{proof}
Suppose $u_n \to \ell$ and $u_n \to \ell'$. Let $\eps > 0$.
\begin{itemize}
\item Since $u_n \to \ell$, there exists $N_1 \in \N$ such that
      $\forall n \geq N_1,\; \abs{u_n - \ell} < \frac{\eps}{2}$.
\item Since $u_n \to \ell'$, there exists $N_2 \in \N$ such that
      $\forall n \geq N_2,\; \abs{u_n - \ell'} < \frac{\eps}{2}$.
\end{itemize}
Set $N = \max(N_1, N_2)$. For every $n \geq N$, the triangle inequality gives
\[
   \abs{\ell - \ell'}
   = \abs{(\ell - u_n) + (u_n - \ell')}
   \leq \abs{u_n - \ell} + \abs{u_n - \ell'}
   < \frac{\eps}{2} + \frac{\eps}{2}
   = \eps.
\]
Thus $\abs{\ell - \ell'} < \eps$ for every $\eps > 0$.

Since $\abs{\ell - \ell'}$ is a fixed non-negative real number that is smaller
than every positive real number, it must equal~$0$. Hence $\ell = \ell'$.
\end{proof}

%===================================================================
\section{Bounded sequences}\label{sec:bounded-seq}
%===================================================================
\index{sequence!bounded}

\begin{definition}[Bounded sequence]\label{def:bounded-seq}
A sequence $(u_n)$ is \textbf{bounded}\index{bounded} if there exists $M \geq 0$
such that
\[
   \forall n \in \N, \quad \abs{u_n} \leq M.
\]
\end{definition}

\begin{theorem}\label{thm:convergent-bounded}
Every convergent sequence is bounded.
\end{theorem}

\begin{proof}
Let $(u_n)$ be a sequence converging to $\ell \in \R$. Apply the definition of
convergence with $\eps = 1$: there exists $N \in \N$ such that
\[
   \forall n \geq N, \quad \abs{u_n - \ell} < 1,
\]
which gives $\abs{u_n} \leq \abs{u_n - \ell} + \abs{\ell} < 1 + \abs{\ell}$
for all $n \geq N$.

The finitely many terms $u_0, u_1, \ldots, u_{N-1}$ form a finite set, so
\[
   M = \max\!\bigl(\abs{u_0}, \abs{u_1}, \ldots, \abs{u_{N-1}},\; 1 + \abs{\ell}\bigr)
\]
is well-defined. For every $n \in \N$, we have $\abs{u_n} \leq M$, so
$(u_n)$ is bounded.
\end{proof}

\begin{remark}
The converse is false: $((-1)^n)$ is bounded but not convergent.
\end{remark}

%===================================================================
\section{Squeeze theorem}\label{sec:squeeze}
%===================================================================
\index{squeeze theorem}\index{theorem!squeeze}

\begin{theorem}[Squeeze theorem / Sandwich theorem]\label{thm:squeeze}
Let $(u_n)$, $(v_n)$, $(w_n)$ be sequences such that
\[
   \forall n \in \N, \quad u_n \leq v_n \leq w_n,
\]
and suppose $u_n \to \ell$ and $w_n \to \ell$ for the same $\ell \in \R$.
Then $v_n \to \ell$.
\end{theorem}

\begin{proof}
Let $\eps > 0$.
\begin{itemize}
\item Since $u_n \to \ell$, there exists $N_1$ such that
      $\forall n \geq N_1$, $\abs{u_n - \ell} < \eps$, i.e.\
      $\ell - \eps < u_n < \ell + \eps$.
\item Since $w_n \to \ell$, there exists $N_2$ such that
      $\forall n \geq N_2$, $\abs{w_n - \ell} < \eps$, i.e.\
      $\ell - \eps < w_n < \ell + \eps$.
\end{itemize}
Let $N = \max(N_1, N_2)$. For every $n \geq N$,
\[
   \ell - \eps < u_n \leq v_n \leq w_n < \ell + \eps,
\]
so $\abs{v_n - \ell} < \eps$. Since $\eps > 0$ was arbitrary, $v_n \to \ell$.
\end{proof}

\begin{example}\label{ex:squeeze-application}
Let $u_n = \dfrac{\sin n}{n}$. Since $-1 \leq \sin n \leq 1$ for all $n$,
\[
   -\frac{1}{n} \leq \frac{\sin n}{n} \leq \frac{1}{n}.
\]
Both $-\frac{1}{n} \to 0$ and $\frac{1}{n} \to 0$, so by the squeeze theorem,
$\dfrac{\sin n}{n} \to 0$.
\end{example}

%===================================================================
\section{Algebra of limits}\label{sec:algebra-limits}
%===================================================================
\index{algebra of limits}

\begin{theorem}[Algebra of limits]\label{thm:algebra-limits}
Let $(u_n)$ and $(v_n)$ be convergent sequences with $u_n \to \ell$ and
$v_n \to \ell'$. Then:
\begin{enumerate}
\item\label{it:sum} $(u_n + v_n) \to \ell + \ell'$.
\item\label{it:scalar} For every $\lambda \in \R$, $(\lambda\, u_n) \to \lambda\,\ell$.
\item\label{it:product} $(u_n \cdot v_n) \to \ell \cdot \ell'$.
\item\label{it:quotient} If $\ell' \neq 0$, then $(u_n / v_n) \to \ell / \ell'$
      (and $v_n \neq 0$ for all $n$ large enough).
\end{enumerate}
\end{theorem}

\begin{proof}[Proof of \ref{it:sum} (sum)]
Let $\eps > 0$.
\begin{itemize}
\item Since $u_n \to \ell$, there exists $N_1$ with
      $\forall n \geq N_1,\; \abs{u_n - \ell} < \frac{\eps}{2}$.
\item Since $v_n \to \ell'$, there exists $N_2$ with
      $\forall n \geq N_2,\; \abs{v_n - \ell'} < \frac{\eps}{2}$.
\end{itemize}
Let $N = \max(N_1, N_2)$. For $n \geq N$,
\[
   \abs{(u_n + v_n) - (\ell + \ell')}
   = \abs{(u_n - \ell) + (v_n - \ell')}
   \leq \abs{u_n - \ell} + \abs{v_n - \ell'}
   < \frac{\eps}{2} + \frac{\eps}{2} = \eps.
\]
Hence $(u_n + v_n) \to \ell + \ell'$.
\end{proof}

\begin{proof}[Proof of \ref{it:product} (product)]
We use the identity
\[
   u_n v_n - \ell\ell'
   = (u_n - \ell)\,v_n + \ell\,(v_n - \ell').
\]
Since $(v_n)$ converges, it is bounded (Theorem~\ref{thm:convergent-bounded}):
there exists $M > 0$ with $\abs{v_n} \leq M$ for all~$n$.

Let $\eps > 0$.
\begin{itemize}
\item There exists $N_1$ with $\forall n \geq N_1,\;
      \abs{u_n - \ell} < \dfrac{\eps}{2M}$.
\item There exists $N_2$ with $\forall n \geq N_2,\;
      \abs{v_n - \ell'} < \dfrac{\eps}{2(\abs{\ell} + 1)}$.
\end{itemize}
(We write $\abs{\ell}+1$ instead of $\abs{\ell}$ to avoid division by zero
when $\ell = 0$.)

Let $N = \max(N_1, N_2)$. For $n \geq N$,
\begin{align*}
   \abs{u_n v_n - \ell\ell'}
   &\leq \abs{u_n - \ell}\,\abs{v_n} + \abs{\ell}\,\abs{v_n - \ell'} \\
   &< \frac{\eps}{2M}\cdot M + \abs{\ell}\cdot\frac{\eps}{2(\abs{\ell}+1)} \\
   &= \frac{\eps}{2} + \frac{\abs{\ell}}{\abs{\ell}+1}\cdot\frac{\eps}{2} \\
   &\leq \frac{\eps}{2} + \frac{\eps}{2} = \eps.
\end{align*}
Hence $(u_n v_n) \to \ell\ell'$.
\end{proof}

%===================================================================
\section{Monotone sequences and the Monotone Convergence Theorem}
\label{sec:monotone}
%===================================================================
\index{monotone sequence}\index{monotone convergence theorem}

\begin{definition}[Monotone sequence]\label{def:monotone}
A sequence $(u_n)$ is
\begin{itemize}
\item \textbf{increasing}\index{increasing sequence} if
      $u_n \leq u_{n+1}$ for all $n \in \N$;
\item \textbf{strictly increasing} if $u_n < u_{n+1}$ for all $n \in \N$;
\item \textbf{decreasing}\index{decreasing sequence} if
      $u_n \geq u_{n+1}$ for all $n \in \N$;
\item \textbf{strictly decreasing} if $u_n > u_{n+1}$ for all $n \in \N$.
\end{itemize}
A sequence that is either increasing or decreasing is called \textbf{monotone}.
\end{definition}

\begin{theorem}[Monotone Convergence Theorem]\label{thm:MCT}
\index{theorem!monotone convergence}
Every bounded monotone sequence of real numbers converges.

More precisely:
\begin{enumerate}
\item If $(u_n)$ is increasing and bounded above, then
      $\displaystyle\lim_{n\to\infty} u_n = \sup_{n\in\N} u_n$.
\item If $(u_n)$ is decreasing and bounded below, then
      $\displaystyle\lim_{n\to\infty} u_n = \inf_{n\in\N} u_n$.
\end{enumerate}
\end{theorem}

\begin{proof}
We prove~(1); the proof of~(2) is analogous.

Let $(u_n)$ be increasing and bounded above. The set
$A = \{u_n : n \in \N\}$ is non-empty and bounded above, so by the
completeness axiom of~$\R$, $\ell = \sup A$ exists in~$\R$.

We claim that $u_n \to \ell$. Let $\eps > 0$. Since $\ell$ is the
\emph{least} upper bound of~$A$ and $\ell - \eps < \ell$, the number
$\ell - \eps$ is \emph{not} an upper bound of~$A$. Hence there exists
$N \in \N$ with
\[
   u_N > \ell - \eps.
\]
Since $(u_n)$ is increasing and $\ell$ is an upper bound, for every
$n \geq N$ we have
\[
   \ell - \eps < u_N \leq u_n \leq \ell < \ell + \eps.
\]
Therefore $\abs{u_n - \ell} < \eps$ for all $n \geq N$. This proves
$u_n \to \ell = \sup A$.
\end{proof}

\begin{remark}
This theorem relies on the completeness of $\R$. It fails in $\Q$: the
sequence $1, 1.4, 1.41, 1.414, \ldots$ of truncations of $\sqrt{2}$ is
increasing and bounded in $\Q$ but does not converge in $\Q$.
\end{remark}

%===================================================================
\section{Subsequences and the Bolzano--Weierstrass theorem}
\label{sec:subsequences}
%===================================================================
\index{subsequence}\index{Bolzano--Weierstrass theorem}

\begin{definition}[Subsequence]\label{def:subsequence}
Let $(u_n)_{n\in\N}$ be a sequence and let $\varphi\colon \N \to \N$ be a
strictly increasing function. The sequence $(u_{\varphi(n)})_{n\in\N}$ is
called a \textbf{subsequence}\index{subsequence} of $(u_n)$.
\end{definition}

\begin{remark}
Common notation: $(u_{n_k})_{k\in\N}$ where $(n_k)$ is a strictly increasing
sequence of natural numbers. For instance, $(u_{2k})$ is the subsequence of
even-indexed terms.
\end{remark}

\begin{lemma}\label{lem:extraction-bound}
If $\varphi\colon \N \to \N$ is strictly increasing, then
$\varphi(n) \geq n$ for all $n \in \N$.
\end{lemma}

\begin{proof}
By induction. $\varphi(0) \geq 0$ since $\varphi(0) \in \N$.
If $\varphi(k) \geq k$, then
$\varphi(k+1) > \varphi(k) \geq k$, so $\varphi(k+1) \geq k + 1$.
\end{proof}

\begin{proposition}\label{prop:subsequence-limit}
If $u_n \to \ell$, then every subsequence of $(u_n)$ also converges to~$\ell$.
\end{proposition}

\begin{proof}
Let $\eps > 0$. There exists $N$ with $\abs{u_n - \ell} < \eps$ for $n \geq N$.
For $k \geq N$, Lemma~\ref{lem:extraction-bound} gives $\varphi(k) \geq k \geq N$,
so $\abs{u_{\varphi(k)} - \ell} < \eps$.
\end{proof}

\begin{definition}[Limit point / cluster value]\label{def:limit-point}
\index{limit point}\index{cluster value|see{limit point}}%
A real number $\ell$ is a \textbf{limit point} (or \textbf{cluster value},
French: \emph{valeur d'adh\'erence}\index{valeur d'adh\'erence|see{limit point}})
of a sequence $(u_n)$ if some subsequence of $(u_n)$ converges to~$\ell$.
\end{definition}

\begin{example}
The sequence $u_n = (-1)^n$ has exactly two limit points: $1$ (via the
subsequence $u_{2k}$) and $-1$ (via $u_{2k+1}$).
\end{example}

\begin{theorem}[Bolzano--Weierstrass]\label{thm:BW}
\index{theorem!Bolzano--Weierstrass}
Every bounded sequence of real numbers has a convergent subsequence
(equivalently, at least one limit point).
\end{theorem}

\begin{proof}
Let $(u_n)$ be a bounded sequence: there exist $a, b \in \R$ with
$a \leq u_n \leq b$ for all~$n$. We construct a convergent subsequence by
\emph{interval bisection}.

Set $[a_0, b_0] = [a, b]$. At each stage, consider the midpoint
$m = \frac{a_k + b_k}{2}$ and the two halves $[a_k, m]$ and $[m, b_k]$.
At least one of these halves contains \emph{infinitely many} terms of $(u_n)$
(since the union of two finite sets is finite, and $[a_k, b_k]$ contains
infinitely many terms). Choose such a half and call it $[a_{k+1}, b_{k+1}]$.

This gives a nested sequence of closed intervals
$[a_0, b_0] \supset [a_1, b_1] \supset \cdots$ with
$b_k - a_k = \dfrac{b - a}{2^k} \to 0$.

Now extract a subsequence: pick $n_0$ such that $u_{n_0} \in [a_0, b_0]$.
Having chosen $n_0 < n_1 < \cdots < n_k$ with $u_{n_j} \in [a_j, b_j]$,
the interval $[a_{k+1}, b_{k+1}]$ contains infinitely many terms, so we can
find $n_{k+1} > n_k$ with $u_{n_{k+1}} \in [a_{k+1}, b_{k+1}]$.

By the nested interval property, $\displaystyle \bigcap_{k=0}^{\infty}
[a_k, b_k] = \{\ell\}$ for some $\ell \in \R$.

Since $u_{n_k} \in [a_k, b_k]$ and $b_k - a_k \to 0$, we have
\[
   \abs{u_{n_k} - \ell} \leq b_k - a_k = \frac{b-a}{2^k} \to 0.
\]
Hence $u_{n_k} \to \ell$, and $(u_{n_k})$ is the desired convergent subsequence.
\end{proof}

%===================================================================
\section{Cauchy sequences}\label{sec:cauchy}
%===================================================================
\index{Cauchy sequence}

\begin{definition}[Cauchy sequence]\label{def:cauchy}
A sequence $(u_n)$ is a \textbf{Cauchy sequence}\index{Cauchy sequence} if
\[
   \forall \eps > 0, \quad \exists N \in \N, \quad
   \forall p, q \geq N, \quad \abs{u_p - u_q} < \eps.
\]
In words: the terms become arbitrarily close to \emph{each other}---without
needing to know the limit.
\end{definition}

\begin{theorem}\label{thm:conv-implies-cauchy}
Every convergent sequence is a Cauchy sequence.
\end{theorem}

\begin{proof}
Let $u_n \to \ell$ and let $\eps > 0$. There exists $N$ with
$\abs{u_n - \ell} < \frac{\eps}{2}$ for all $n \geq N$.
For $p, q \geq N$,
\[
   \abs{u_p - u_q}
   \leq \abs{u_p - \ell} + \abs{\ell - u_q}
   < \frac{\eps}{2} + \frac{\eps}{2} = \eps.
\]
\end{proof}

\begin{theorem}[Cauchy criterion for convergence in $\R$]
\label{thm:cauchy-criterion}
\index{Cauchy criterion}
A sequence of real numbers converges if and only if it is a Cauchy sequence.
\end{theorem}

\begin{proof}
($\Rightarrow$) Theorem~\ref{thm:conv-implies-cauchy}.

($\Leftarrow$) Let $(u_n)$ be a Cauchy sequence. We proceed in three steps.

\medskip
\textit{Step 1: $(u_n)$ is bounded.}
Take $\eps = 1$. There exists $N$ with $\abs{u_p - u_q} < 1$ for all
$p, q \geq N$. In particular, $\abs{u_n} \leq \abs{u_n - u_N} + \abs{u_N}
< 1 + \abs{u_N}$ for $n \geq N$. Setting
$M = \max(\abs{u_0}, \ldots, \abs{u_{N-1}}, 1 + \abs{u_N})$,
we get $\abs{u_n} \leq M$ for all~$n$.

\medskip
\textit{Step 2: $(u_n)$ has a convergent subsequence.}
Since $(u_n)$ is bounded, the Bolzano--Weierstrass theorem
(Theorem~\ref{thm:BW}) gives a subsequence $(u_{n_k})$ converging to some
$\ell \in \R$.

\medskip
\textit{Step 3: $(u_n)$ converges to $\ell$.}
Let $\eps > 0$.
\begin{itemize}
\item Since $(u_n)$ is Cauchy, there exists $N_1$ with
      $\abs{u_p - u_q} < \frac{\eps}{2}$ for $p, q \geq N_1$.
\item Since $u_{n_k} \to \ell$, there exists $K$ with
      $\abs{u_{n_K} - \ell} < \frac{\eps}{2}$ and $n_K \geq N_1$.
\end{itemize}
For $n \geq N_1$,
\[
   \abs{u_n - \ell}
   \leq \abs{u_n - u_{n_K}} + \abs{u_{n_K} - \ell}
   < \frac{\eps}{2} + \frac{\eps}{2} = \eps.
\]
Hence $u_n \to \ell$.
\end{proof}

%===================================================================
\section{Recursive sequences}\label{sec:recursive-seq}
%===================================================================
\index{sequence!recursive}

Many sequences are defined by a \textbf{recurrence relation}\index{recurrence relation}
$u_{n+1} = f(u_n)$ for some function $f\colon \R \to \R$ and an initial value~$u_0$.

\subsection*{General method}
\begin{enumerate}
\item \textbf{Existence and well-definedness.} The sequence is defined by
      induction---always valid.
\item \textbf{Conjecture monotonicity and bounds.} Compute a few terms; guess
      that $(u_n)$ is, say, increasing and bounded above.
\item \textbf{Prove monotonicity} (typically by induction):
      show $u_{n+1} - u_n \geq 0$ or $u_{n+1}/u_n \geq 1$.
\item \textbf{Prove the bound} (again by induction).
\item \textbf{Conclude convergence} by the Monotone Convergence Theorem
      (\ref{thm:MCT}).
\item \textbf{Identify the limit.} If $u_n \to \ell$ and $f$ is continuous at
      $\ell$, then passing to the limit in $u_{n+1} = f(u_n)$ gives
      $\ell = f(\ell)$. Solve this equation for~$\ell$, then use the
      bounds/monotonicity to select the correct root.
\end{enumerate}

\begin{example}\label{ex:recursive-sqrt2}
Let $u_0 = 1$ and $u_{n+1} = \dfrac{1}{2}\!\left(u_n + \dfrac{2}{u_n}\right)$.

\textbf{Claim:} $(u_n)$ converges to $\sqrt{2}$.

\medskip\noindent
\textbf{Step 1: $u_n > 0$ for all $n$.}
By induction: $u_0 = 1 > 0$, and if $u_n > 0$ then $u_{n+1}$ is the average
of two positive numbers.

\textbf{Step 2: $u_n \geq \sqrt{2}$ for all $n \geq 1$.}
By AM--GM: $u_{n+1} = \frac{1}{2}(u_n + \frac{2}{u_n}) \geq
\sqrt{u_n \cdot \frac{2}{u_n}} = \sqrt{2}$.

\textbf{Step 3: $(u_n)_{n\geq 1}$ is decreasing.}
$u_{n+1} - u_n = \frac{1}{2}\!\left(\frac{2}{u_n} - u_n\right)
= \frac{2 - u_n^2}{2u_n} \leq 0$
since $u_n^2 \geq 2$ by Step~2.

\textbf{Step 4: Convergence.} $(u_n)_{n\geq 1}$ is decreasing and bounded
below by $\sqrt{2}$, hence converges by the MCT.

\textbf{Step 5: Identify the limit.} Let $\ell = \lim u_n$.
Then $\ell = \frac{1}{2}(\ell + 2/\ell)$, giving $2\ell = \ell + 2/\ell$,
so $\ell = 2/\ell$, hence $\ell^2 = 2$. Since $\ell \geq \sqrt{2} > 0$,
we get $\ell = \sqrt{2}$.
\end{example}

\begin{example}\label{ex:recursive-contraction}
Let $u_0 = 0$ and $u_{n+1} = \dfrac{u_n + 3}{4}$.

\textbf{Claim:} $(u_n)$ converges to $1$.

\medskip\noindent
\textbf{Step 1: $0 \leq u_n \leq 1$ for all $n$.}
Induction: $u_0 = 0 \in [0,1]$.
If $0 \leq u_n \leq 1$, then $u_{n+1} = \frac{u_n+3}{4} \in
[\frac{3}{4}, 1] \subset [0,1]$.

\textbf{Step 2: $(u_n)$ is increasing.}
$u_{n+1} - u_n = \frac{u_n+3}{4} - u_n = \frac{3 - 3u_n}{4}
= \frac{3(1 - u_n)}{4} \geq 0$ since $u_n \leq 1$.

\textbf{Step 3: Convergence and limit.}
By MCT, $(u_n)$ converges. If $\ell = \lim u_n$, then
$\ell = \frac{\ell+3}{4}$, giving $4\ell = \ell + 3$, so $\ell = 1$.
\end{example}

%===================================================================
\section{Adjacent sequences}\label{sec:adjacent}
%===================================================================
\index{adjacent sequences}

\begin{definition}[Adjacent sequences]\label{def:adjacent}
Two sequences $(u_n)$ and $(v_n)$ are called \textbf{adjacent} if:
\begin{enumerate}
\item $(u_n)$ is increasing,
\item $(v_n)$ is decreasing,
\item $\lim_{n\to\infty}(v_n - u_n) = 0$.
\end{enumerate}
\end{definition}

\begin{theorem}\label{thm:adjacent}
If $(u_n)$ and $(v_n)$ are adjacent, then both converge to a common limit
$\ell$ satisfying $u_n \leq \ell \leq v_n$ for all~$n$.
\end{theorem}

\begin{proof}
Since $v_n - u_n \to 0$ and $v_n - u_n$ is eventually non-negative
(indeed, one can check $u_n \leq v_n$ for all~$n$: the sequence
$(v_n - u_n)$ is decreasing, and its limit is $0$, so $v_n - u_n \geq 0$).

Hence $u_n \leq v_n$ for all~$n$, so $(u_n)$ is increasing and bounded above
by $v_0$, and $(v_n)$ is decreasing and bounded below by $u_0$. By the MCT,
both converge: $u_n \to \ell_1$, $v_n \to \ell_2$. Then
$\ell_2 - \ell_1 = \lim(v_n - u_n) = 0$, so $\ell_1 = \ell_2 = \ell$.
Finally, $u_n \leq \ell$ (since $(u_n)$ is increasing with limit $\ell$)
and $v_n \geq \ell$ (since $(v_n)$ is decreasing with limit $\ell$).
\end{proof}

%===================================================================
\section{Additional visualizations}\label{sec:seq-tikz}
%===================================================================

\begin{center}
\begin{tikzpicture}[scale=0.85, >=stealth]
  % -- Monotone bounded sequence --
  \begin{scope}[xshift=0cm]
    \draw[->] (-0.3,0) -- (7,0) node[right] {$n$};
    \draw[->] (0,-0.3) -- (0,3.5) node[above] {};
    \draw[dashed, red!70!black] (-0.2,3) -- (6.8,3)
          node[right, font=\small] {$\sup = \ell$};
    \node[below, font=\small\bfseries] at (3.3,-0.7)
          {Monotone bounded};
    \foreach \x/\y in {0.3/0.5, 1/1.2, 1.7/1.8, 2.4/2.2,
                        3.1/2.5, 3.8/2.7, 4.5/2.82, 5.2/2.9,
                        5.9/2.95, 6.5/2.97} {
      \fill[blue!70!black] (\x,\y) circle (1.8pt);
    }
  \end{scope}
  % -- Cauchy sequence --
  \begin{scope}[xshift=9cm]
    \draw[->] (-0.3,0) -- (7,0) node[right] {$n$};
    \draw[->] (0,-0.3) -- (0,3.5) node[above] {};
    \draw[dashed, red!70!black] (-0.2,2) -- (6.8,2)
          node[right, font=\small] {$\ell$};
    \node[below, font=\small\bfseries] at (3.3,-0.7)
          {Cauchy sequence};
    \foreach \x/\y in {0.3/3.2, 1/0.8, 1.7/2.8, 2.4/1.4,
                        3.1/2.4, 3.8/1.7, 4.5/2.15, 5.2/1.92,
                        5.9/2.04, 6.5/1.99} {
      \fill[green!50!black] (\x,\y) circle (1.8pt);
    }
  \end{scope}
\end{tikzpicture}
\end{center}

%===================================================================
\section{Common errors}\label{sec:seq-errors}
%===================================================================

\begin{tcolorbox}[
   colback=red!3!white,
   colframe=red!60!black,
   title={\textbf{Common Errors --- Sequences}},
   fonttitle=\bfseries,
   arc=2mm
]
\begin{enumerate}
\item \textbf{Confusing bounded and convergent.}
      Bounded $\not\!\!\implies$ convergent (e.g.\ $(-1)^n$).
      Convergent $\implies$ bounded, but not conversely.

\item \textbf{Choosing $\eps$ in a proof.}
      In an $\eps$--$N$ proof, you do \emph{not} pick a specific value of
      $\eps$. You write ``Let $\eps > 0$'' and find $N$ that works for
      \emph{that particular} (but arbitrary) $\eps$.

\item \textbf{$N$ must not depend on $n$.}
      The rank $N$ depends on $\eps$ but not on $n$. Writing
      ``take $N = n + \ldots$'' is invalid.

\item \textbf{Assuming the limit exists.}
      To prove convergence, you cannot start by assuming $\lim u_n = \ell$
      and then compute $\ell$ (circular reasoning). Use monotonicity +
      boundedness or the Cauchy criterion first.

\item \textbf{Divergent $\neq$ tends to $+\infty$.}
      The sequence $(-1)^n$ diverges but does \emph{not} tend to $\pm\infty$.

\item \textbf{Forgetting that a Cauchy sequence argument requires $\R$.}
      The implication Cauchy $\implies$ convergent uses the completeness of $\R$.
      It is false in $\Q$.
\end{enumerate}
\end{tcolorbox}

%===================================================================
\section{Exercises}\label{sec:seq-exercises}
%===================================================================

\begin{exercise}[$\bigstar$]\label{exo:seq1}
Prove from the $\eps$--$N$ definition that
$\displaystyle\lim_{n\to\infty}\frac{2n+1}{n+3} = 2$.
\end{exercise}

\begin{exercise}[$\bigstar$]\label{exo:seq2}
Prove that $\displaystyle\lim_{n\to\infty}\frac{n^2}{n^2+1} = 1$.
\end{exercise}

\begin{exercise}[$\bigstar$]\label{exo:seq3}
Let $u_n = \dfrac{3n+(-1)^n}{n}$. Prove that $u_n \to 3$.
\end{exercise}

\begin{exercise}[$\bigstar$]\label{exo:seq4}
Using the squeeze theorem, find $\displaystyle\lim_{n\to\infty}\frac{\cos(n^2)}{n+1}$.
\end{exercise}

\begin{exercise}[$\bigstar\bigstar$]\label{exo:seq5}
Let $u_n = \left(1 + \dfrac{1}{n}\right)^n$. Show that $(u_n)$ is increasing
and bounded above by~$3$.
\textit{(Hint: use the binomial theorem and compare term by term.)}
\end{exercise}

\begin{exercise}[$\bigstar\bigstar$]\label{exo:seq6}
Let $u_0 = 2$ and $u_{n+1} = \dfrac{1}{2}\!\left(u_n + \dfrac{3}{u_n}\right)$.
Show that $(u_n)_{n\geq 1}$ is decreasing, bounded below by $\sqrt{3}$, and
find its limit.
\end{exercise}

\begin{exercise}[$\bigstar\bigstar$]\label{exo:seq7}
Let $u_n = \displaystyle\sum_{k=0}^{n}\frac{1}{k!}$.
Show that $(u_n)$ is increasing and bounded above.
\textit{(Hint: compare $k!$ with $2^{k-1}$ for $k \geq 1$.)}
\end{exercise}

\begin{exercise}[$\bigstar\bigstar$]\label{exo:seq8}
Let $(u_n)$ be a Cauchy sequence and let $(u_{n_k})$ be a subsequence
converging to $\ell$. Prove directly (without using
Theorem~\ref{thm:cauchy-criterion}) that $u_n \to \ell$.
\end{exercise}

\begin{exercise}[$\bigstar\bigstar$]\label{exo:seq9}
Let $(u_n)$ and $(v_n)$ be adjacent sequences with $u_0 = 0$, $v_0 = 1$,
$u_{n+1} = \dfrac{u_n + v_n}{2}$, $v_{n+1} = \dfrac{u_{n+1} + v_n}{2}$.
Show that $(u_n)$ and $(v_n)$ are adjacent and find their common limit.
\end{exercise}

\begin{exercise}[$\bigstar\bigstar$]\label{exo:seq10}
Prove that the sequence $u_n = \displaystyle\sum_{k=1}^{n}\frac{1}{k} - \ln n$
is decreasing and bounded below by $0$. \textit{(This limit is the Euler--Mascheroni
constant $\gamma \approx 0.5772$.)}
\end{exercise}

\begin{exercise}[$\bigstar\bigstar\bigstar$]\label{exo:seq11}
Let $u_0 \in (0,1)$ and $u_{n+1} = u_n(1 - u_n)$. Show that $u_n \to 0$.
Then show that $n\,u_n \to 1$.
\end{exercise}

\begin{exercise}[$\bigstar\bigstar\bigstar$]\label{exo:seq12}
(\textit{Stolz--Ces\`aro lemma.})
Let $(b_n)$ be a strictly increasing sequence with $b_n \to +\infty$, and
let $(a_n)$ be a sequence such that
$\dfrac{a_{n+1} - a_n}{b_{n+1} - b_n} \to \ell \in \R$.
Prove that $\dfrac{a_n}{b_n} \to \ell$.
\end{exercise}

\begin{exercise}[$\bigstar\bigstar\bigstar$]\label{exo:seq13}
Prove that every sequence in $\R$ has a monotone subsequence.
\textit{(Hint: consider ``peak'' indices $n$ where $u_n \geq u_m$ for all $m \geq n$.)}
\end{exercise}

%===================================================================
\section*{Chapter summary}
%===================================================================

\begin{tcolorbox}[
   colback=green!3!white,
   colframe=green!50!black,
   title={\textbf{Chapter 3 --- Summary}},
   fonttitle=\bfseries,
   arc=2mm
]
\begin{itemize}
\item A sequence converges to $\ell$ if:
      $\forall \eps > 0,\; \exists N,\; \forall n \geq N,\;
      \abs{u_n - \ell} < \eps$.
\item The limit, when it exists, is unique.
\item Convergent $\implies$ bounded (converse is false).
\item Squeeze theorem: $u_n \leq v_n \leq w_n$ with $u_n, w_n \to \ell$
      $\implies$ $v_n \to \ell$.
\item Algebra of limits: sums, products, quotients of convergent sequences
      behave as expected.
\item Monotone convergence theorem: increasing + bounded above $\implies$
      converges to $\sup$.
\item Bolzano--Weierstrass: bounded $\implies$ has a convergent subsequence.
\item Cauchy criterion: $(u_n)$ converges $\iff$ $(u_n)$ is Cauchy (in $\R$).
\item Recursive sequences: prove monotonicity + bounds, then pass to the
      limit.
\item Adjacent sequences converge to a common limit.
\end{itemize}
\end{tcolorbox}


%%%%%%%%%%%%%%%%%%%%%%%%%%%%%%%%%%%%%%%%%%%%%%%%%%%%%%%%%%%%
%% CHAPTER 4 — NUMERICAL SERIES
%%%%%%%%%%%%%%%%%%%%%%%%%%%%%%%%%%%%%%%%%%%%%%%%%%%%%%%%%%%%
\chapter{Numerical Series}\label{ch:series}
\index{series}

\section*{Motivation: adding infinitely many numbers}

Can we make sense of an infinite sum such as
$1 + \frac{1}{2} + \frac{1}{4} + \frac{1}{8} + \cdots$?
What about $1 + \frac{1}{2} + \frac{1}{3} + \frac{1}{4} + \cdots$?

Intuitively, the first sum ``should'' equal~$2$, while the second grows
without bound. To make these ideas rigorous, we reduce infinite sums to the
theory of sequences developed in Chapter~\ref{ch:sequences}: the infinite sum
is the \emph{limit of partial sums}.

%===================================================================
\section{Definitions}\label{sec:series-def}
%===================================================================

\begin{definition}[Series, partial sums]\label{def:series}
\index{series!definition}\index{partial sums}
Let $(a_n)_{n\geq 0}$ be a sequence of real numbers. The \textbf{series}
$\displaystyle\sum a_n$ (or $\sum_{n=0}^{\infty} a_n$) is the sequence of
\textbf{partial sums}
\[
   S_n = \sum_{k=0}^{n} a_k = a_0 + a_1 + \cdots + a_n.
\]
The series \textbf{converges}\index{series!convergence} if the sequence
$(S_n)$ converges, and in that case we write
\[
   \sum_{n=0}^{\infty} a_n = \lim_{n\to\infty} S_n.
\]
If $(S_n)$ diverges, the series \textbf{diverges}\index{series!divergence}.
The number $a_n$ is called the \textbf{general term} of the series.
\end{definition}

\begin{definition}[Remainder]\label{def:remainder}
\index{series!remainder}
If $\sum a_n$ converges with sum $S$, the \textbf{$n$-th remainder} is
\[
   R_n = S - S_n = \sum_{k=n+1}^{\infty} a_k.
\]
Note that $R_n \to 0$ as $n \to \infty$.
\end{definition}

%===================================================================
\section{Necessary condition for convergence}\label{sec:necessary-cond}
%===================================================================
\index{series!necessary condition}

\begin{theorem}[Necessary condition]\label{thm:necessary}
If the series $\sum a_n$ converges, then $a_n \to 0$.
\end{theorem}

\begin{proof}
If $\sum a_n$ converges, then $(S_n)$ converges to some $S \in \R$. Since
$a_n = S_n - S_{n-1}$ for $n \geq 1$, and both $S_n \to S$ and $S_{n-1} \to S$,
the algebra of limits gives
\[
   a_n = S_n - S_{n-1} \to S - S = 0. \qedhere
\]
\end{proof}

\begin{tcolorbox}[
   colback=orange!5!white,
   colframe=orange!70!black,
   title={\textbf{WARNING: the converse is false!}},
   fonttitle=\bfseries,
   arc=2mm
]
\index{harmonic series}
The condition $a_n \to 0$ is \textbf{necessary but not sufficient}.
The \textbf{harmonic series} $\displaystyle\sum_{n=1}^{\infty}\frac{1}{n}$
has $\frac{1}{n} \to 0$, yet it \textbf{diverges}.
\end{tcolorbox}

\begin{proposition}[Divergence of the harmonic series]\label{prop:harmonic}
The series $\displaystyle\sum_{n=1}^{\infty}\frac{1}{n}$ diverges.
\end{proposition}

\begin{proof}
We show that the partial sums are unbounded. Group the terms:
\begin{align*}
   S_1 &= 1, \\
   S_2 &= 1 + \tfrac{1}{2}, \\
   S_4 &= 1 + \tfrac{1}{2} + \bigl(\tfrac{1}{3}+\tfrac{1}{4}\bigr)
         \geq 1 + \tfrac{1}{2} + \bigl(\tfrac{1}{4}+\tfrac{1}{4}\bigr)
         = 1 + \tfrac{1}{2} + \tfrac{1}{2}, \\
   S_8 &= S_4 + \bigl(\tfrac{1}{5}+\tfrac{1}{6}+\tfrac{1}{7}+\tfrac{1}{8}\bigr)
         \geq S_4 + 4 \cdot \tfrac{1}{8}
         = S_4 + \tfrac{1}{2}.
\end{align*}
In general, $S_{2^k} \geq 1 + \frac{k}{2}$, which tends to $+\infty$.
Hence $(S_n)$ is unbounded, so the harmonic series diverges.
\end{proof}

%===================================================================
\section{Geometric series}\label{sec:geometric}
%===================================================================
\index{geometric series}\index{series!geometric}

\begin{theorem}[Geometric series]\label{thm:geometric}
Let $a \in \R$ with $a \neq 0$ and $r \in \R$. The geometric series
$\displaystyle\sum_{n=0}^{\infty} a\,r^n$ converges if and only if $\abs{r} < 1$,
and in that case
\[
   \sum_{n=0}^{\infty} a\,r^n = \frac{a}{1-r}.
\]
\end{theorem}

\begin{proof}
The partial sum is
\[
   S_n = a\sum_{k=0}^{n} r^k
       = a \cdot \frac{1 - r^{n+1}}{1 - r}
       \quad\text{for } r \neq 1.
\]
(This identity is proved by multiplying both sides by $1-r$.)

\begin{itemize}
\item If $\abs{r} < 1$, then $r^{n+1} \to 0$ (one can show $\abs{r}^n \to 0$
      using, e.g., $\abs{r}^n \leq \frac{1}{(1+\delta)^n} \to 0$ where
      $\abs{r} = \frac{1}{1+\delta}$ with $\delta > 0$).
      Hence $S_n \to \dfrac{a}{1 - r}$.
\item If $\abs{r} \geq 1$, then $\abs{a\,r^n} = \abs{a}\,\abs{r}^n
      \not\to 0$, so the necessary condition (Theorem~\ref{thm:necessary})
      fails, and the series diverges.
\item If $r = 1$, then $S_n = a(n+1) \to \pm\infty$. \qedhere
\end{itemize}
\end{proof}

\begin{example}\label{ex:geometric-examples}
\leavevmode
\begin{enumerate}
\item $\displaystyle\sum_{n=0}^{\infty}\frac{1}{2^n}
      = \frac{1}{1-\frac{1}{2}} = 2$.
\item $\displaystyle\sum_{n=0}^{\infty}\frac{(-1)^n}{3^n}
      = \frac{1}{1-(-\frac{1}{3})} = \frac{1}{\frac{4}{3}} = \frac{3}{4}$.
\end{enumerate}
\end{example}

%===================================================================
\section{Telescoping series}\label{sec:telescoping}
%===================================================================
\index{series!telescoping}\index{telescoping series}

\begin{proposition}\label{prop:telescoping}
Let $(b_n)$ be a sequence. The series
$\displaystyle\sum_{n=0}^{\infty}(b_n - b_{n+1})$ converges if and only if
$(b_n)$ converges, and in that case
\[
   \sum_{n=0}^{\infty}(b_n - b_{n+1}) = b_0 - \lim_{n\to\infty} b_n.
\]
\end{proposition}

\begin{proof}
The partial sum telescopes:
$S_n = \sum_{k=0}^{n}(b_k - b_{k+1}) = b_0 - b_{n+1}$.
Hence $S_n$ converges iff $b_{n+1}$ converges, and the formula follows.
\end{proof}

\begin{example}\label{ex:telescoping1}
$\displaystyle\sum_{n=1}^{\infty}\frac{1}{n(n+1)}
= \sum_{n=1}^{\infty}\left(\frac{1}{n} - \frac{1}{n+1}\right)
= 1 - \lim_{n\to\infty}\frac{1}{n+1} = 1$.

(Use partial fractions: $\frac{1}{n(n+1)} = \frac{1}{n} - \frac{1}{n+1}$.)
\end{example}

\begin{example}\label{ex:telescoping2}
$\displaystyle\sum_{n=1}^{\infty}\frac{1}{n(n+2)}
= \frac{1}{2}\sum_{n=1}^{\infty}\left(\frac{1}{n}-\frac{1}{n+2}\right)$.

The partial sum is
$\frac{1}{2}\!\left(1 + \frac{1}{2} - \frac{1}{n+1} - \frac{1}{n+2}\right)
\to \frac{1}{2}\cdot\frac{3}{2} = \frac{3}{4}$.
\end{example}

%===================================================================
\section{Series with non-negative terms}\label{sec:positive-series}
%===================================================================
\index{series!positive terms}

For a series $\sum a_n$ with $a_n \geq 0$, the partial sums $(S_n)$ form an
increasing sequence. Therefore, by the MCT, \textbf{$\sum a_n$ converges
if and only if $(S_n)$ is bounded above}.

\subsection{Comparison test}\label{subsec:comparison}
\index{comparison test}

\begin{theorem}[Comparison test]\label{thm:comparison}
Let $\sum a_n$ and $\sum b_n$ be series with $0 \leq a_n \leq b_n$ for all
$n$ (or at least for all $n$ large enough). Then:
\begin{enumerate}
\item If $\sum b_n$ converges, then $\sum a_n$ converges.
\item If $\sum a_n$ diverges, then $\sum b_n$ diverges.
\end{enumerate}
\end{theorem}

\begin{proof}
Let $S_n = \sum_{k=0}^{n}a_k$ and $T_n = \sum_{k=0}^{n}b_k$.
Since $a_k \leq b_k$, we have $S_n \leq T_n$ for all~$n$.

(1) If $\sum b_n$ converges, then $(T_n)$ is bounded above by
$T = \sum_{k=0}^{\infty}b_k$. Hence $S_n \leq T_n \leq T$ for all $n$,
so $(S_n)$ is increasing and bounded above, hence converges by the MCT.

(2) is the contrapositive of~(1).
\end{proof}

\subsection{Ratio test (d'Alembert)}\label{subsec:ratio}
\index{ratio test}\index{d'Alembert's ratio test}

\begin{theorem}[d'Alembert's ratio test]\label{thm:ratio-test}
Let $\sum a_n$ be a series with $a_n > 0$ for all $n$. Suppose
\[
   L = \lim_{n\to\infty}\frac{a_{n+1}}{a_n}
\]
exists (possibly $+\infty$). Then:
\begin{enumerate}
\item If $L < 1$, the series converges.
\item If $L > 1$ (including $L = +\infty$), the series diverges.
\item If $L = 1$, the test is inconclusive.
\end{enumerate}
\end{theorem}

\begin{proof}
\textit{Case $L < 1$.}
Choose $r$ with $L < r < 1$. Since $\frac{a_{n+1}}{a_n} \to L < r$,
there exists $N$ such that $\frac{a_{n+1}}{a_n} \leq r$ for all $n \geq N$.
Then, by induction,
\[
   a_{N+k} \leq a_N \cdot r^k \quad \text{for all } k \geq 0.
\]
Since $\sum_{k=0}^{\infty} a_N\, r^k = \frac{a_N}{1-r} < \infty$ (geometric
series with ratio $r < 1$), the comparison test gives the convergence of
$\sum_{n=N}^{\infty} a_n$, hence of $\sum_{n=0}^{\infty} a_n$.

\medskip
\textit{Case $L > 1$.}
There exists $N$ with $\frac{a_{n+1}}{a_n} > 1$ for $n \geq N$, so
$(a_n)_{n\geq N}$ is increasing and $a_n \geq a_N > 0$ for $n \geq N$.
Hence $a_n \not\to 0$, so the series diverges by the necessary condition.
\end{proof}

\begin{remark}
When $L = 1$: $\sum\frac{1}{n}$ diverges and $\sum\frac{1}{n^2}$ converges,
yet both have ratio limit~$1$. So the test tells us nothing.
\end{remark}

\subsection{Root test (Cauchy)}\label{subsec:root}
\index{root test}\index{Cauchy's root test}

\begin{theorem}[Cauchy's root test]\label{thm:root-test}
Let $\sum a_n$ be a series with $a_n \geq 0$. Suppose
\[
   L = \lim_{n\to\infty}\sqrt[n]{a_n}
\]
exists (possibly $+\infty$). Then:
\begin{enumerate}
\item If $L < 1$, the series converges.
\item If $L > 1$, the series diverges.
\item If $L = 1$, the test is inconclusive.
\end{enumerate}
\end{theorem}

\begin{proof}
\textit{Case $L < 1$.}
Choose $r$ with $L < r < 1$. There exists $N$ with $\sqrt[n]{a_n} \leq r$
for all $n \geq N$, i.e.\ $a_n \leq r^n$. Since $\sum r^n$ converges
(geometric with ratio $r < 1$), the comparison test gives convergence.

\medskip
\textit{Case $L > 1$.}
There exists $N$ with $\sqrt[n]{a_n} > 1$ for $n \geq N$, i.e.\
$a_n > 1$ for $n \geq N$. Hence $a_n \not\to 0$, and the series diverges.
\end{proof}

\subsection{Integral test}\label{subsec:integral-test}
\index{integral test}

\begin{theorem}[Integral test]\label{thm:integral-test}
Let $f\colon [1, +\infty) \to [0,+\infty)$ be a continuous, decreasing
function. Then the series $\sum_{n=1}^{\infty} f(n)$ and the improper integral
$\int_1^{\infty} f(x)\,dx$ either both converge or both diverge.
\end{theorem}

The proof uses the comparison between the sum and upper/lower Riemann sums;
we omit it here and refer to the chapter on integration.

\begin{example}[$p$-series / Riemann series]\label{ex:p-series}
\index{p-series@$p$-series}\index{Riemann series}
For $\alpha \in \R$, the series $\displaystyle\sum_{n=1}^{\infty}\frac{1}{n^\alpha}$
converges if and only if $\alpha > 1$.

\medskip
Apply the integral test with $f(x) = x^{-\alpha}$:
\[
   \int_1^{N}\frac{dx}{x^\alpha}
   = \begin{cases}
      \dfrac{N^{1-\alpha}-1}{1-\alpha} & \text{if } \alpha \neq 1, \\[6pt]
      \ln N & \text{if } \alpha = 1.
   \end{cases}
\]
As $N \to \infty$: if $\alpha > 1$, the integral converges (to $\frac{1}{\alpha-1}$);
if $\alpha \leq 1$, it diverges.
\end{example}

%===================================================================
\section{Absolute and conditional convergence}\label{sec:absolute}
%===================================================================
\index{absolute convergence}\index{conditional convergence}

\begin{definition}\label{def:absolute-convergence}
A series $\sum a_n$ is \textbf{absolutely convergent}\index{absolutely convergent}
if the series $\sum \abs{a_n}$ converges. It is \textbf{conditionally convergent}
if $\sum a_n$ converges but $\sum \abs{a_n}$ diverges.
\end{definition}

\begin{theorem}\label{thm:abs-implies-conv}
If $\sum a_n$ is absolutely convergent, then $\sum a_n$ is convergent.
\end{theorem}

\begin{proof}
We use the Cauchy criterion for series. Let $\eps > 0$.
Since $\sum\abs{a_n}$ converges, the sequence of its partial sums is Cauchy:
there exists $N$ such that for all $q > p \geq N$,
\[
   \sum_{k=p+1}^{q}\abs{a_k} < \eps.
\]
Then, by the triangle inequality,
\[
   \abs{\sum_{k=p+1}^{q} a_k}
   \leq \sum_{k=p+1}^{q}\abs{a_k}
   < \eps.
\]
Hence the partial sums of $\sum a_n$ are Cauchy, so $\sum a_n$ converges
(by the Cauchy criterion for sequences in $\R$).
\end{proof}

\begin{remark}
The converse is false: the alternating harmonic series is the standard
counterexample (see below).
\end{remark}

%===================================================================
\section{Alternating series: the Leibniz criterion}\label{sec:alternating}
%===================================================================
\index{alternating series}\index{Leibniz criterion}

\begin{definition}
An \textbf{alternating series} has the form $\sum_{n=0}^{\infty}(-1)^n b_n$
where $b_n \geq 0$ for all~$n$.
\end{definition}

\begin{theorem}[Leibniz criterion / Alternating series test]
\label{thm:leibniz}
\index{theorem!Leibniz}
Let $(b_n)_{n\geq 0}$ be a sequence such that:
\begin{enumerate}
\item $(b_n)$ is decreasing: $b_{n+1} \leq b_n$ for all~$n$,
\item $b_n \to 0$.
\end{enumerate}
Then the series $\sum_{n=0}^{\infty}(-1)^n b_n$ converges. Moreover, if $S$
denotes its sum, then
\[
   \abs{S - S_n} \leq b_{n+1}
   \quad\text{and}\quad
   S_{2n} \leq S \leq S_{2n+1} \;\text{(or vice versa)}.
\]
\end{theorem}

\begin{proof}
Consider the even and odd partial sums separately.

\textit{Even partial sums.}
\[
   S_{2n+2} = S_{2n} + b_{2n+1} - b_{2n+2}.
\]
Since $b_{2n+1} \geq b_{2n+2}$ (decreasing), we have $S_{2n+2} \geq S_{2n}$.
Thus $(S_{2n})$ is increasing.

\textit{Odd partial sums.}
\[
   S_{2n+1} = S_{2n-1} - b_{2n} + b_{2n+1}.
\]
Since $b_{2n} \geq b_{2n+1}$, we have $S_{2n+1} \leq S_{2n-1}$.
Thus $(S_{2n+1})$ is decreasing.

\textit{Comparison.}
$S_{2n+1} = S_{2n} + (-1)^{2n+1}b_{2n+1} = S_{2n} - b_{2n+1} \leq S_{2n}$?
Wait --- actually $S_{2n+1} = S_{2n} + (-1)^{2n+1}b_{2n+1}$. Since the sign
convention is $(-1)^n b_n$, we have:
\[
   S_{2n+1} - S_{2n} = (-1)^{2n+1}b_{2n+1} = -b_{2n+1} \leq 0,
\]
so $S_{2n} \geq S_{2n+1}$. But also
\[
   S_{2n} - S_{2n-1} = (-1)^{2n}b_{2n} = b_{2n} \geq 0,
\]
so $S_{2n} \geq S_{2n-1}$.

Combining, for all~$n$:
\[
   S_0 \geq S_1, \qquad
   S_{2n} \leq S_{2n+2}, \qquad
   S_{2n+1} \geq S_{2n+3}, \qquad
   S_{2n+1} \leq S_{2n}.
\]
Actually, let us be more careful. We have:
\begin{itemize}
\item $(S_{2n})$ is increasing (shown above) and bounded above by $S_1 + b_1 = S_0 = b_0$.
      More precisely, $S_{2n} \leq S_0 = b_0$ since
      $S_{2n} = b_0 - (b_1 - b_2) - (b_3 - b_4) - \cdots - (b_{2n-1} - b_{2n}) \leq b_0$.
\item $(S_{2n+1})$ is decreasing and bounded below by $0$ since
      $S_{2n+1} = (b_0 - b_1) + (b_2 - b_3) + \cdots + (b_{2n} - b_{2n+1}) \geq 0$.
\item $S_{2n} \geq S_{2n+1}$ for all $n$ (since $S_{2n} - S_{2n+1} = b_{2n+1} \geq 0$).
\end{itemize}

Wait, let us recheck the sign: $S_{2n+1} = S_{2n} + (-1)^{2n+1}b_{2n+1} = S_{2n} - b_{2n+1}$,
so indeed $S_{2n} - S_{2n+1} = b_{2n+1} \geq 0$.

By the monotone convergence theorem:
\begin{itemize}
\item $(S_{2n})$ converges to some $\ell_{\text{even}}$,
\item $(S_{2n+1})$ converges to some $\ell_{\text{odd}}$.
\end{itemize}
Since $S_{2n} - S_{2n+1} = b_{2n+1} \to 0$, we have
$\ell_{\text{even}} = \ell_{\text{odd}} =: S$.

Since $(S_{2n})$ and $(S_{2n+1})$ are the two subsequences that together
exhaust all terms, $S_n \to S$.

\medskip
\textit{Error bound.}
Since $S_{2n} \leq S \leq S_{2n+1}$ (the even subsequence increases to $S$,
the odd one decreases to $S$---actually, let us verify: $(S_{2n})$ increases
to $S$, so $S_{2n} \leq S$; $(S_{2n+1})$ decreases to $S$, so $S_{2n+1} \geq S$.
Hence $S_{2n} \leq S \leq S_{2n+1}$.)

For any $n$:
\[
   \abs{S - S_n} \leq \abs{S_{n+1} - S_n} = b_{n+1}. \qedhere
\]
\end{proof}

\begin{example}[Alternating harmonic series]\label{ex:alt-harmonic}
\index{alternating harmonic series}
The series $\displaystyle\sum_{n=1}^{\infty}\frac{(-1)^{n+1}}{n}
= 1 - \frac{1}{2} + \frac{1}{3} - \frac{1}{4} + \cdots$ converges
by the Leibniz criterion (with $b_n = \frac{1}{n}$, which is decreasing
and tends to $0$). Its sum is $\ln 2$.

However, $\displaystyle\sum_{n=1}^{\infty}\frac{1}{n}$ diverges, so the
series is \textbf{conditionally convergent}, not absolutely convergent.
\end{example}

%===================================================================
\section{Power series}\label{sec:power-series}
%===================================================================
\index{power series}

\begin{definition}[Power series]\label{def:power-series}
A \textbf{power series} centered at $0$ is a series of the form
\[
   \sum_{n=0}^{\infty} a_n\, x^n,
\]
where $(a_n)_{n\geq 0}$ is a sequence of real numbers (the \textbf{coefficients})
and $x$ is a real variable.

More generally, a power series centered at $c \in \R$ is
$\sum_{n=0}^{\infty} a_n (x - c)^n$.
\end{definition}

\begin{theorem}[Radius of convergence]\label{thm:radius}
\index{radius of convergence}
For every power series $\sum a_n x^n$, there exists a unique
$R \in [0, +\infty]$ (called the \textbf{radius of convergence}) such that:
\begin{enumerate}
\item The series converges absolutely for $\abs{x} < R$.
\item The series diverges for $\abs{x} > R$.
\item At $\abs{x} = R$, anything can happen (convergence, divergence, or
      conditional convergence).
\end{enumerate}
The interval $(-R, R)$ is called the \textbf{interval of convergence}
\index{interval of convergence} (behavior at the endpoints $\pm R$ must be
checked separately).
\end{theorem}

\begin{center}
\begin{tikzpicture}[scale=1.1, >=stealth]
  \draw[->] (-5,0) -- (5,0) node[right] {$x$};
  \fill[green!15] (-3,-.35) rectangle (3,.35);
  \draw[very thick, green!60!black] (-3,0) -- (3,0);
  \draw[very thick, green!60!black, fill=white] (-3,0) circle (3pt);
  \draw[very thick, green!60!black, fill=white] (3,0) circle (3pt);
  \node[below=8pt, green!50!black, font=\small\bfseries] at (0,0)
        {absolute convergence};
  \draw[<->, thick, blue!60!black] (-3,-0.7) -- (3,-0.7);
  \node[below, blue!60!black, font=\small] at (0,-0.7) {$2R$};
  \node[above=4pt] at (-3,0) {$-R$};
  \node[above=4pt] at (3,0) {$R$};
  \node[above=4pt, red!60!black, font=\small] at (-4.2,0) {diverges};
  \node[above=4pt, red!60!black, font=\small] at (4.2,0) {diverges};
  \node[below=2pt, font=\small] at (-3,-0.1) {?};
  \node[below=2pt, font=\small] at (3,-0.1) {?};
  \node[below, font=\footnotesize] at (0,-1.2)
        {(endpoints require separate analysis)};
\end{tikzpicture}
\end{center}

\begin{theorem}[Hadamard formula]\label{thm:hadamard}
\index{Hadamard formula}
The radius of convergence of $\sum a_n x^n$ is
\[
   R = \frac{1}{\limsup_{n\to\infty}\sqrt[n]{\abs{a_n}}},
\]
with the conventions $\frac{1}{0} = +\infty$ and $\frac{1}{+\infty} = 0$.
\end{theorem}

\begin{remark}
In practice, one often computes $R$ via the ratio test:
$R = \displaystyle\lim_{n\to\infty}\abs{\frac{a_n}{a_{n+1}}}$
when this limit exists.
\end{remark}

\subsection{Classical power series}\label{subsec:classical-power}

\begin{example}[Exponential]\label{ex:exp-series}
\index{exponential!series}
\[
   e^x = \sum_{n=0}^{\infty}\frac{x^n}{n!}, \quad R = +\infty.
\]
Ratio: $\frac{a_{n+1}}{a_n}\cdot\abs{x} = \frac{\abs{x}}{n+1} \to 0 < 1$
for all $x \in \R$.
\end{example}

\begin{example}[Sine and cosine]\label{ex:sincos-series}
\index{sine!series}\index{cosine!series}
\[
   \sin x = \sum_{n=0}^{\infty}\frac{(-1)^n\, x^{2n+1}}{(2n+1)!},
   \qquad
   \cos x = \sum_{n=0}^{\infty}\frac{(-1)^n\, x^{2n}}{(2n)!},
   \quad R = +\infty.
\]
\end{example}

\begin{example}[Logarithm]\label{ex:log-series}
\index{logarithm!series}
\[
   \ln(1+x) = \sum_{n=1}^{\infty}\frac{(-1)^{n+1}\,x^n}{n}
   = x - \frac{x^2}{2} + \frac{x^3}{3} - \cdots, \quad R = 1.
\]
Converges for $x \in (-1, 1]$ (converges at $x=1$ by Leibniz; diverges at $x = -1$).
\end{example}

\begin{example}[Binomial series]\label{ex:binomial-series}
\index{binomial series}
For $\alpha \in \R$,
\[
   (1+x)^\alpha
   = \sum_{n=0}^{\infty}\binom{\alpha}{n}x^n
   \quad\text{where } \binom{\alpha}{n}
   = \frac{\alpha(\alpha-1)\cdots(\alpha-n+1)}{n!}, \quad R = 1.
\]
\end{example}

%===================================================================
\section{Cauchy product}\label{sec:cauchy-product}
%===================================================================
\index{Cauchy product}

\begin{theorem}[Cauchy product / Mertens' theorem]\label{thm:cauchy-product}
If $\sum a_n$ and $\sum b_n$ are absolutely convergent series with sums $A$
and $B$ respectively, then the \textbf{Cauchy product}
\[
   c_n = \sum_{k=0}^{n} a_k\, b_{n-k}
\]
defines an absolutely convergent series $\sum c_n$ with
$\displaystyle\sum_{n=0}^{\infty} c_n = A \cdot B$.
\end{theorem}

\begin{example}[Product $e^x \cdot e^y = e^{x+y}$]\label{ex:exp-product}
Let $A = \sum \frac{x^n}{n!}$ and $B = \sum \frac{y^n}{n!}$. The Cauchy
product has general term
\[
   c_n = \sum_{k=0}^{n}\frac{x^k}{k!}\cdot\frac{y^{n-k}}{(n-k)!}
       = \frac{1}{n!}\sum_{k=0}^{n}\binom{n}{k}x^k y^{n-k}
       = \frac{(x+y)^n}{n!}.
\]
Hence $e^x \cdot e^y = \sum \frac{(x+y)^n}{n!} = e^{x+y}$.
\end{example}

%===================================================================
\section{Visualizations}\label{sec:series-tikz}
%===================================================================

\begin{center}
\begin{tikzpicture}[scale=0.9, >=stealth]
  % -- Partial sums of geometric series --
  \draw[->] (-0.5,0) -- (8.5,0) node[right] {$n$};
  \draw[->] (0,-0.3) -- (0,3.5) node[above] {$S_n$};
  \draw[dashed, red!70!black] (-0.3,3) -- (8.3,3)
        node[right, font=\small] {$S = 2$};
  \node[below, font=\small\bfseries] at (4,-0.7)
        {Partial sums of $\sum \frac{1}{2^n}$};
  % S_0=1, S_1=1.5, S_2=1.75, S_3=1.875, ...
  \foreach \n/\y in {0/1.5, 1/2.25, 2/2.625, 3/2.8125,
                      4/2.90625, 5/2.953, 6/2.977, 7/2.988} {
    \fill[blue!70!black] (\n,{\y*3/3}) circle (2pt);
  }
  % rescale: S=2 maps to y=3, so multiply by 3/2
  % Redo: S_0=1->1.5, S_1=1.5->2.25, S_2=1.75->2.625...
  % Actually let me just plot correctly
\end{tikzpicture}
\end{center}

\begin{center}
\begin{tikzpicture}[scale=0.9, >=stealth]
  % -- Partial sums of alternating series --
  \draw[->] (-0.5,0) -- (10.5,0) node[right] {$n$};
  \draw[->] (0,-0.3) -- (0,3.5) node[above] {$S_n$};
  \draw[dashed, red!70!black] (-0.3,2.08) -- (10.3,2.08)
        node[right, font=\small] {$S = \ln 2$};
  \node[below, font=\small\bfseries] at (5,-0.7)
        {Partial sums of $\sum (-1)^{n+1}/n$};
  % S_1=1, S_2=0.5, S_3=0.833, S_4=0.583, S_5=0.783, S_6=0.617,
  % S_7=0.760, S_8=0.635, S_9=0.746, S_10=0.646
  % scale: multiply by 3 for plot
  \foreach \n/\y in {1/3.0, 2/1.5, 3/2.5, 4/1.75, 5/2.35,
                      6/1.85, 7/2.28, 8/1.905, 9/2.24, 10/1.94} {
    \fill[blue!70!black] (\n,\y) circle (2pt);
  }
  % connect with lines
  \draw[blue!40!black, thin]
    (1,3.0) -- (2,1.5) -- (3,2.5) -- (4,1.75) -- (5,2.35) --
    (6,1.85) -- (7,2.28) -- (8,1.905) -- (9,2.24) -- (10,1.94);
\end{tikzpicture}
\end{center}

%===================================================================
\section{Common errors}\label{sec:series-errors}
%===================================================================

\begin{tcolorbox}[
   colback=red!3!white,
   colframe=red!60!black,
   title={\textbf{Common Errors --- Series}},
   fonttitle=\bfseries,
   arc=2mm
]
\begin{enumerate}
\item \textbf{``$a_n \to 0$ so the series converges.''}
      NO! This is a necessary condition, not sufficient. The harmonic series
      is the classic counterexample.

\item \textbf{Confusing the series $\sum a_n$ with the sequence $(a_n)$.}
      A series is the \emph{sequence of partial sums}, not the sequence of
      terms. When we say ``the series converges,'' we mean $(S_n)$ converges.

\item \textbf{Applying the ratio/root test when $L = 1$.}
      When $L = 1$, the test gives \textbf{no conclusion}. You must use a
      different method.

\item \textbf{Forgetting to check endpoints of power series.}
      The radius of convergence tells you about the open interval $(-R,R)$.
      At $x = \pm R$, you must analyze convergence separately.

\item \textbf{Conditionally convergent $\neq$ divergent.}
      A conditionally convergent series \emph{does} converge; it is just not
      absolutely convergent. (Riemann's rearrangement theorem shows its terms
      can be rearranged to give any sum.)

\item \textbf{Comparison test requires non-negative terms (or absolute values).}
      Make sure $a_n \geq 0$ and $b_n \geq 0$ before applying the comparison test.
\end{enumerate}
\end{tcolorbox}

%===================================================================
\section{Exercises}\label{sec:series-exercises}
%===================================================================

\begin{exercise}[$\bigstar$]\label{exo:series1}
Determine whether $\displaystyle\sum_{n=1}^{\infty}\frac{n}{n+1}$ converges.
\end{exercise}

\begin{exercise}[$\bigstar$]\label{exo:series2}
Compute $\displaystyle\sum_{n=0}^{\infty}\frac{3}{4^n}$.
\end{exercise}

\begin{exercise}[$\bigstar$]\label{exo:series3}
Compute $\displaystyle\sum_{n=1}^{\infty}\frac{1}{n(n+1)(n+2)}$.
\textit{(Hint: partial fractions and telescoping.)}
\end{exercise}

\begin{exercise}[$\bigstar$]\label{exo:series4}
Using the ratio test, determine convergence of
$\displaystyle\sum_{n=0}^{\infty}\frac{n!}{3^n}$.
\end{exercise}

\begin{exercise}[$\bigstar$]\label{exo:series5}
Using the root test, determine convergence of
$\displaystyle\sum_{n=1}^{\infty}\left(\frac{n}{2n+1}\right)^n$.
\end{exercise}

\begin{exercise}[$\bigstar\bigstar$]\label{exo:series6}
Prove that $\displaystyle\sum_{n=1}^{\infty}\frac{1}{n^2}$ converges.
\textit{(Hint: compare with $\sum\frac{1}{n(n-1)}$ for $n \geq 2$.)}
\end{exercise}

\begin{exercise}[$\bigstar\bigstar$]\label{exo:series7}
Study the convergence of
$\displaystyle\sum_{n=2}^{\infty}\frac{1}{n(\ln n)^2}$.
\end{exercise}

\begin{exercise}[$\bigstar\bigstar$]\label{exo:series8}
Show that $\displaystyle\sum_{n=0}^{\infty}\frac{(-1)^n}{2n+1} = \frac{\pi}{4}$.
\textit{(Hint: relate to the power series of $\arctan x$ evaluated at $x=1$.)}
\end{exercise}

\begin{exercise}[$\bigstar\bigstar$]\label{exo:series9}
Find the radius of convergence and determine convergence at the endpoints for
$\displaystyle\sum_{n=1}^{\infty}\frac{x^n}{n\cdot 2^n}$.
\end{exercise}

\begin{exercise}[$\bigstar\bigstar$]\label{exo:series10}
Prove that if $\sum a_n$ converges absolutely and $(b_n)$ is bounded, then
$\sum a_n b_n$ converges absolutely.
\end{exercise}

\begin{exercise}[$\bigstar\bigstar\bigstar$]\label{exo:series11}
(\textit{Abel's summation.})
Let $(a_n)$ be a decreasing sequence with $a_n \to 0$, and let $(b_n)$ be a
sequence whose partial sums $B_n = \sum_{k=0}^{n}b_k$ are bounded. Prove that
$\sum a_n b_n$ converges.
\textit{(Hint: Abel summation by parts: $\sum_{k=p}^{q}a_k b_k
= a_q B_q - a_p B_{p-1} - \sum_{k=p}^{q-1}(a_{k+1}-a_k)B_k$.)}
\end{exercise}

\begin{exercise}[$\bigstar\bigstar\bigstar$]\label{exo:series12}
Find the radius of convergence of
$\displaystyle\sum_{n=0}^{\infty}\frac{(2n)!}{(n!)^2}\,x^n$.
\textit{(Hint: use the ratio test and Stirling's approximation, or compute
$a_{n+1}/a_n$ directly.)}
\end{exercise}

\begin{exercise}[$\bigstar\bigstar\bigstar$]\label{exo:series13}
Prove the \textbf{condensation test}\index{condensation test}: if $(a_n)$ is a
decreasing sequence with $a_n \geq 0$, then $\sum a_n$ converges if and only if
$\sum 2^n\, a_{2^n}$ converges. Use this to give another proof that
$\sum 1/n^\alpha$ converges iff $\alpha > 1$.
\end{exercise}

%===================================================================
\section*{Chapter summary}
%===================================================================

\begin{tcolorbox}[
   colback=green!3!white,
   colframe=green!50!black,
   title={\textbf{Chapter 4 --- Summary}},
   fonttitle=\bfseries,
   arc=2mm
]
\begin{itemize}
\item A series $\sum a_n$ converges iff its partial sums $(S_n)$ converge.
\item Necessary condition: $a_n \to 0$ (not sufficient---harmonic series!).
\item Geometric series $\sum r^n$: converges iff $\abs{r}<1$, sum $= \frac{1}{1-r}$.
\item Telescoping series: $\sum(b_n - b_{n+1}) = b_0 - \lim b_n$.
\item For series with non-negative terms:
   \begin{itemize}
   \item Comparison: $0 \leq a_n \leq b_n$ and $\sum b_n$ conv.\ $\implies$ $\sum a_n$ conv.
   \item Ratio test: $\frac{a_{n+1}}{a_n} \to L$; $L<1 \implies$ conv., $L>1 \implies$ div.
   \item Root test: $\sqrt[n]{a_n} \to L$; $L<1 \implies$ conv., $L>1 \implies$ div.
   \item Integral test: compare with $\int f$.
   \end{itemize}
\item Absolutely convergent $\implies$ convergent (converse false).
\item Leibniz: alternating series with $b_n \searrow 0$ converges; error $\leq b_{n+1}$.
\item Power series $\sum a_n x^n$: radius $R$ (Hadamard), converges absolutely on $(-R,R)$.
\item Cauchy product: if $\sum a_n$, $\sum b_n$ abs.\ conv., then
      $(\sum a_n)(\sum b_n) = \sum c_n$ with $c_n = \sum_{k=0}^{n}a_k b_{n-k}$.
\end{itemize}
\end{tcolorbox}
%% ====================================================================
%%  PART 3 — Chapters 5–6, Bibliography, Index
%%  Real Analysis I  (English)
%% ====================================================================

%% ====================================================================
%%  CHAPTER 5 — LIMITS AND CONTINUITY OF FUNCTIONS
%% ====================================================================
\chapter{Limits and Continuity of Functions}\label{ch:limits-continuity}

\minitoc\bigskip

\lettrine[lines=3,lhang=0.15]{T}{his chapter} develops the theory of
limits and continuity for real-valued functions of a real variable.
These ideas sit at the very heart of analysis: differentiation,
integration, and the study of series all depend on a precise
understanding of what it means for a function to ``approach a value''
or to ``have no jumps.''  We give every $\varepsilon$--$\delta$
argument in complete detail, including the scratch work that motivates
the choice of~$\delta$.

%% ────────────────────────────────────────────────────────────────────
\section{Limits of Functions}
%% ────────────────────────────────────────────────────────────────────

\subsection{The \texorpdfstring{$\varepsilon$--$\delta$}{epsilon--delta}
  Definition}

Let us fix the setting.  We consider a function
$f\colon D\to\mathbb{R}$, where $D\subset\mathbb{R}$, and a point
$a\in\mathbb{R}$ that is a \emph{limit point} (accumulation point)
of~$D$: every open interval containing~$a$ meets~$D\setminus\{a\}$.

\begin{definition}[Limit of a function]\label{def:limit-function}
\index{limit!of a function}
Let $f\colon D\to\mathbb{R}$ and let $a$ be a limit point of~$D$.
We say that
\[
  \lim_{x\to a} f(x)=L
\]
if and only if
\[
  \forall\,\varepsilon>0,\;\exists\,\delta>0,\;
  \forall\, x\in D,\quad
  0<|x-a|<\delta \;\Longrightarrow\; |f(x)-L|<\varepsilon.
\]
\end{definition}

\begin{tcolorbox}[colback=blue!3,colframe=blue!50!black,
  title={Reading the definition}]
\textbf{In words:} no matter how small the tolerance
$\varepsilon>0$ is, we can find a ``response'' $\delta>0$ so that
every $x$ in the domain within distance $\delta$ of~$a$ (but
$x\neq a$) has its image $f(x)$ within distance $\varepsilon$
of~$L$.

\medskip\noindent
\textbf{Key points:}
\begin{enumerate}[nosep]
\item The value $f(a)$, if it exists, plays \emph{no role}; we require
  $0<|x-a|$.
\item $\delta$ is allowed to depend on~$\varepsilon$ (and on $a$
  and~$f$).
\item The quantifier order matters:
  $\forall\,\varepsilon\;\exists\,\delta\;\forall\, x$.
\end{enumerate}
\end{tcolorbox}

\begin{center}
\begin{tikzpicture}[scale=1.1,>=Stealth]
  % axes
  \draw[->] (-0.5,0) -- (6,0) node[right] {$x$};
  \draw[->] (0,-0.5) -- (0,4.5) node[above] {$y$};
  % function
  \draw[thick,blue,domain=0.3:5.5,samples=100]
    plot(\x,{0.4*\x*\x - 1.6*\x + 3.2});
  % point a, L
  \def\aval{3}
  \pgfmathsetmacro{\Lval}{0.4*\aval*\aval - 1.6*\aval + 3.2}
  % epsilon band
  \def\eps{0.6}
  \fill[red!10] (-0.3,{\Lval-\eps}) rectangle (5.7,{\Lval+\eps});
  \draw[red,dashed] (-0.3,{\Lval-\eps}) -- (5.7,{\Lval-\eps})
    node[right,font=\footnotesize] {$L-\varepsilon$};
  \draw[red,dashed] (-0.3,{\Lval+\eps}) -- (5.7,{\Lval+\eps})
    node[right,font=\footnotesize] {$L+\varepsilon$};
  % delta band
  \def\del{0.7}
  \fill[green!10] ({\aval-\del},0) rectangle ({\aval+\del},4.3);
  \draw[green!50!black,dashed] ({\aval-\del},0) -- ({\aval-\del},4.3);
  \draw[green!50!black,dashed] ({\aval+\del},0) -- ({\aval+\del},4.3);
  % labels
  \draw[<->] ({\aval-\del},-0.35) -- node[below,font=\footnotesize]
    {$\delta$} (\aval,-0.35);
  \draw[<->] ({\aval},-0.35) -- node[below,font=\footnotesize]
    {$\delta$} ({\aval+\del},-0.35);
  \draw[<->] (-0.25,\Lval) -- node[left,font=\footnotesize]
    {$\varepsilon$} (-0.25,{\Lval+\eps});
  \draw[<->] (-0.25,{\Lval-\eps}) -- node[left,font=\footnotesize]
    {$\varepsilon$} (-0.25,\Lval);
  % point
  \filldraw (\aval,\Lval) circle (1.5pt);
  \draw[dotted] (\aval,0) node[below,font=\footnotesize] {$a$}
    -- (\aval,\Lval) -- (0,\Lval) node[left,font=\footnotesize] {$L$};
  % caption
  \node[font=\footnotesize,text width=3.5cm,align=center]
    at (5,3.8) {All of the curve inside the green strip lies inside the
    red strip.};
\end{tikzpicture}
\captionof{figure}{The $\varepsilon$--$\delta$ definition of
  $\displaystyle\lim_{x\to a}f(x)=L$.}\label{fig:eps-delta-limit}
\end{center}

\subsection{Sequential Characterization (Heine's Criterion)}

The following theorem connects the limit of a function with sequential
limits, which are often easier to work with.

\begin{theorem}[Heine's criterion]\label{thm:heine}
\index{Heine's criterion}\index{sequential characterization!of function limit}
Let $f\colon D\to\mathbb{R}$ and let $a$ be a limit point of~$D$.
Then
\[
  \lim_{x\to a}f(x)=L
\]
if and only if for every sequence $(x_n)$ in $D\setminus\{a\}$
with $x_n\to a$ we have $f(x_n)\to L$.
\end{theorem}

\begin{proof}
We prove both directions in full detail.

\medskip\noindent
$(\Rightarrow)$\; Assume $\lim_{x\to a}f(x)=L$.
Let $(x_n)$ be any sequence in $D\setminus\{a\}$ with $x_n\to a$.
We must show $f(x_n)\to L$.

Fix $\varepsilon>0$.  By the limit hypothesis, there exists $\delta>0$
such that
\[
  \forall\, x\in D,\quad
  0<|x-a|<\delta \implies |f(x)-L|<\varepsilon.
\]
Since $x_n\to a$, there exists $N\in\mathbb{N}$ such that for all
$n\ge N$ we have $|x_n - a|<\delta$.  Because $x_n\neq a$ for every~$n$
(they belong to $D\setminus\{a\}$), we have $0<|x_n - a|<\delta$ for
$n\ge N$, and therefore $|f(x_n)-L|<\varepsilon$ for $n\ge N$.  This
proves $f(x_n)\to L$.

\medskip\noindent
$(\Leftarrow)$\; We prove the contrapositive.  Assume that
$\lim_{x\to a}f(x)\neq L$; we construct a sequence $(x_n)$ in
$D\setminus\{a\}$ with $x_n\to a$ but $f(x_n)\not\to L$.

The negation of the $\varepsilon$--$\delta$ definition gives:
\[
  \exists\,\varepsilon_0>0,\;\forall\,\delta>0,\;
  \exists\, x\in D,\quad
  0<|x-a|<\delta \;\text{and}\; |f(x)-L|\ge\varepsilon_0.
\]
Apply this with $\delta=\tfrac{1}{n}$ for each $n\ge 1$: we obtain
$x_n\in D$ with $0<|x_n-a|<\tfrac{1}{n}$ and
$|f(x_n)-L|\ge\varepsilon_0$.  Then $x_n\in D\setminus\{a\}$,
$x_n\to a$ (by the squeeze theorem, since $|x_n-a|<1/n\to 0$), yet
$f(x_n)\not\to L$ because $|f(x_n)-L|\ge\varepsilon_0$ for every~$n$.
\end{proof}

\begin{remark}
Heine's criterion is especially useful for \emph{disproving} the
existence of a limit: find two sequences converging to~$a$ whose images
converge to different values.
\end{remark}

%% ────────────────────────────────────────────────────────────────────
\subsection{Worked Example: \texorpdfstring{$\lim_{x\to 2}(3x-1)=5$}{lim(3x-1)=5}}

\begin{example}\label{ex:limit-linear}
\textbf{Claim.}\; $\displaystyle\lim_{x\to 2}(3x-1)=5$.

\medskip\noindent
\textbf{Scratch work (not part of the proof).}\;
We need $|f(x)-L|<\varepsilon$, i.e.\
$|(3x-1)-5|<\varepsilon$, i.e.\ $|3x-6|<\varepsilon$, i.e.\
$3|x-2|<\varepsilon$, i.e.\ $|x-2|<\varepsilon/3$.
So $\delta=\varepsilon/3$ should work.

\medskip\noindent
\textbf{Formal proof.}\;
Let $\varepsilon>0$.  Set $\delta=\varepsilon/3>0$.  Then for every
$x\in\mathbb{R}$ with $0<|x-2|<\delta$,
\[
  |(3x-1)-5| = |3x-6| = 3|x-2| < 3\delta = 3\cdot\frac{\varepsilon}{3}
  = \varepsilon.
\]
By definition, $\lim_{x\to 2}(3x-1)=5$.  \qed
\end{example}

\subsection{Worked Example: \texorpdfstring{$\lim_{x\to 0}\frac{\sin x}{x}=1$}{lim sin(x)/x = 1}}

\begin{example}\label{ex:sinx-over-x}
\textbf{Claim.}\; $\displaystyle\lim_{x\to 0}\frac{\sin x}{x}=1$.

\medskip\noindent
\textbf{Geometric argument.}\;
For $0<x<\pi/2$, consider the unit circle.  The area of the inscribed
triangle, the circular sector, and the circumscribed triangle give the
chain of inequalities
\[
  \frac{\sin x}{2}
  \;\le\; \frac{x}{2}
  \;\le\; \frac{\tan x}{2}.
\]
Dividing by $\frac{\sin x}{2}>0$:
\[
  1 \;\le\; \frac{x}{\sin x} \;\le\; \frac{1}{\cos x}.
\]
Taking reciprocals (all terms positive):
\[
  \cos x \;\le\; \frac{\sin x}{x} \;\le\; 1.
\]
Since $\cos x\to 1$ as $x\to 0^+$, the squeeze theorem gives
$\frac{\sin x}{x}\to 1$ as $x\to 0^+$.

For $x<0$, set $u=-x>0$; then $\frac{\sin x}{x}
=\frac{\sin(-u)}{-u}=\frac{\sin u}{u}\to 1$.

Therefore $\displaystyle\lim_{x\to 0}\frac{\sin x}{x}=1$.  \qed
\end{example}

%% ────────────────────────────────────────────────────────────────────
\section{One-Sided Limits, Infinite Limits, Limits at Infinity}

\begin{definition}[One-sided limits]\label{def:one-sided}
\index{limit!one-sided}
\leavevmode
\begin{itemize}
\item \textbf{Right-hand limit:}
$\displaystyle\lim_{x\to a^+}f(x)=L$ means
\[
  \forall\,\varepsilon>0,\;\exists\,\delta>0,\;
  \forall\, x\in D,\quad
  a<x<a+\delta \implies |f(x)-L|<\varepsilon.
\]
\item \textbf{Left-hand limit:}
$\displaystyle\lim_{x\to a^-}f(x)=L$ means
\[
  \forall\,\varepsilon>0,\;\exists\,\delta>0,\;
  \forall\, x\in D,\quad
  a-\delta<x<a \implies |f(x)-L|<\varepsilon.
\]
\end{itemize}
\end{definition}

\begin{proposition}\label{prop:two-sided-from-one-sided}
$\displaystyle\lim_{x\to a}f(x)=L$ if and only if both one-sided
limits exist and equal~$L$.
\end{proposition}

\begin{proof}
The forward direction is immediate (restrict the quantifier).
For the converse, given $\varepsilon>0$, let $\delta_1$ work for the
right-hand limit and $\delta_2$ for the left-hand limit; take
$\delta=\min(\delta_1,\delta_2)$.
\end{proof}

\begin{definition}[Infinite limits and limits at infinity]
\label{def:infinite-limits}\index{limit!infinite}\index{limit!at infinity}
\leavevmode
\begin{itemize}
\item $\displaystyle\lim_{x\to a}f(x)=+\infty$ means
\[
  \forall\, M>0,\;\exists\,\delta>0,\;
  \forall\, x\in D,\quad
  0<|x-a|<\delta \implies f(x)>M.
\]
\item $\displaystyle\lim_{x\to +\infty}f(x)=L$ means
\[
  \forall\,\varepsilon>0,\;\exists\, A>0,\;
  \forall\, x\in D,\quad
  x>A \implies |f(x)-L|<\varepsilon.
\]
\item $\displaystyle\lim_{x\to +\infty}f(x)=+\infty$ means
\[
  \forall\, M>0,\;\exists\, A>0,\;
  \forall\, x\in D,\quad
  x>A \implies f(x)>M.
\]
\end{itemize}
Analogous definitions hold for $-\infty$.
\end{definition}

\begin{example}
$\displaystyle\lim_{x\to 0^+}\frac{1}{x}=+\infty$.

\medskip\noindent
\textbf{Proof.}\; Let $M>0$.  Set $\delta=1/M$.  If $0<x<\delta$,
then $1/x>1/\delta=M$.  \qed
\end{example}

%% ────────────────────────────────────────────────────────────────────
\section{Continuity at a Point}

\begin{definition}[Continuity at a point]\label{def:continuity-point}
\index{continuity!at a point}
Let $f\colon D\to\mathbb{R}$ and $a\in D$.  We say $f$ is
\textbf{continuous at~$a$} if
\[
  \forall\,\varepsilon>0,\;\exists\,\delta>0,\;
  \forall\, x\in D,\quad
  |x-a|<\delta \implies |f(x)-f(a)|<\varepsilon.
\]
Equivalently, if $a$ is a limit point of $D$: $\lim_{x\to a}f(x)=f(a)$.
If $a$ is an isolated point of $D$, then $f$ is automatically continuous
at~$a$.
\end{definition}

\begin{tcolorbox}[colback=yellow!5,colframe=orange!60!black,
  title={Limit vs.\ continuity}]
Comparing Definition~\ref{def:limit-function} and
Definition~\ref{def:continuity-point}:
\begin{itemize}[nosep]
\item For the \emph{limit}, we require $0<|x-a|$
  (we exclude~$a$ itself).
\item For \emph{continuity}, we do \emph{not} exclude~$a$,
  and the target is $f(a)$.
\item Continuity at~$a$ \emph{requires} $a\in D$ (so $f(a)$ exists).
\end{itemize}
\end{tcolorbox}

\begin{proposition}[Sequential characterization of continuity]
\label{prop:seq-continuity}\index{continuity!sequential characterization}
Let $f\colon D\to\mathbb{R}$ and $a\in D$.  Then $f$ is continuous
at~$a$ if and only if for every sequence $(x_n)$ in $D$ with $x_n\to a$,
we have $f(x_n)\to f(a)$.
\end{proposition}

\begin{proof}
The proof is essentially the same as for Theorem~\ref{thm:heine}, but
now we allow $x_n=a$ and compare with $f(a)$ instead of~$L$.  We omit
the repetition and invite the reader to write the details as an exercise.
\end{proof}

\subsection{Examples of Continuous Functions}

\begin{example}[Polynomials are continuous]\label{ex:poly-cont}
\index{continuity!of polynomials}
Every polynomial $p(x)=a_0+a_1 x+\cdots+a_n x^n$ is continuous on
$\mathbb{R}$.

\medskip\noindent
\textbf{Proof sketch.}\;
The constant function $x\mapsto c$ is continuous (take any $\delta$),
and $x\mapsto x$ is continuous (take $\delta=\varepsilon$).
Continuity is preserved by sums and products (limit laws), so every
polynomial is continuous.
\end{example}

\begin{example}[$|x|$ is continuous]\label{ex:abs-cont}
\index{continuity!of absolute value}
The function $f(x)=|x|$ is continuous on $\mathbb{R}$.

\medskip\noindent
\textbf{$\varepsilon$--$\delta$ proof.}\;
Fix $a\in\mathbb{R}$ and $\varepsilon>0$.  Set $\delta=\varepsilon$.
For $|x-a|<\delta$:
\[
  \bigl||x|-|a|\bigr| \;\le\; |x-a| < \delta = \varepsilon,
\]
where the first inequality is the \emph{reverse triangle inequality}.
Hence $f$ is continuous at~$a$.  \qed
\end{example}

\begin{example}[Dirichlet function --- nowhere continuous]
\label{ex:dirichlet}\index{Dirichlet function}
Define $\mathbf{1}_{\mathbb{Q}}\colon\mathbb{R}\to\mathbb{R}$ by
\[
  \mathbf{1}_{\mathbb{Q}}(x) =
  \begin{cases} 1 & \text{if } x\in\mathbb{Q},\\
                 0 & \text{if } x\notin\mathbb{Q}.
  \end{cases}
\]
This function is \textbf{nowhere continuous}.

\medskip\noindent
\textbf{Proof.}\;
Fix any $a\in\mathbb{R}$.  Take $\varepsilon_0=\tfrac{1}{2}$.
By the density of $\mathbb{Q}$ and $\mathbb{R}\setminus\mathbb{Q}$,
every interval $(a-\delta,a+\delta)$ contains both a rational~$r$ and an
irrational~$s$.  Then $|\mathbf{1}_{\mathbb{Q}}(r)-\mathbf{1}_{\mathbb{Q}}(s)|=1\ge\varepsilon_0$.
So either $|\mathbf{1}_{\mathbb{Q}}(r)-\mathbf{1}_{\mathbb{Q}}(a)|$
or $|\mathbf{1}_{\mathbb{Q}}(s)-\mathbf{1}_{\mathbb{Q}}(a)|$ is
$\ge\tfrac12=\varepsilon_0$.  Hence $f$ is not continuous at~$a$.\qed

\medskip
Alternatively, using the sequential characterization: if $a$ is
rational, pick an irrational sequence $s_n\to a$; then
$f(s_n)=0\not\to 1=f(a)$.  Similarly if $a$ is irrational.
\end{example}

%% ────────────────────────────────────────────────────────────────────
\section{Algebra of Continuous Functions}

\begin{proposition}\label{prop:algebra-cont}
If $f,g$ are continuous at $a$, then so are $f+g$, $f\cdot g$,
$\lambda f$ ($\lambda\in\mathbb{R}$), and $f/g$ (provided $g(a)\neq 0$).
The composition $g\circ f$ is continuous at $a$ whenever $f$ is
continuous at $a$ and $g$ is continuous at $f(a)$.
\end{proposition}

\begin{proof}
Each statement reduces to the corresponding limit law via
$\lim_{x\to a}f(x)=f(a)$.  For the composition: let $\varepsilon>0$.
By continuity of $g$ at $f(a)$, find $\eta>0$ with
$|y-f(a)|<\eta\Rightarrow|g(y)-g(f(a))|<\varepsilon$.  By continuity
of $f$ at $a$, find $\delta>0$ with
$|x-a|<\delta\Rightarrow|f(x)-f(a)|<\eta$.  Then
$|x-a|<\delta\Rightarrow|g(f(x))-g(f(a))|<\varepsilon$.
\end{proof}

%% ────────────────────────────────────────────────────────────────────
\section{The Intermediate Value Theorem}

\begin{theorem}[Intermediate Value Theorem (IVT)]
\label{thm:ivt}\index{Intermediate Value Theorem}
Let $f\colon[a,b]\to\mathbb{R}$ be continuous with $f(a)<f(b)$
(the case $f(a)>f(b)$ is analogous).  Then for every value
$\gamma$ with $f(a)<\gamma<f(b)$, there exists $c\in(a,b)$ with
$f(c)=\gamma$.
\end{theorem}

\begin{proof}
Define
\[
  S = \{x\in[a,b] : f(x)<\gamma\}.
\]
$S$ is nonempty ($a\in S$) and bounded above (by~$b$), so
$c:=\sup S$ exists.  We show $f(c)=\gamma$.

\medskip\noindent
\textbf{Step 1: $f(c)\le\gamma$.}\;
Since $c=\sup S$, for each $n\ge 1$ there exists $x_n\in S$ with
$c-\tfrac{1}{n}<x_n\le c$.  Then $x_n\to c$ and
$f(x_n)<\gamma$ for every~$n$.  By continuity at~$c$,
$f(c)=\lim f(x_n)\le\gamma$.

\medskip\noindent
\textbf{Step 2: $f(c)\ge\gamma$.}\;
If $c=b$, then $f(c)=f(b)>\gamma$ and we are done.  If $c<b$, then
for every $n$ large enough that $c+\tfrac{1}{n}\le b$, the point
$c+\tfrac{1}{n}\notin S$ (it exceeds the supremum), so
$f(c+\tfrac{1}{n})\ge\gamma$.  Letting $n\to\infty$, by continuity,
$f(c)\ge\gamma$.

\medskip\noindent
Steps 1 and 2 together give $f(c)=\gamma$.  Moreover $c\neq a$
(because $f(a)<\gamma$ and continuity would force nearby values
$<\gamma$, giving $c>a$) and similarly $c\neq b$, so $c\in(a,b)$.
\end{proof}

\begin{center}
\begin{tikzpicture}[scale=0.9,>=Stealth]
  \draw[->] (-0.3,0) -- (6.5,0) node[right] {$x$};
  \draw[->] (0,-0.5) -- (0,4.5) node[above] {$y$};
  \draw[thick,blue,domain=0.5:6,samples=80]
    plot(\x,{0.3*(\x-1)*(\x-3.5)*(\x-5.5)+2.5});
  \draw[dashed,red] (-0.2,2) -- (6.3,2)
    node[right,font=\footnotesize]{$\gamma$};
  \draw (0.5,0.02) -- (0.5,-0.1) node[below,font=\footnotesize]{$a$};
  \draw (6,0.02) -- (6,-0.1) node[below,font=\footnotesize]{$b$};
  \pgfmathsetmacro{\fa}{0.3*(0.5-1)*(0.5-3.5)*(0.5-5.5)+2.5}
  \draw[dotted] (0.5,0) -- (0.5,\fa) -- (0,\fa)
    node[left,font=\footnotesize]{$f(a)$};
  \pgfmathsetmacro{\fb}{0.3*(6-1)*(6-3.5)*(6-5.5)+2.5}
  \draw[dotted] (6,0) -- (6,\fb) -- (0,\fb)
    node[left,font=\footnotesize]{$f(b)$};
  % mark the c's
  \node[font=\footnotesize,below] at (1.55,-0.1) {$c$};
  \draw[dotted] (1.55,0) -- (1.55,2);
  \filldraw[red] (1.55,2) circle (2pt);
\end{tikzpicture}
\captionof{figure}{The Intermediate Value Theorem: the graph must
  cross the horizontal line $y=\gamma$.}\label{fig:ivt}
\end{center}

\begin{corollary}[Existence of $n$th roots]
\label{cor:nth-roots}\index{nth root!existence}
For every $y>0$ and every $n\in\mathbb{N}_{\ge 1}$, there exists a
unique $x>0$ with $x^n=y$.
\end{corollary}

\begin{proof}
The function $f(x)=x^n$ is continuous on $[0,\infty)$ with $f(0)=0$
and $f(x)\to+\infty$.  Choosing $b$ large enough so that $b^n>y$,
the IVT on $[0,b]$ gives $c$ with $c^n=y$.  Uniqueness follows
because $f$ is strictly increasing on $[0,\infty)$.
\end{proof}

\begin{corollary}[Odd-degree polynomials have a real root]
\label{cor:odd-poly}\index{polynomial!odd-degree root}
Every polynomial of odd degree with real coefficients has at least
one real root.
\end{corollary}

\begin{proof}
If $p(x)=x^{2k+1}+\text{lower-order terms}$, then $p(x)\to+\infty$
as $x\to+\infty$ and $p(x)\to-\infty$ as $x\to-\infty$.  So there
exist $a<b$ with $p(a)<0<p(b)$.  Apply the IVT.
\end{proof}

%% ────────────────────────────────────────────────────────────────────
\section{The Extreme Value Theorem}

\begin{theorem}[Extreme Value Theorem]
\label{thm:evt}\index{Extreme Value Theorem}
Let $f\colon[a,b]\to\mathbb{R}$ be continuous.  Then $f$ is bounded
and attains its maximum and minimum: there exist $x_m,x_M\in[a,b]$
with
\[
  f(x_m) \le f(x) \le f(x_M)
  \quad\text{for all } x\in[a,b].
\]
\end{theorem}

\begin{proof}
We prove that $f$ attains its supremum (the infimum is analogous).

\medskip\noindent
\textbf{Step 1: $f$ is bounded above.}\;
Suppose not.  Then for each $n\ge 1$, there exists $x_n\in[a,b]$ with
$f(x_n)>n$.  By the Bolzano--Weierstrass theorem, $(x_n)$ has a
convergent subsequence $x_{n_k}\to c\in[a,b]$.  By continuity,
$f(x_{n_k})\to f(c)$, contradicting $f(x_{n_k})>n_k\to\infty$.

\medskip\noindent
\textbf{Step 2: The supremum is attained.}\;
Let $M=\sup_{x\in[a,b]}f(x)<\infty$ (Step~1).  By definition of
supremum, for each $n\ge 1$ there exists $x_n\in[a,b]$ with
$M-\tfrac{1}{n}<f(x_n)\le M$.  Again by Bolzano--Weierstrass,
extract a convergent subsequence $x_{n_k}\to x_M\in[a,b]$.  By
continuity, $f(x_{n_k})\to f(x_M)$.  Since
$M-\tfrac{1}{n_k}<f(x_{n_k})\le M$ and both sides converge to~$M$,
the squeeze theorem gives $f(x_M)=M$.
\end{proof}

%% ────────────────────────────────────────────────────────────────────
\section{Uniform Continuity}

\begin{definition}[Uniform continuity]\label{def:unif-cont}
\index{uniform continuity}
A function $f\colon D\to\mathbb{R}$ is \textbf{uniformly continuous}
on~$D$ if
\[
  \forall\,\varepsilon>0,\;\exists\,\delta>0,\;
  \forall\, x,y\in D,\quad
  |x-y|<\delta \implies |f(x)-f(y)|<\varepsilon.
\]
\end{definition}

\begin{tcolorbox}[colback=red!3,colframe=red!60!black,
  title={Uniform vs.\ pointwise continuity}]
\begin{itemize}[nosep]
\item \textbf{Pointwise:}
  $\forall\, a\in D,\;\forall\,\varepsilon>0,\;\exists\,\delta>0,\;
   \forall\, x\in D,\;\ldots$\;
  ($\delta$ may depend on~$a$).
\item \textbf{Uniform:}
  $\forall\,\varepsilon>0,\;\exists\,\delta>0,\;
   \forall\, x,y\in D,\;\ldots$\;
  (\emph{one} $\delta$ works for \emph{all} pairs).
\end{itemize}
Uniform $\Rightarrow$ pointwise (set $y=a$).  The converse is false.
\end{tcolorbox}

\begin{center}
\begin{tikzpicture}[scale=0.85,>=Stealth]
  % Pointwise
  \begin{scope}
    \draw[->] (-0.3,0) -- (5,0) node[right]{$x$};
    \draw[->] (0,-0.3) -- (0,4) node[above]{$y$};
    \draw[thick,blue,domain=0.15:4.5,samples=80]
      plot(\x,{0.2*\x*\x});
    % different deltas at different points
    \foreach \a/\d in {1/1.2, 3/0.4}{
      \draw[green!50!black,dashed] ({\a-\d},0) -- ({\a-\d},3.8);
      \draw[green!50!black,dashed] ({\a+\d},0) -- ({\a+\d},3.8);
      \draw[<->,green!50!black] ({\a-\d},-0.25) --
        node[below,font=\tiny]{$\delta$} ({\a+\d},-0.25);
    }
    \node[font=\footnotesize] at (2.3,4.2)
      {Pointwise: $\delta$ shrinks as $a$ grows};
  \end{scope}
  % Uniform
  \begin{scope}[xshift=7cm]
    \draw[->] (-0.3,0) -- (5,0) node[right]{$x$};
    \draw[->] (0,-0.3) -- (0,4) node[above]{$y$};
    \draw[thick,blue,domain=0:4.5,samples=80]
      plot(\x,{0.7*\x});
    % same delta everywhere
    \foreach \a in {1,2.5,4}{
      \draw[green!50!black,dashed] ({\a-0.5},0) -- ({\a-0.5},3.8);
      \draw[green!50!black,dashed] ({\a+0.5},0) -- ({\a+0.5},3.8);
    }
    \node[font=\footnotesize] at (2.3,4.2)
      {Uniform: same $\delta$ works everywhere};
  \end{scope}
\end{tikzpicture}
\captionof{figure}{Pointwise continuity (left) vs.\ uniform
  continuity (right).}\label{fig:uniform-vs-pointwise}
\end{center}

\begin{example}[$x^2$ is not uniformly continuous on $\mathbb{R}$]
\label{ex:x2-not-uc}\index{uniform continuity!counterexample}
Take $\varepsilon_0=1$.  For any $\delta>0$, choose $x=1/\delta$ and
$y=x+\delta/2$.  Then $|x-y|=\delta/2<\delta$ but
\[
  |x^2-y^2| = |x-y|\cdot|x+y| = \frac{\delta}{2}\Bigl(\frac{2}{\delta}
  +\frac{\delta}{2}\Bigr) = 1+\frac{\delta^2}{4} > 1 = \varepsilon_0.
\]
So $f(x)=x^2$ is not uniformly continuous on~$\mathbb{R}$.
\end{example}

\begin{theorem}[Heine's theorem]\label{thm:heine-uc}
\index{Heine's theorem (uniform continuity)}
If $f\colon[a,b]\to\mathbb{R}$ is continuous, then $f$ is
uniformly continuous on $[a,b]$.
\end{theorem}

\begin{proof}
Suppose for contradiction that $f$ is \emph{not} uniformly continuous.
Then there exists $\varepsilon_0>0$ such that for every $\delta>0$ there
exist $x,y\in[a,b]$ with $|x-y|<\delta$ but
$|f(x)-f(y)|\ge\varepsilon_0$.

Applying this with $\delta=1/n$, we get sequences $(x_n),(y_n)$ in
$[a,b]$ with $|x_n-y_n|<1/n$ and $|f(x_n)-f(y_n)|\ge\varepsilon_0$.

By Bolzano--Weierstrass, $(x_n)$ has a subsequence
$x_{n_k}\to c\in[a,b]$.  Since $|y_{n_k}-x_{n_k}|<1/n_k\to 0$,
we also have $y_{n_k}\to c$.  By continuity at~$c$:
\[
  f(x_{n_k})\to f(c) \quad\text{and}\quad f(y_{n_k})\to f(c).
\]
Hence $|f(x_{n_k})-f(y_{n_k})|\to 0$, contradicting
$|f(x_{n_k})-f(y_{n_k})|\ge\varepsilon_0>0$.
\end{proof}

\begin{definition}[Lipschitz continuity]\label{def:lipschitz}
\index{Lipschitz continuity}
$f\colon D\to\mathbb{R}$ is \textbf{Lipschitz} (with constant~$K$) if
\[
  \forall\, x,y\in D,\quad |f(x)-f(y)|\le K|x-y|.
\]
\end{definition}

\begin{proposition}\label{prop:lip-implies-uc}
Lipschitz $\Longrightarrow$ uniformly continuous.
\end{proposition}

\begin{proof}
Given $\varepsilon>0$, take $\delta=\varepsilon/K$.
Then $|x-y|<\delta$ implies $|f(x)-f(y)|\le K|x-y|<K\delta=\varepsilon$.
\end{proof}

\begin{example}
$f(x)=\sqrt{x}$ is uniformly continuous on $[0,\infty)$ but \emph{not}
Lipschitz on $[0,\infty)$ (the derivative $1/(2\sqrt{x})\to\infty$ as
$x\to 0^+$).  This shows the converse of
Proposition~\ref{prop:lip-implies-uc} is false.
\end{example}

%% ────────────────────────────────────────────────────────────────────
\section{Common Errors}

\begin{tcolorbox}[colback=red!4,colframe=red!70!black,
  title={Frequent mistakes in Chapter~5}]
\begin{enumerate}[nosep]
\item \textbf{Forgetting $0<|x-a|$ in the limit definition.}\;
  The limit is about behaviour \emph{near} $a$, not \emph{at} $a$.
\item \textbf{Choosing $\delta$ that depends on $x$.}\;
  In $\varepsilon$--$\delta$ proofs, $\delta$ is chosen \emph{before}
  $x$ is specified.
\item \textbf{Confusing pointwise and uniform continuity.}\;
  Check the quantifier order!
\item \textbf{Applying the IVT to a function that is not continuous.}\;
  The step function $\lfloor x\rfloor$ has range $\mathbb{Z}$ despite
  $f(0)=0,\;f(3/2)=1$; it skips non-integer values because it is
  discontinuous.
\item \textbf{Applying the Extreme Value Theorem on an open interval.}\;
  $f(x)=1/x$ on $(0,1)$ is continuous but unbounded.
\item \textbf{Writing ``$f$ is continuous so $f$ is uniformly
  continuous'' without specifying a closed bounded interval.}\;
  Heine's theorem requires $[a,b]$.
\end{enumerate}
\end{tcolorbox}

%% ────────────────────────────────────────────────────────────────────
\section{Exercises}

\begin{exercise}\label{exo:ch5-1}
Prove from the $\varepsilon$--$\delta$ definition that
$\displaystyle\lim_{x\to 3}(x^2)=9$.

\emph{Hint:} factor $|x^2-9|=|x-3|\cdot|x+3|$ and bound $|x+3|$
by restricting $\delta\le 1$.
\end{exercise}

\begin{exercise}\label{exo:ch5-2}
Prove that $\displaystyle\lim_{x\to 1}\frac{1}{x}=1$ using
$\varepsilon$--$\delta$.
\end{exercise}

\begin{exercise}\label{exo:ch5-3}
Use Heine's criterion to show that
$\displaystyle\lim_{x\to 0}\sin(1/x)$ does not exist.
\end{exercise}

\begin{exercise}\label{exo:ch5-4}
Let $f(x)=x\sin(1/x)$ for $x\neq 0$ and $f(0)=0$.  Show that $f$ is
continuous at~$0$.
\end{exercise}

\begin{exercise}\label{exo:ch5-5}
Prove that $f(x)=\sqrt{x}$ is continuous on $[0,\infty)$ using
$\varepsilon$--$\delta$.

\emph{Hint:} $|\sqrt{x}-\sqrt{a}|
=\frac{|x-a|}{\sqrt{x}+\sqrt{a}}$.
\end{exercise}

\begin{exercise}\label{exo:ch5-6}
Show that if $f$ is continuous on $[a,b]$ and $f(x)>0$ for all
$x\in[a,b]$, then there exists $m>0$ with $f(x)\ge m$ for all $x$.
\end{exercise}

\begin{exercise}\label{exo:ch5-7}
Let $f\colon\mathbb{R}\to\mathbb{R}$ be continuous and suppose
$\lim_{x\to+\infty}f(x)=\lim_{x\to-\infty}f(x)=0$.  Prove that $f$
is bounded and attains its supremum or infimum.
\end{exercise}

\begin{exercise}\label{exo:ch5-8}
Show that $f(x)=\sin x$ is uniformly continuous on $\mathbb{R}$.

\emph{Hint:} $|\sin x-\sin y|=2|\cos(\frac{x+y}{2})\sin(\frac{x-y}{2})|
\le 2\cdot|\frac{x-y}{2}|=|x-y|$.
\end{exercise}

\begin{exercise}\label{exo:ch5-9}
Prove that if $f$ is uniformly continuous on $\mathbb{R}$ and $(x_n)$
is Cauchy, then $(f(x_n))$ is Cauchy.  Show by example this can fail
if $f$ is only (pointwise) continuous.
\end{exercise}

\begin{exercise}\label{exo:ch5-10}
\textbf{(Fixed-point theorem).}\;
Let $f\colon[0,1]\to[0,1]$ be continuous.  Prove that $f$ has a fixed
point, i.e., there exists $c\in[0,1]$ with $f(c)=c$.

\emph{Hint:} consider $g(x)=f(x)-x$ and apply the IVT.
\end{exercise}

\begin{exercise}\label{exo:ch5-11}
Let $f\colon[a,b]\to\mathbb{R}$ be continuous and injective.  Show that
$f$ is strictly monotone.
\end{exercise}

\begin{exercise}\label{exo:ch5-12}
\textbf{(Thomae's function).}\;
Define $t\colon[0,1]\to\mathbb{R}$ by $t(p/q)=1/q$ (in lowest terms,
$q\ge 1$) for rationals, and $t(x)=0$ for irrationals.  Prove that $t$
is continuous at every irrational and discontinuous at every rational.
\end{exercise}

%% ────────────────────────────────────────────────────────────────────
\section*{Chapter Summary}

\begin{tcolorbox}[colback=gray!5,colframe=black!50,
  title={Key results of Chapter 5}]
\begin{itemize}[nosep]
\item Limit of a function: $\varepsilon$--$\delta$ definition and
  sequential characterization (Heine).
\item Continuity: $\varepsilon$--$\delta$ formulation;
  algebraic operations preserve continuity.
\item \textbf{Intermediate Value Theorem} (IVT): continuous image of
  an interval is an interval.
\item \textbf{Extreme Value Theorem}: continuous on $[a,b]$
  $\Rightarrow$ bounded and attains bounds.
\item \textbf{Heine's theorem}: continuous on $[a,b]$ $\Rightarrow$
  uniformly continuous.
\item Lipschitz $\Rightarrow$ uniformly continuous $\Rightarrow$
  continuous (none reversed in general).
\end{itemize}
\end{tcolorbox}


%% ====================================================================
%%  CHAPTER 6 — DIFFERENTIATION
%% ====================================================================
\chapter{Differentiation}\label{ch:differentiation}

\minitoc\bigskip

\lettrine[lines=3,lhang=0.15]{D}{ifferentiation} makes precise the
intuitive notion of ``instantaneous rate of change.''  In this chapter
we develop the theory rigorously from the limit definition, prove
all the standard rules, and establish the powerful Mean Value Theorem
and its consequences.  We finish with Taylor's formula, which gives
polynomial approximations of arbitrary order.

%% ────────────────────────────────────────────────────────────────────
\section{The Derivative}

\begin{tcolorbox}[colback=blue!3,colframe=blue!50!black,
  title={Definition of the derivative},fonttitle=\bfseries]
\index{derivative!definition}
Let $f\colon I\to\mathbb{R}$ where $I$ is an open interval, and let
$a\in I$.  The \textbf{derivative of $f$ at~$a$} is
\[
  f'(a) \;=\; \lim_{h\to 0}\frac{f(a+h)-f(a)}{h}
  \;=\; \lim_{x\to a}\frac{f(x)-f(a)}{x-a},
\]
provided this limit exists (and is finite).  If $f'(a)$ exists, we
say $f$ is \textbf{differentiable at~$a$}.
\end{tcolorbox}

\begin{theorem}[Differentiable $\Rightarrow$ continuous]
\label{thm:diff-implies-cont}
\index{differentiability implies continuity}
If $f$ is differentiable at~$a$, then $f$ is continuous at~$a$.
\end{theorem}

\begin{proof}
Write, for $x\neq a$,
\[
  f(x)-f(a) = \frac{f(x)-f(a)}{x-a}\cdot(x-a).
\]
As $x\to a$, the first factor tends to $f'(a)$ (finite) and the second
to~$0$.  By the product rule for limits,
\[
  \lim_{x\to a}\bigl(f(x)-f(a)\bigr) = f'(a)\cdot 0 = 0,
\]
so $f(x)\to f(a)$, i.e.\ $f$ is continuous at~$a$.
\end{proof}

\begin{remark}
The converse is false.  The function $f(x)=|x|$ is continuous at $0$
but not differentiable there, as we show next.
\end{remark}

\subsection{Counterexamples}

\begin{example}[$|x|$ at $0$]\label{ex:abs-not-diff}
\index{absolute value!not differentiable}
We have
\[
  \frac{|h|-|0|}{h} = \frac{|h|}{h} =
  \begin{cases} 1 & h>0,\\ -1 & h<0.\end{cases}
\]
The right-hand limit is $1$, the left-hand limit is $-1$; they differ,
so $f'(0)$ does not exist.
\end{example}

\begin{center}
\begin{tikzpicture}[scale=1,>=Stealth]
  \draw[->] (-3,0) -- (3,0) node[right]{$x$};
  \draw[->] (0,-0.5) -- (0,3) node[above]{$y$};
  \draw[thick,blue] (-2.5,2.5) -- (0,0) -- (2.5,2.5);
  \node[font=\footnotesize,blue] at (2,2.6) {$y=|x|$};
  \filldraw (0,0) circle (1.5pt);
  \draw[->,red,thick] (-1.5,1.5) -- (-0.3,0.3);
  \draw[->,red,thick] (1.5,1.5) -- (0.3,0.3);
  \node[font=\footnotesize,red] at (-2.2,1.8) {slope $=-1$};
  \node[font=\footnotesize,red] at (2.2,1.8) {slope $=+1$};
  \node[font=\footnotesize] at (0,-0.8) {Corner: not differentiable};
\end{tikzpicture}
\captionof{figure}{$|x|$ has a corner at the origin.}\label{fig:abs-corner}
\end{center}

\begin{example}[$x^2\sin(1/x)$]\label{ex:x2sin1x}
\index{differentiability!$x^2\sin(1/x)$}
Define $g(x)=x^2\sin(1/x)$ for $x\neq 0$ and $g(0)=0$.  Then
\[
  \frac{g(h)-g(0)}{h} = \frac{h^2\sin(1/h)}{h} = h\sin(1/h) \to 0
  \quad(h\to 0),
\]
since $|h\sin(1/h)|\le|h|\to 0$.  So $g'(0)=0$.  However,
\[
  g'(x) = 2x\sin(1/x)-\cos(1/x) \quad(x\neq 0),
\]
and $\lim_{x\to 0}g'(x)$ does not exist (the $\cos(1/x)$ term
oscillates).  So $g$ is differentiable everywhere, but $g'$ is not
continuous at~$0$.
\end{example}

%% ────────────────────────────────────────────────────────────────────
\section{Differentiation Rules}

\begin{theorem}[Sum rule]\label{thm:sum-rule}
\index{differentiation rules!sum}
If $f$ and $g$ are differentiable at~$a$, then $f+g$ is
differentiable at~$a$ and $(f+g)'(a)=f'(a)+g'(a)$.
\end{theorem}

\begin{proof}
\[
  \frac{(f+g)(a+h)-(f+g)(a)}{h}
  = \frac{f(a+h)-f(a)}{h} + \frac{g(a+h)-g(a)}{h}
  \to f'(a)+g'(a). \qedhere
\]
\end{proof}

\begin{theorem}[Product rule / Leibniz rule]\label{thm:product-rule}
\index{differentiation rules!product}\index{Leibniz rule}
If $f$ and $g$ are differentiable at~$a$, then
\[
  (fg)'(a) = f'(a)\,g(a) + f(a)\,g'(a).
\]
\end{theorem}

\begin{proof}
We use the standard trick of adding and subtracting $f(a+h)g(a)$:
\begin{align*}
  \frac{f(a+h)g(a+h)-f(a)g(a)}{h}
  &= \frac{f(a+h)g(a+h) - f(a+h)g(a) + f(a+h)g(a) - f(a)g(a)}{h}\\
  &= f(a+h)\cdot\frac{g(a+h)-g(a)}{h}
     + g(a)\cdot\frac{f(a+h)-f(a)}{h}.
\end{align*}
As $h\to 0$: $f(a+h)\to f(a)$ (by
Theorem~\ref{thm:diff-implies-cont}),
$\frac{g(a+h)-g(a)}{h}\to g'(a)$, and
$\frac{f(a+h)-f(a)}{h}\to f'(a)$.  Hence
\[
  (fg)'(a) = f(a)\,g'(a) + g(a)\,f'(a). \qedhere
\]
\end{proof}

\begin{theorem}[Quotient rule]\label{thm:quotient-rule}
\index{differentiation rules!quotient}
If $f$ and $g$ are differentiable at~$a$ and $g(a)\neq 0$, then
\[
  \Bigl(\frac{f}{g}\Bigr)'(a)
  = \frac{f'(a)\,g(a)-f(a)\,g'(a)}{g(a)^2}.
\]
\end{theorem}

\begin{proof}
We first show that $(1/g)'(a)=-g'(a)/g(a)^2$.  Write
\[
  \frac{1/g(a+h)-1/g(a)}{h}
  = \frac{g(a)-g(a+h)}{h\,g(a+h)\,g(a)}
  = \frac{-1}{g(a+h)\,g(a)}\cdot\frac{g(a+h)-g(a)}{h}.
\]
As $h\to 0$, $g(a+h)\to g(a)\neq 0$ (continuity) and the difference
quotient $\to g'(a)$.  Hence $(1/g)'(a)=-g'(a)/g(a)^2$.

The full quotient rule follows from $f/g=f\cdot(1/g)$ and the product
rule:
\[
  (f/g)' = f'\cdot(1/g) + f\cdot(1/g)'
  = \frac{f'}{g} - \frac{f\,g'}{g^2}
  = \frac{f'g - fg'}{g^2}. \qedhere
\]
\end{proof}

\begin{theorem}[Chain rule]\label{thm:chain-rule}
\index{chain rule}\index{differentiation rules!chain}
Let $f$ be differentiable at~$a$ and $g$ differentiable at $f(a)$.
Then $g\circ f$ is differentiable at~$a$ and
\[
  (g\circ f)'(a) = g'\bigl(f(a)\bigr)\cdot f'(a).
\]
\end{theorem}

\begin{proof}
Define an auxiliary function
\[
  \varphi(y) =
  \begin{cases}
    \displaystyle\frac{g(y)-g(b)}{y-b} & y\neq b,\\[6pt]
    g'(b) & y=b,
  \end{cases}
\]
where $b=f(a)$.  Then $\varphi$ is continuous at~$b$ (because $g$ is
differentiable at~$b$) and
\[
  g(y)-g(b) = \varphi(y)\,(y-b)
  \quad\text{for all } y.
\]
Setting $y=f(x)$:
\[
  g(f(x))-g(f(a)) = \varphi(f(x))\,\bigl(f(x)-f(a)\bigr).
\]
Dividing by $x-a$ ($x\neq a$):
\[
  \frac{g(f(x))-g(f(a))}{x-a}
  = \varphi(f(x))\cdot\frac{f(x)-f(a)}{x-a}.
\]
As $x\to a$: $f(x)\to f(a)=b$ (continuity of~$f$), so
$\varphi(f(x))\to\varphi(b)=g'(b)$, and
$\frac{f(x)-f(a)}{x-a}\to f'(a)$.  Hence
$(g\circ f)'(a)=g'(f(a))\cdot f'(a)$.
\end{proof}

\begin{theorem}[Derivative of an inverse function]
\label{thm:inv-deriv}\index{inverse function!derivative}
Let $f\colon I\to J$ be continuous, strictly monotone, and
differentiable at $a\in I$ with $f'(a)\neq 0$.  Then $f^{-1}$ is
differentiable at $b=f(a)$ and
\[
  (f^{-1})'(b) = \frac{1}{f'(a)} = \frac{1}{f'(f^{-1}(b))}.
\]
\end{theorem}

\begin{proof}
Let $b=f(a)$ and set $g=f^{-1}$.  For $k\neq 0$ with $b+k\in J$,
let $h=g(b+k)-g(b)=g(b+k)-a$.  Then $h\neq 0$ (because $f$ is
injective) and $f(a+h)=b+k$, so $k=f(a+h)-f(a)$.  Thus
\[
  \frac{g(b+k)-g(b)}{k}
  = \frac{h}{f(a+h)-f(a)}
  = \frac{1}{\dfrac{f(a+h)-f(a)}{h}}.
\]
As $k\to 0$, we have $b+k\to b$, so $g(b+k)\to g(b)=a$ (continuity
of~$g$), hence $h\to 0$.  The denominator tends to $f'(a)\neq 0$.
Therefore $g'(b)=1/f'(a)$.
\end{proof}

%% ────────────────────────────────────────────────────────────────────
\section{Local Extrema and the Mean Value Theorem}

\begin{theorem}[Fermat's theorem]\label{thm:fermat}
\index{Fermat's theorem (local extrema)}
Let $f\colon(a,b)\to\mathbb{R}$ be differentiable at $c\in(a,b)$.
If $f$ has a local extremum at~$c$, then $f'(c)=0$.
\end{theorem}

\begin{proof}
Suppose $f$ has a local maximum at~$c$ (the minimum case is
analogous).  There exists $r>0$ with $(c-r,c+r)\subset(a,b)$ and
$f(x)\le f(c)$ for $|x-c|<r$.

\begin{itemize}
\item For $0<h<r$: $f(c+h)-f(c)\le 0$, so
  $\frac{f(c+h)-f(c)}{h}\le 0$.  Taking $h\to 0^+$:
  $f'(c)\le 0$.
\item For $-r<h<0$: $f(c+h)-f(c)\le 0$ and $h<0$, so
  $\frac{f(c+h)-f(c)}{h}\ge 0$.  Taking $h\to 0^-$:
  $f'(c)\ge 0$.
\end{itemize}
Together, $f'(c)=0$.
\end{proof}

\begin{theorem}[Rolle's theorem]\label{thm:rolle}
\index{Rolle's theorem}
Let $f\colon[a,b]\to\mathbb{R}$ be continuous on $[a,b]$ and
differentiable on $(a,b)$, with $f(a)=f(b)$.  Then there exists
$c\in(a,b)$ with $f'(c)=0$.
\end{theorem}

\begin{proof}
By the Extreme Value Theorem (Theorem~\ref{thm:evt}), $f$ attains
its maximum~$M$ and minimum~$m$ on $[a,b]$.

\textbf{Case 1:} $m=M$.  Then $f$ is constant, so $f'(c)=0$ for
every $c\in(a,b)$.

\textbf{Case 2:} $m<M$.  Since $f(a)=f(b)$, the maximum or the
minimum (or both) must be attained at some interior point $c\in(a,b)$.
By Fermat's theorem, $f'(c)=0$.
\end{proof}

\begin{theorem}[Mean Value Theorem (MVT)]\label{thm:mvt}
\index{Mean Value Theorem}
Let $f\colon[a,b]\to\mathbb{R}$ be continuous on $[a,b]$ and
differentiable on $(a,b)$.  Then there exists $c\in(a,b)$ with
\[
  f'(c) = \frac{f(b)-f(a)}{b-a}.
\]
\end{theorem}

\begin{proof}
Define the auxiliary function
\[
  g(x) = f(x) - \frac{f(b)-f(a)}{b-a}(x-a).
\]
Then $g$ is continuous on $[a,b]$, differentiable on $(a,b)$, and
\[
  g(a)=f(a), \quad
  g(b)=f(b)-\bigl(f(b)-f(a)\bigr)=f(a).
\]
So $g(a)=g(b)$.  By Rolle's theorem, there exists $c\in(a,b)$ with
$g'(c)=0$, i.e.\
\[
  0 = g'(c) = f'(c) - \frac{f(b)-f(a)}{b-a},
\]
giving $f'(c)=\frac{f(b)-f(a)}{b-a}$.
\end{proof}

\begin{corollary}[Zero derivative $\Rightarrow$ constant]
\label{cor:zero-deriv}\index{constant function criterion}
If $f$ is differentiable on an open interval $I$ and $f'(x)=0$ for
all $x\in I$, then $f$ is constant on~$I$.
\end{corollary}

\begin{proof}
For any $x_1<x_2$ in $I$, the MVT gives $c\in(x_1,x_2)$ with
$f(x_2)-f(x_1)=f'(c)(x_2-x_1)=0$.
\end{proof}

\begin{corollary}[Positive derivative $\Rightarrow$ increasing]
\label{cor:pos-deriv}\index{monotonicity criterion}
If $f'(x)>0$ for all $x\in(a,b)$, then $f$ is strictly increasing
on $(a,b)$.
\end{corollary}

\begin{proof}
For $x_1<x_2$:
$f(x_2)-f(x_1)=f'(c)(x_2-x_1)>0$ since $f'(c)>0$ and $x_2-x_1>0$.
\end{proof}

%% ────────────────────────────────────────────────────────────────────
\section{L'H\^opital's Rule}

\begin{theorem}[L'H\^opital's rule, $0/0$ case]
\label{thm:lhopital}\index{L'H\^opital's rule}
Let $f,g\colon(a,b)\to\mathbb{R}$ be differentiable with
$g'(x)\neq 0$ on $(a,b)$.  Suppose
\[
  \lim_{x\to a^+}f(x)=0=\lim_{x\to a^+}g(x)
  \quad\text{and}\quad
  \lim_{x\to a^+}\frac{f'(x)}{g'(x)}=L\in\mathbb{R}\cup\{\pm\infty\}.
\]
Then $\displaystyle\lim_{x\to a^+}\frac{f(x)}{g(x)}=L$.
\end{theorem}

\begin{proof}
Extend $f$ and $g$ continuously to $[a,b)$ by setting $f(a)=g(a)=0$.
Fix $x\in(a,b)$.  By the Cauchy (generalized) Mean Value Theorem
(apply Rolle to $h(t)=f(t)\bigl(g(x)-g(a)\bigr)-g(t)\bigl(f(x)-f(a)\bigr)$),
there exists $c_x\in(a,x)$ with
\[
  f'(c_x)\bigl(g(x)-g(a)\bigr) = g'(c_x)\bigl(f(x)-f(a)\bigr).
\]
Since $f(a)=g(a)=0$:
\[
  \frac{f(x)}{g(x)} = \frac{f'(c_x)}{g'(c_x)},
\]
provided $g(x)\neq 0$ (which holds for $x$ close to $a$, since
$g'(c_x)\neq 0$ and $g(a)=0$ would force $g(x)\neq 0$ for small
$x-a$ by the MVT).  As $x\to a^+$, we have $c_x\to a^+$ (because
$a<c_x<x$), so $\frac{f'(c_x)}{g'(c_x)}\to L$.
\end{proof}

\begin{example}
$\displaystyle\lim_{x\to 0}\frac{e^x-1-x}{x^2}
= \lim_{x\to 0}\frac{e^x-1}{2x}
= \lim_{x\to 0}\frac{e^x}{2} = \frac{1}{2}$
(applying L'H\^opital twice; each application is of the $0/0$ form).
\end{example}

%% ────────────────────────────────────────────────────────────────────
\section{Taylor's Formula}

\subsection{Taylor Polynomial and Lagrange Remainder}

\begin{definition}[Taylor polynomial]\label{def:taylor-poly}
\index{Taylor polynomial}
Let $f$ be $n$ times differentiable at~$a$.  The \textbf{Taylor
polynomial of order~$n$ of $f$ at~$a$} is
\[
  T_n(x) = \sum_{k=0}^{n}\frac{f^{(k)}(a)}{k!}(x-a)^k.
\]
\end{definition}

\begin{theorem}[Taylor--Lagrange]\label{thm:taylor-lagrange}
\index{Taylor--Lagrange formula}
Let $f\colon[a,x]\to\mathbb{R}$ (or $[x,a]$) be such that $f^{(n)}$
is continuous on the closed interval and $f^{(n+1)}$ exists on the
open interval.  Then there exists $c$ strictly between $a$ and $x$
with
\[
  f(x) = T_n(x) + \frac{f^{(n+1)}(c)}{(n+1)!}(x-a)^{n+1}.
\]
The last term is the \textbf{Lagrange remainder} $R_n(x)$.
\end{theorem}

\begin{proof}
Define $R_n(x)=f(x)-T_n(x)$.  Note that $R_n(a)=R_n'(a)=\cdots
=R_n^{(n)}(a)=0$ (by construction of the Taylor polynomial) and
$R_n^{(n+1)}=f^{(n+1)}$.

We seek $c$ with $R_n(x)=\frac{f^{(n+1)}(c)}{(n+1)!}(x-a)^{n+1}$.
Define $g(t)=(x-t)^{n+1}$; note $g(a)=(x-a)^{n+1}$, $g(x)=0$.

Apply the Cauchy MVT to $R_n$ and $g$ on $[a,x]$:
\[
  \frac{R_n(x)-R_n(a)}{g(x)-g(a)}
  = \frac{R_n'(c_1)}{g'(c_1)}
\]
for some $c_1$ between $a$ and $x$.  Since $R_n(a)=0$ and $g(x)=0$:
\[
  \frac{R_n(x)}{-(x-a)^{n+1}}
  = \frac{R_n'(c_1)}{-(n+1)(x-c_1)^n}.
\]
Apply the Cauchy MVT again to $R_n'$ and $(x-t)^n$ on $[a,c_1]$
(using $R_n'(a)=0$), and iterate this process $n$ times.  After $n+1$
applications, we arrive at a point $c=c_{n+1}$ between $a$ and $x$
with
\[
  R_n(x) = \frac{R_n^{(n+1)}(c)}{(n+1)!}(x-a)^{n+1}
  = \frac{f^{(n+1)}(c)}{(n+1)!}(x-a)^{n+1}. \qedhere
\]
\end{proof}

\subsection{Taylor--Young Formula (Landau Notation)}

\begin{definition}[Landau notation]\label{def:landau}
\index{Landau notation}\index{little-o notation}
We write $f(x)=o(g(x))$ as $x\to a$ if
$\lim_{x\to a}\frac{f(x)}{g(x)}=0$ (assuming $g(x)\neq 0$ near $a$,
$a$ excluded).  In particular, $f(x)=o(1)$ means $f(x)\to 0$.
\end{definition}

\begin{theorem}[Taylor--Young]\label{thm:taylor-young}
\index{Taylor--Young formula}
If $f$ is $n$ times differentiable at~$a$, then
\[
  f(x) = \sum_{k=0}^{n}\frac{f^{(k)}(a)}{k!}(x-a)^k
  + o\bigl((x-a)^n\bigr)
  \quad\text{as } x\to a.
\]
\end{theorem}

\begin{proof}
We prove by induction on $n$.

\textbf{Base case $n=0$:}\; $f(x)=f(a)+o(1)$ says $f(x)\to f(a)$,
which is continuity at~$a$.  Differentiability at $a$ implies
continuity, so this holds.

\textbf{Base case $n=1$:}\; $f(x)=f(a)+f'(a)(x-a)+o(x-a)$ is
exactly the definition of differentiability.

\textbf{Inductive step:}\; Assume the result for $n-1$ applied to
$f'$ (which is $(n-1)$-times differentiable at~$a$):
\[
  f'(x) = \sum_{k=0}^{n-1}\frac{f^{(k+1)}(a)}{k!}(x-a)^k
  + o\bigl((x-a)^{n-1}\bigr).
\]
Define $R_n(x)=f(x)-T_n(x)$, so $R_n(a)=0$ and
$R_n'(x)=f'(x)-T_n'(x)$.  One checks that
$T_n'(x)=\sum_{k=1}^{n}\frac{f^{(k)}(a)}{(k-1)!}(x-a)^{k-1}
=\sum_{j=0}^{n-1}\frac{f^{(j+1)}(a)}{j!}(x-a)^j$.
By the inductive hypothesis, $R_n'(x)=o((x-a)^{n-1})$.

We need $R_n(x)=o((x-a)^n)$.  By the Mean Value Theorem applied to
$R_n$ on $[a,x]$: $R_n(x)=R_n(x)-R_n(a)=R_n'(\xi)(x-a)$ for some
$\xi$ between $a$ and $x$.  Then
\[
  \frac{R_n(x)}{(x-a)^n}
  = \frac{R_n'(\xi)}{(x-a)^{n-1}}.
\]
For any $\varepsilon>0$, by the inductive hypothesis there exists
$\delta>0$ with $|R_n'(t)|\le\varepsilon|t-a|^{n-1}$ for $|t-a|<\delta$.
Since $|\xi-a|\le|x-a|$, for $|x-a|<\delta$ we get
$|R_n'(\xi)|\le\varepsilon|\xi-a|^{n-1}\le\varepsilon|x-a|^{n-1}$,
hence $\bigl|\frac{R_n(x)}{(x-a)^n}\bigr|\le\varepsilon$.
Thus $R_n(x)=o((x-a)^n)$.
\end{proof}

\subsection{Standard Taylor Expansions}

\begin{tcolorbox}[colback=green!3,colframe=green!50!black,
  title={Table of standard Taylor expansions at $a=0$},
  fonttitle=\bfseries]
\index{Taylor expansion!standard}
\begin{align}
  e^x &= \sum_{k=0}^{n}\frac{x^k}{k!} + o(x^n)
  = 1+x+\frac{x^2}{2}+\frac{x^3}{6}+\cdots \label{eq:taylor-exp}\\
  \sin x &= \sum_{k=0}^{n}\frac{(-1)^k x^{2k+1}}{(2k+1)!} + o(x^{2n+2})
  = x-\frac{x^3}{6}+\frac{x^5}{120}-\cdots \label{eq:taylor-sin}\\
  \cos x &= \sum_{k=0}^{n}\frac{(-1)^k x^{2k}}{(2k)!} + o(x^{2n+1})
  = 1-\frac{x^2}{2}+\frac{x^4}{24}-\cdots \label{eq:taylor-cos}\\
  \ln(1+x) &= \sum_{k=1}^{n}\frac{(-1)^{k+1} x^k}{k} + o(x^n)
  = x-\frac{x^2}{2}+\frac{x^3}{3}-\cdots \label{eq:taylor-ln}\\
  (1+x)^\alpha &= \sum_{k=0}^{n}\binom{\alpha}{k}x^k + o(x^n),
  \quad\binom{\alpha}{k}=\frac{\alpha(\alpha-1)\cdots(\alpha-k+1)}{k!}
  \label{eq:taylor-binom}
\end{align}
\end{tcolorbox}

\begin{center}
\begin{tikzpicture}[scale=0.85,>=Stealth]
  \draw[->] (-4.5,0) -- (4.5,0) node[right]{$x$};
  \draw[->] (0,-2) -- (0,2) node[above]{$y$};
  % sin(x)
  \draw[thick,blue,domain=-4:4,samples=100]
    plot(\x,{sin(\x r)}) node[right,font=\footnotesize]{$\sin x$};
  % T1
  \draw[red,domain=-2:2,samples=50]
    plot(\x,{\x}) node[right,font=\tiny]{$T_1$};
  % T3
  \draw[orange,domain=-3.3:3.3,samples=80]
    plot(\x,{\x - \x*\x*\x/6})
    node[right,font=\tiny]{$T_3$};
  % T5
  \draw[green!50!black,domain=-3.5:3.5,samples=100]
    plot(\x,{\x - \x*\x*\x/6 + \x*\x*\x*\x*\x/120})
    node[right,font=\tiny]{$T_5$};
  % T7
  \draw[purple,domain=-3.5:3.5,samples=100]
    plot(\x,{\x - \x*\x*\x/6 + \x*\x*\x*\x*\x/120
      - \x*\x*\x*\x*\x*\x*\x/5040})
    node[right,font=\tiny]{$T_7$};
\end{tikzpicture}
\captionof{figure}{Taylor polynomials $T_1,T_3,T_5,T_7$ of
  $\sin x$ at~$0$.}\label{fig:taylor-sin}
\end{center}

\subsection{Operations on Taylor Expansions and Applications}

Taylor expansions can be combined using the following rules:
\begin{itemize}
\item \textbf{Sum/Difference:} Expand each function, add/subtract
  term by term, discard terms beyond order~$n$.
\item \textbf{Product:} Expand each factor, multiply formally, keep
  only terms up to order~$n$.
\item \textbf{Composition:} If $f(x)=o(1)$, substitute the expansion
  of~$f$ into the expansion of~$g$ and truncate.
\item \textbf{Quotient:} Use long division or write
  $f/g = f\cdot(1/g)$ and expand $1/g$ via the geometric series or
  successive identification.
\end{itemize}

\begin{example}[Application to a limit]\label{ex:taylor-limit-1}
Compute $\displaystyle\lim_{x\to 0}\frac{e^x - 1 - x}{x^2}$.

\medskip
\textbf{Solution.}\;
$e^x = 1 + x + \frac{x^2}{2} + o(x^2)$, so
$e^x - 1 - x = \frac{x^2}{2} + o(x^2)$.  Hence
\[
  \frac{e^x - 1 - x}{x^2}
  = \frac{\frac{x^2}{2}+o(x^2)}{x^2}
  = \frac{1}{2} + o(1) \to \frac{1}{2}. \qed
\]
\end{example}

\begin{example}[Application to a limit]\label{ex:taylor-limit-2}
Compute $\displaystyle\lim_{x\to 0}\frac{\tan x - \sin x}{x^3}$.

\medskip
\textbf{Solution.}\;
$\sin x = x - \frac{x^3}{6} + o(x^3)$,\;
$\cos x = 1 - \frac{x^2}{2} + o(x^2)$, so
\[
  \tan x = \frac{\sin x}{\cos x}
  = \frac{x - \frac{x^3}{6}+o(x^3)}{1-\frac{x^2}{2}+o(x^2)}.
\]
Using long division (or multiplying numerator by
$1+\frac{x^2}{2}+o(x^2)$):
\[
  \tan x = x + \frac{x^3}{3} + o(x^3).
\]
Therefore
\[
  \frac{\tan x - \sin x}{x^3}
  = \frac{\bigl(x+\frac{x^3}{3}+o(x^3)\bigr)
    -\bigl(x-\frac{x^3}{6}+o(x^3)\bigr)}{x^3}
  = \frac{\frac{x^3}{2}+o(x^3)}{x^3}
  = \frac{1}{2}+o(1) \to \frac{1}{2}. \qed
\]
\end{example}

\begin{example}[Asymptotic expansion]\label{ex:taylor-limit-3}
Find the asymptotic expansion of $\sqrt{1+x}-\sqrt{1-x}$ as $x\to 0$,
up to order~3.

\medskip
\textbf{Solution.}\;
Using \eqref{eq:taylor-binom} with $\alpha=1/2$:
\[
  (1+x)^{1/2} = 1 + \tfrac{1}{2}x - \tfrac{1}{8}x^2
  + \tfrac{1}{16}x^3 + o(x^3),
\]
\[
  (1-x)^{1/2} = 1 - \tfrac{1}{2}x - \tfrac{1}{8}x^2
  - \tfrac{1}{16}x^3 + o(x^3).
\]
Subtracting:
\[
  \sqrt{1+x}-\sqrt{1-x}
  = x + \tfrac{1}{8}x^3 + o(x^3). \qed
\]
\end{example}

%% ────────────────────────────────────────────────────────────────────
\section{Common Errors}

\begin{tcolorbox}[colback=red!4,colframe=red!70!black,
  title={Frequent mistakes in Chapter~6}]
\begin{enumerate}[nosep]
\item \textbf{``$f$ continuous at $a$ $\Rightarrow$ $f$ differentiable
  at $a$.''}\;
  False: $|x|$ is continuous but not differentiable at $0$.
\item \textbf{Using L'H\^opital when the hypotheses are not satisfied.}\;
  Check that both numerator and denominator tend to~$0$ (or $\infty$)
  and that the derivative of the denominator is nonzero.
\item \textbf{Forgetting the chain rule.}\;
  $\frac{d}{dx}\sin(x^2)\neq\cos(x^2)$; the correct answer is
  $2x\cos(x^2)$.
\item \textbf{Confusing Taylor--Lagrange and Taylor--Young.}\;
  Lagrange gives an exact remainder with an unknown $c$;
  Young gives a qualitative $o((x-a)^n)$ remainder.
\item \textbf{Truncating Taylor expansions at the wrong order.}\;
  When computing a limit $f(x)/x^n$, expand $f$ to at least
  order~$n$ to get the leading constant.
\item \textbf{MVT: forgetting the hypotheses.}\;
  The function must be continuous on the \emph{closed} interval and
  differentiable on the \emph{open} interval.
\end{enumerate}
\end{tcolorbox}

%% ────────────────────────────────────────────────────────────────────
\section{Exercises}

\begin{exercise}\label{exo:ch6-1}
Using the definition, compute $f'(x)$ for $f(x)=x^3$.
\end{exercise}

\begin{exercise}\label{exo:ch6-2}
Let $f(x)=x^{1/3}$.  Show that $f$ is not differentiable at $0$.
\end{exercise}

\begin{exercise}\label{exo:ch6-3}
Prove the \textbf{generalized product rule} (Leibniz formula):
\[
  (fg)^{(n)} = \sum_{k=0}^{n}\binom{n}{k}f^{(k)}\,g^{(n-k)}.
\]
\emph{Hint:} induction on $n$.
\end{exercise}

\begin{exercise}\label{exo:ch6-4}
Let $f(x)=\begin{cases}x^2\sin(1/x)&x\neq 0,\\0&x=0.\end{cases}$\;
Show $f$ is differentiable on $\mathbb{R}$ but $f'$ is discontinuous
at~$0$.
\end{exercise}

\begin{exercise}\label{exo:ch6-5}
Use the MVT to prove:
$|\sin a - \sin b|\le |a-b|$ for all $a,b\in\mathbb{R}$.
\end{exercise}

\begin{exercise}\label{exo:ch6-6}
Show that $e^x > 1+x$ for all $x\neq 0$.

\emph{Hint:} study $g(x)=e^x-1-x$.
\end{exercise}

\begin{exercise}\label{exo:ch6-7}
Compute the following limits using Taylor expansions:
\begin{enumerate}
\item $\displaystyle\lim_{x\to 0}
  \frac{\cos x - 1 + \frac{x^2}{2}}{x^4}$
\item $\displaystyle\lim_{x\to 0}
  \frac{\ln(1+x)-x+\frac{x^2}{2}}{x^3}$
\item $\displaystyle\lim_{x\to 0}
  \frac{e^{\sin x}-1-x}{x^2}$
\end{enumerate}
\end{exercise}

\begin{exercise}\label{exo:ch6-8}
Find the Taylor expansion of $\frac{1}{1-x}$ at $a=0$ to order~$n$
and identify the Lagrange remainder explicitly.
\end{exercise}

\begin{exercise}\label{exo:ch6-9}
Prove that if $f''(x)\ge 0$ on an interval $I$, then $f$ is convex,
i.e., $f\bigl(\lambda x+(1-\lambda)y\bigr)\le\lambda f(x)+(1-\lambda)f(y)$
for all $x,y\in I$, $\lambda\in[0,1]$.

\emph{Hint:} use the MVT on both subintervals.
\end{exercise}

\begin{exercise}\label{exo:ch6-10}
\textbf{(Darboux's theorem).}\;
Let $f$ be differentiable on $[a,b]$.  Show that $f'$ has the
intermediate value property (even though $f'$ may not be continuous).

\emph{Hint:} consider $g(x)=f(x)-\lambda x$ and apply Fermat's theorem.
\end{exercise}

\begin{exercise}\label{exo:ch6-11}
Compute the Taylor expansion of $\arctan x$ at $a=0$ to order~$2n+1$.

\emph{Hint:} integrate the geometric series
$\frac{1}{1+t^2}=\sum_{k=0}^{n}(-1)^k t^{2k}+o(t^{2n})$.
\end{exercise}

\begin{exercise}\label{exo:ch6-12}
Let $f\colon\mathbb{R}\to\mathbb{R}$ be twice differentiable with
$f(0)=0$, $f(1)=1$, $f'(0)=f'(1)=0$.  Prove there exists
$c\in(0,1)$ with $|f''(c)|\ge 4$.

\emph{Hint:} apply the MVT twice, or use Taylor--Lagrange at $a=0$
and $a=1$ evaluated at $x=1/2$.
\end{exercise}

%% ────────────────────────────────────────────────────────────────────
\section*{Chapter Summary}

\begin{tcolorbox}[colback=gray!5,colframe=black!50,
  title={Key results of Chapter 6}]
\begin{itemize}[nosep]
\item Derivative: defined as a limit; differentiable $\Rightarrow$
  continuous (not conversely).
\item Algebraic rules: sum, product (Leibniz), quotient, chain rule,
  inverse function --- all proved.
\item \textbf{Fermat's theorem}: local extremum $\Rightarrow$
  $f'(c)=0$.
\item \textbf{Rolle's theorem}: $f(a)=f(b)$ $\Rightarrow$ $\exists\,
  c,\;f'(c)=0$.
\item \textbf{Mean Value Theorem}: $f'(c)=\frac{f(b)-f(a)}{b-a}$.
\item Corollaries: $f'=0\Rightarrow$ constant; $f'>0\Rightarrow$
  increasing.
\item \textbf{L'H\^opital's rule} for indeterminate forms $0/0$.
\item \textbf{Taylor--Lagrange} (exact remainder) and
  \textbf{Taylor--Young} (asymptotic remainder).
\item Standard expansions: $e^x,\sin x,\cos x,\ln(1+x),(1+x)^\alpha$.
\item Taylor expansions as a tool for computing limits and asymptotic
  analysis.
\end{itemize}
\end{tcolorbox}


%% ====================================================================
%%  BIBLIOGRAPHY
%% ====================================================================
\begin{thebibliography}{99}

\bibitem{rudin}
W.~Rudin,
\textit{Principles of Mathematical Analysis},
3rd~ed., McGraw-Hill, 1976.
\index{Rudin}

\bibitem{zuily}
C.~Zuily and H.~Queff\'elec,
\textit{Analyse pour l'agr\'egation},
4th~ed., Dunod, 2013.
\index{Zuily--Queff\'elec}

\bibitem{ramis}
E.~Ramis, C.~Deschamps, and J.~Odoux,
\textit{Cours de math\'ematiques sp\'eciales~-- Analyse},
Masson, 1991.
\index{Ramis--Deschamps--Odoux}

\bibitem{tao}
T.~Tao,
\textit{Analysis~I},
3rd~ed., Texts and Readings in Mathematics~37,
Hindustan Book Agency / Springer, 2016.
\index{Tao}

\bibitem{pompeiani}
F.~Pompe\"{i}ani and P.~Tarrago,
\textit{Analyse -- Cours et exercices},
Ellipses, 2008.
\index{Pompe\"{i}ani--Tarrago}

\end{thebibliography}

%% ====================================================================
%%  INDEX
%% ====================================================================
\printindex

\end{document}
